% ============================================================
%  R. Nielsen, Om „Den gode Villie" som Magt i Videnskaben
%  [On "the Good Will" as a Power in Science]
%  Copenhagen: Gyldendal (F. Hegel), printed by J. H. Schultz, 1867.
%  141 pp. + [1] p. Rettelser.
%  TRANSLATION (English)
%  Source of truth: transcription.tex (Danish), COMPLETE.
%  Mirrors transcription.tex 1:1; every printed page carries \opage{N}
%  at the corresponding point in the English. See TRANSLATION-PLAYBOOK.md.
%  Translated 22 Sep 2026 in ten batches (the transcription's batches).
%  Latin, Greek and German kept as printed; the book's Rettelser are not
%  applied in the text but flagged at each site by a % Errata: comment.
% ============================================================
\documentclass[11pt]{book}
\usepackage[utf8]{inputenc}
\usepackage[T1]{fontenc}
\usepackage{libertinus}
\usepackage{libertinust1math}
\usepackage{textalpha}
\usepackage{graphicx}
\DeclareUnicodeCharacter{0254}{\reflectbox{c}}
\usepackage{amsmath}
\usepackage[english]{babel}
\usepackage[protrusion=true,expansion=true]{microtype}
\usepackage{enumitem}
\usepackage{fancyhdr}
\setlength{\headheight}{28pt}
\addtolength{\topmargin}{-13.5pt}
\pagestyle{fancy}
\fancyhf{}
\fancyhead[LE,RO]{\thepage}
\fancyhead[RE]{\textit{On ``the Good Will'' as a Power in Science}}
\fancyhead[LO]{\textit{\leftmark}}
\renewcommand{\thefootnote}{\ifcase\value{footnote}\or *)\or **)\or ***)\else
  \arabic{footnote})\fi}
\usepackage{setspace}
\setstretch{1.3}
\usepackage{marginnote}
\newcommand{\opage}[1]{\marginnote{\footnotesize\textit{[#1]}}}
\usepackage{hyperref}
\hypersetup{pdftitle={On "the Good Will" as a Power in Science},pdfauthor={R. Nielsen},hidelinks}

\begin{document}

\frontmatter
\begin{titlepage}
\centering
\vspace*{3cm}
{\LARGE On}\\[1.5em]
{\Huge ``the Good Will''}\\[1.5em]
{\LARGE as a Power in Science.}\\[2em]
{\small By}\\[1em]
{\large R. Nielsen.}\\[5em]
Copenhagen.\\[0.5em]
Published by the Gyldendal Bookshop (F. Hegel.)\\[0.3em]
{\footnotesize Printed by J. H. Schultz.}\\[0.5em]
1867.\\[4em]
{\small\itshape Translated from the Danish}
\end{titlepage}

\tableofcontents

\mainmatter
\markboth{On ``the Good Will''}{On ``the Good Will''}

% --- p. 1 ---
\opage{1}\setcounter{footnote}{0}%
``Οὐκ ἀγαθὸν πολυκοιρανίη, εἷς κοίρανος ἔστω.
`\emph{The rule of many is never of benefit; let One alone be ruler}'. By giving this Homeric word a scientific application, Aristotle meant to say that the dominion of many principles is not of benefit but of ruin, not only for states but also for scientific systems, where it brings about anarchy and inner dissolution, but that one principle ought to be the ruling and all-governing one, just as an army ought to have only one supreme general, since otherwise everything falls into disorder. He adds: `that which is does not wish to be governed badly.' (Metaphys. XII, 10).'' With these words Bishop Martensen has introduced what he calls the ``Concluding Reflections'' to his occasional piece \emph{Om Tro og Viden}; they are in truth remarkable words, words that well deserve that we seriously take them to heart.

What would the antagonists of theology, the spokesmen of negative criticism, what would a D. Strau\ss{}, a Bruno Bauer, an L. Feuerbach think and say when they read these introductory words in a book whose author they did not yet know? They would probably say: ``Look, now here is a fresh and resolute beginning; this is how a reasonable man should take hold of the matter from the ground up: no ambiguities, no circumlocutions! The dominion of many principles is not of benefit but of ruin for scientific systems; let One be ruler! That is reason itself, royal reason, the self-appointed sole ruler, for whose indisputable right we have so long fought and
% --- p. 2 ---
\opage{2}\setcounter{footnote}{0}%
labored. Only when the principle of reason is openly acknowledged as the solely ruling principle in the realm of spirit can one get the better of the half-scientific nuisance that calls itself theology, and at last hope for deliverance from all the murky sentimentality, the moldy, half-reflected immediacy, with which speculative seminarians and thoughtless catechism-priests have all too long been serving us. The dominion of many principles is not of benefit; let One alone be ruler: boldly, indeed with youthful freshness, the author has spoken up for criticism and science, for the autonomy of reason, the absolute unity of scientific self-consciousness!'' Indeed, so apt, so rousing are Bishop Martensen's words, that the most zealous defenders of the autonomy of reason could simply claim them for their own.

But of course, when the champions of autonomy went further into the text and to their great surprise discovered what is really aimed at here, they would surely change their tune. ``So it is not the principle of reason but the principle of faith, and hence the good old principle of miracles, that, with a few modern circumlocutions, is offering to keep the sciences in order. Let One be ruler: faith in theology, theology in faith! A scientific monarchy of faith, a theological absolute universal monarchy: what a monster of a thought, what a thought-monster! But is theology in its dotage, is it entering its second childhood? It is not content, then, to be, as hitherto, a mislaid, cramped, bedraggled science by itself, a science which, when angry criticism challenged it to draw its blade, pleaded so movingly: Let me be, I am too humble for such an honor --- no, memories of the past haunt its old brain; it nourishes ambitious plans, it wants to rise aloft, it wants, as science with the supernatural principle, and so after all as a science of miracles --- \textit{risum teneatis} --- to be what the philosophy of the Absolute, despite its profound speculation, could not be, the central science
% --- p. 3 ---
\opage{3}\setcounter{footnote}{0}%
of the sciences. What, then, has happened here? Are the conjunctures extraordinary, the constellations auspicious; did the sun perhaps enter the sign of the Lion just as Mars was in conjunction with Venus? Or is it considerations of prudence, fine calculation, which now sees the quotient of the favorable and the possible conditions approaching unity? The extraordinary must be presupposed if the universal monarchy is to succeed; if the free-born sciences are really to submit to a science of miracles, then a scientific miracle must take place.''

We have imagined what theology's tormenting spirits, the ``freethinkers,'' ``the children of unbelief who have surfeited themselves on the fruit of knowledge,'' would say when they regarded the matter from the standpoint of reason. But if they now did not confine themselves exclusively to looking at the wrong done to reason, but also set themselves to consider the circumstances of theology, how, I wonder, would they then feel? They would feel badly; for it would vex them if theology, the theology they have so often mocked and scorned, should really take such an upswing, seat itself on the throne, seize the scepter of faith, and say Homerically: Let One be ruler! Yes, it would vex them, and the vexation would increase when in their imagination they could do nothing but depict to themselves all the honor and joy such a rise must procure for the theologians. ``What song, what ringing! Does it not sound in our ears as though we heard hymns of victory from the long-besieged fortress? See, the banner of dogmatics waves high from every tower; hermeneutics is illuminating; exegetics hangs out its flag; even the Introduction to the New Testament hangs out a rag --- it has only a rag. Yes, there is a stir in the theologians' camp: the old, we had nearly said headless, faculty lifts its head, shades its eyes, glimpses the new light, listens to the glad tidings, and the glad tidings are the gospel of the universal monarchy.''

% --- p. 4 ---
\opage{4}\setcounter{footnote}{0}%
Such moods and utterances from a standpoint that would have not only theology but also religion abolished are by no means suggested in order to cast a shadow on the principle set up by Bishop Martensen; on the contrary, they are suggested precisely in order to set the principle in a clearer light, inasmuch as they make it recognizable that here, in the name of theology, a great and resolute thought has been uttered. We have quite rightly come to a final dilemma: either theology must go, or it must take up its original place as the highest science. In the character of the highest science, theology must then be determined as the ruling science, as the science that --- naturally not arbitrarily but with a well-founded right, not immediately but mediately, not despotically but organically --- asserts primacy over all the sciences, since the light in which it sees existence is the highest and purest light, is the holy light of revelation, the light of truth in which all finite cognitions must find their final and highest explanation. The thought is, as one sees, both great and beautiful in itself.

In the absolutely divine idea, living truth is one with holy love, and this love is not only truth in itself but the source of all truth. In the divine knowledge there cannot possibly be two absolutely heterogeneous truths, an impersonal and a personal. The more human knowledge takes on likeness to the divine, the more will the fragmentary be overcome, the more will the impersonal subordinate itself to the personal, and the separation between the heterogeneous truths disappear. But human knowledge can become a true reflection of the divine only on the condition that it proceeds in the highest sense from faith, is and wishes to be a knowledge in faith, from faith.

The highest, purest knowledge in faith is a speculative knowledge. ``Speculation,'' says Schelling, ``is everything; it is to behold in God, to contemplate that which is. Science itself has
% --- p. 5 ---
\opage{5}\setcounter{footnote}{0}%
worth only insofar as it is speculative, that is, contemplation of God as he is. But the time will come when the sciences will more and more cease, and immediate cognition will take their place. Only in the highest science does the mortal eye close, as it is no longer man who sees, but the eternal seeing itself that sees in man.'' The lofty goal aimed at here cannot possibly be reached in abstract speculation; it can only --- by approximation --- be reached in the speculation whose point of departure is the living God, and so only in the full, substantial speculation that begins and ends with faith.

So understood, the principle set up by Bishop Martensen then also acquires a pervasive significance for the development of the encyclopedia of the sciences. ``The encyclopedia,'' says Hegel, ``is really a presentation of the general content of philosophy; for what in the sciences is grounded in reason depends on philosophy, whereas what in them rests on arbitrary and external determinations, or, as it is called, is positive and statutory, lies outside it.'' With this it is stated that all the sciences make up one single system of the sciences. The systematic unity of the sciences is conditioned by a distinction being made, by right of the idea, between lower and higher, subordinate and superordinate sciences, and rests in the last instance on there being an absolute science of the sciences. Had abstract logical knowledge been capable of mastering the Absolute, then it would certainly have been philosophy that became the central science. But since, as indicated above, it is only believing speculation, and consequently not autonomous but theonomous speculation, the speculation that begins and ends with the ethical concept of God, that can give human knowledge its highest content, it is not philosophy but theology that we must acknowledge as the real central science.

% --- p. 6 ---
\opage{6}\setcounter{footnote}{0}%
A doctrine of ideas that aims at Christianity ``must,'' says Bishop Martensen (Leilighedsskr. p. 122), ``first and foremost aim at the idea of the Good, as that which is ἐπέκεινα τῆς οὐσίας, and see to it that a thorough subordination comes into the realm of ideas.'' With this agrees what the Bishop states in an earlier context (p. 102): ``If the \emph{philosophy} of religion or of revelation --- for a philosophy of religion that is not a philosophy of revelation is only a half-measure --- is the \emph{universal} science, insofar as it cognizes revelation in its relation to mythology, to the whole manifold of nature and history, then theology is the \emph{central} science, in that it immerses itself exclusively in the One, in the Kingdom of God as such.''

When, with this information, we rightly consider what weight Bishop Martensen lays on his Homeric-Aristotelian word: Let One be ruler! we cannot for a moment doubt that he, literally and in dead earnest, would demand a subordination, albeit a mediate one, of all relative sciences under his central science; that he, in other words, would establish a new theological universal monarchy.

But a new theological universal monarchy, what an enormous thought! Can a thought of this height, I wonder, go in through the door of actuality without stooping so low that it becomes quite bent? I too have met the lofty thought; it has many times come toward me when I regarded dogmatics, not from its shadow side but from its bright side; for although, in accordance with my task, I have most often had to wander through its shadowed parts, I have not on that account forgotten where I can still find the sunlit places. But when, after having beheld the ideal in all its glory, I then seriously asked myself: can such an ideal be realized? I have always received a negative answer.

It is true enough that the idea of the Good, understood as
% --- p. 7 ---
\opage{7}\setcounter{footnote}{0}%
the idea of love, of the holy will, is in divine knowledge one with the eternal idea of truth; but from this it by no means follows that the idea of truth in a human knowledge is one with the idea of the Good in human existential willing. It is true enough that in wisdom itself no chasm can be thought between faith and science: but why is God no man of science? Because God is no believer. To think of the divine knowledge as a scientific knowledge, be it philosophical or theological knowledge, is just as blasphemous as to speak of a believing God. Why must a human being apprehend the Absolute in faith? Because the purest and highest development of human knowledge is not divine but scientific. Why must a human being apprehend the Absolute in knowledge? Because faith is a human, not a divine, cognition. Why does the Absolute that is apprehended in human knowledge appear so heterogeneous with the Absolute that is apprehended in faith? Because the relative forms of apprehension grasp the Absolute, not in its central depth, as it must be thought by us \textit{in abstracto} to be, with infinite concretion, in the divine Self, but in two polar extremes, of which the one is designated by universal thought grounded in rational necessity, the other by the punctual life of the will springing up in the conscience. Here there is in the Absolute something I absolutely must believe, because there is in the Absolute something I absolutely must know.

As for the relation of the sciences to one another, the strict subordination is, encyclopedically regarded, really a peculiar matter. Each science is autonomous in its own domain, is a sovereign within its own boundaries. Certainly the one sovereign can subordinate itself to the other, but always with a reservation that secures its independence. Should the special sciences gather in subordination under a highest, solely ruling central science, they would surely begin their greeting with the well-known words:
% --- p. 8 ---
\opage{8}\setcounter{footnote}{0}%
We, who each of us are as good as you and together better than you, acknowledge, etc. I do understand that theology will not set out to rule the other sciences immediately, that it will not, for example, compel either astronomy or geology to, if I may put it so, parrot the Bible; I understand that it will only draw the real, subordinate sciences in under the light of the idea of holy love, of revealed wisdom. But I also understand that if such a theological illumination were really to have scientific significance, then the subordinated sciences would have to be brought to see that by means of supernatural presuppositions a superordinate knowledge is attained, a knowledge that is capable of applying ideal correctives and of scientifically explaining what, from the natural standpoint of real inquiries and exact analyses, appears absurd, self-contradictory, incomprehensible. But such a concession the inquiring sciences surely cannot and will not make.

Since, then, Bishop Martensen in his Homeric-Aristotelian word: ``Let One be ruler!'' has given us his promise to take care ``that a thorough subordination comes into the realm of ideas,'' and in this way to establish a theological universal monarchy, we have, with reservation of our misgivings and our own point of view, to see in what follows on what conditions he will then be able to carry out his plan, how he actually intends to go about establishing the monarchy. With constant regard to the relation between knowledge and will, theoretical and practical cognition, we shall open negotiations with the sciences destined for subordination, in such a way that we draw particular attention to: \emph{the fate of natural science}, \emph{the conditions of history}, and \emph{the position of philosophy in the monarchy}.

% --- p. 9 ---
\opage{9}\setcounter{footnote}{0}%
\section*{I.\\ The Fate of Natural Science in the Monarchy.}
\addcontentsline{toc}{section}{I. The Fate of Natural Science in the Monarchy}
\markboth{The Fate of Natural Science in the Monarchy}{The Fate of Natural Science in the Monarchy}

There is something harsh in the word fate. In using so harsh a word, we must at once ward off a misunderstanding, as if the meaning were that theology, when it now again comes into possession of absolute power, expressly intended to treat natural science harshly. Far from it! How could a science whose principle is love itself and ``the good will'' treat any of its subjects harshly --- that is, when the subject voluntarily subordinates itself: theology will not be a blind fate, it will on the contrary be a gracious providence. But, of course, the subjects must appreciate the grace; for otherwise the gracious providence can easily turn into fate.

The first great principle that must be impressed upon natural science in the monarchy is this: that since the practical is theoretical and the theoretical practical, theory must essentially conform to the will: to a good theory corresponds a good will, to an evil theory an evil will. Without appropriating this principle, natural science will not even be able to distinguish between truth and falsehood. For, says Bishop Martensen (Leilighedsskr. p. 122): ``The opposition between truth and falsehood is at once practical and theoretical, in that, as affirmation or negation of the Good as the perfect \emph{actuality}, it is grounded in the good or evil or at any rate sinful will. The good will has the theory corresponding to the Good. But the evil will, which has stepped out of the dependence of love, must devise a theory for itself, strive in Tantalus fashion to put an imagined actuality in place of the objective one.'' If natural science now cannot agree to make its theories either mediately or immediately dependent on theology's ``good will,''
% --- p. 10 ---
\opage{10}\setcounter{footnote}{0}%
what follows from this? It follows from this that, even if investigations of nature may be permitted when they are understood as subordinate attempts to find one's way in the sensibly actual, the merely logical-physical, the natural sciences themselves are nonetheless to be regarded as spurious sciences, untrue in their innermost essence.

The second great principle that must be impressed upon natural science in the monarchy is this: that the consciousness of culture, with the real sciences emancipated from church doctrine, rests on a false, criminal foundation. It is no use wanting in practice to acknowledge freedom, in theory necessity; in practice to eat of the tree of life, in theory of the tree of knowledge; in practice to assert ``the ethical will,'' in theory the objectively valid law. It is not the fault of actuality, not of the facts, not of the laws, that natural science has stepped out of the dependence of love; no, it is the fault of the will, solely and entirely the fault of the will. For, says Bishop Martensen (Leilighedsskr. pp. 123--124): ``The will wants to have the whole, the all-embracing primacy in consciousness, and if it did not want this, it would not be will. As ethical will it wants to select or reject the theoretical that presents itself to consciousness, to reserve to itself the decisive Yes and No, placet and veto, since it is itself to be responsible and to render account, also for the whole inner household.'' Further (pp. 125--26): ``To eat of the tree of knowledge rests on a will and is an action whose most important consequences are consequences for the personality. And Prometheus, who taught men science and every art and called forth the consciousness of culture on earth, must, chained to the rock, atone for the \emph{action} by which he broke through to his knowledge, because he procured it for himself by a way that was not legitimate, and because he only made men clever but not good --- in this a prototype of the present-day consciousness of culture with its many sufferings, both open and hidden.''
% --- p. 11 ---
\opage{11}\setcounter{footnote}{0}%
If natural science cannot get this principle into its head, what follows from this? Then it naturally follows from this that no true Christian can be a natural scientist, and no natural scientist a true Christian.

But are the natural sciences, then, not actual and true sciences? One might be tempted to take such a question for an untimely joke, a literary indecency, if one did not expressly take account of the statements above. For if we may otherwise call a doctrine that gives us a systematically coherent knowledge of something actual an actual science, then natural science with all its divisions is not only an actual science, but one might be tempted to say the most actual of all actual sciences. Now can any science be thought to be actual without being true? One would not think so. If we separate actuality from truth, then actuality becomes untrue, and truth unactual. Now if an unactual truth is untrue, and an untrue actuality unactual, then it necessarily follows from this that truth and actuality must be united at every stage of actual cognition; and then it follows in turn from this that natural science, as surely as it is an actual science, must be a true science. If natural science is henceforth to keep its character as an actual and true science, this can only happen by its entering into open contradiction with the demand uttered by Bishop Martensen: Let One be ruler, namely the good will.

From the standpoint of natural science, it is so far from the case that our scientific cognitions of truth may conform to our good will, that the honest inquirer on the contrary feels himself ethically obliged to let the will bow to the strict laws, to respect facts and sound logic, and without circumlocution to hold to the logical-physical. ``It is not subjective feelings, not pious wishes, but insight into the essence of things that diverts the lightning bolt and sets the steam carriage
% --- p. 12 ---
\opage{12}\setcounter{footnote}{0}%
in motion. Certainly science is impossible without subjectivity; but subjectivity is in service. Here there is sensing, here feeling, here thinking; but what is merely subjective in sensation and thought must everywhere be purged away as the accidental, the disturbing, the untrue: every presupposition is objective.'' (Phil. Propæd. pp. 74--75).

How will it now actually look, under these ``objective presuppositions,'' when Bishop Martensen, with his Homeric-Aristotelian-dogmatic fundamental thought, appears in the midst of the circle of natural scientists, and with the demand: Let One be ruler! makes the Yes and No of the good will, its placet and veto, the supreme court for the inquiring sciences? Will not those learned in nature, I wonder, be struck dumb at such a surprise, when the earnest man, having addressed such a summons to the assembly, adds an explanation such as this: ``Only when one has ascended from all subordinate ideas to the Good as to the royal principle can there be talk of the living, positive truth as its revelation, in contrast to a merely negative, abstract-logical truth, even if this latter is determined as the unity of knowledge and power. And only when the Good is determined as absolute love, when the opposition between good and evil is recognized as the first fundamental opposition, can there be talk of the living opposition between truth and falsehood, which is something quite other than the merely logical opposition between right and wrong, essence and appearance''? (Leilighedsskr. p. 121).

What natural science in its several branches gets out of its individual inquiries, how the astronomer, for example, goes about the perturbation problems, how the geognost determines his minerals and rock masses, about that theology will not ask. Theology is, as regards facts and their analysis, not exacting, not petty; it does not plague its subjects with negotiations about the
% --- p. 13 ---
\opage{13}\setcounter{footnote}{0}%
particular; it looks to the total. It is thus not for the astronomical details, but for the fundamental astronomical view, not for the geognostic specialties, but for the whole geognostic total view, that the sciences concerned are to render account to theology.

But just here lies the knot. In the inquiring, observing, and calculating special sciences, the apprehension of the whole must necessarily be determined by the apprehension of the particular, and vice versa. What does it really mean, then, that a subordinated special science is to render account to theology's ``good will'' for a total view that it is not to justify by special investigations? It looks strange! Theology seats itself on the throne and says to the subordinated science: Come and render account for your total view! The subordinate approaches and says with a deep bow: ``Your Majesty, here is my account, my analyses, my calculations!'' Theology answers: ``I am not asking, of course, about your accounts and specialties; but you are to vouch to me for a total view that is a fruit, not of the tree of knowledge, but of the tree of life.'' The subordinate excuses itself: ``I cannot get anything into my total view but what follows with rational necessity from my special investigations; have pity on me! My will is not evil; but I cannot do otherwise.'' --- Gracious theology begins to become a little ungracious: ``You offend me with your importunity. You know very well that I will never condescend to understand your accounts.''

Those who think they discover in this silhouette either distortions or exaggerations will perhaps be pleased to tell us wherein these distortions, these exaggerations are supposed to consist. The natural scientists must in any case be able to expect an explanation of how Bishop Martensen has actually conceived the dependence of natural science on ``the royal principle.'' The matter is important; the moment is tense. Will the Bishop really make

% --- p. 14 ---
\opage{14}\setcounter{footnote}{0}%
earnest of carrying through, with ecclesiastical severity, the ecclesiastically strict judgment he has pronounced upon Prometheus and the cultural consciousness of the present age, then the natural sciences all together must prepare themselves for unconditional submission.

Natural science boasts of its great discoveries; but these great discoveries are, after all, fruits of the tree of knowledge, forbidden fruits. Those so radiant lights in the firmament of the sciences, those brilliant stars that encircle names such as Copernicus and Tycho Brahe, Galilei and Rømer, Kepler and Newton, Volta, Ørsted, Laplace with wreaths of immortality --- what are they, really? Fulgurations of the Promethean spark, flashes of flame from that fire of the gods which the Titan brought from heaven to earth. What judgment, what stern judgment must Bishop Martensen, in consequence of his royal principle --- on behalf of the monarchy --- not now find himself compelled to pass upon these sinful productions, these signs of perishable greatness, these deceptive mirages! Those flashes of flame are, after all, flashes of cosmic, not of hypercosmic, of demonic, not of canonical origin: the Promethean spark is a theft from the gods; the fruit of the tree of knowledge is a forbidden fruit, the eating of which has brought a curse upon all nature.

Natural science is, in sap and marrow, in trunk and branches, autonomous, and never gets so far as to become theonomous. Natural science cannot, will not, say \textit{credo, ut intelligam}; it cannot, never could, and never will be able to develop out of any ecclesiastical confession; natural science stands, and must remain standing, as ``a secular science, wholly independent of faith,'' one that contradicts the science of faith, ``flaunting and parading in fullest bloom alongside faith''; natural science --- there is no use trying to hide it --- must, according to the royal principle, be referred to ``the evil will, which has stepped out of the dependence of love,'' and which therefore ``must devise a theory for itself, strive tantalically to put an imagined reality
% --- p. 15 ---
\opage{15}\setcounter{footnote}{0}%
in place of the real one.'' There is, to be sure, something uncannily cutting in such a sentence of condemnation; but no one who is capable of putting two thoughts together will be able to deny that the sentence is an irrefutable consequence of ``the good will,'' the ``royal'' ``principle'' of the sciences, with which Bishop Martensen now means to carry through his Homeric-Aristotelian-dogmatic universal monarchy.

What are the natural sciences now to do: submit to the judgment, or plead lack of jurisdiction? ``They are,'' someone will perhaps say, ``to repent; they are sincerely to renounce `the evil will, which has stepped out of the dependence of love'; they are to consider that when it is said of the Devil, as the father of lies, that he did not abide in the truth, this has as its presupposition that he did not abide in love and obedience''; they are to consider what ``can be seen from the story of man's fall and the two absolutely heterogeneous trees,'' that the opposition between good and evil is the first fundamental opposition, and that there can be talk of ``the living opposition between truth and falsehood'' only when the former is recognized as the first; they are to consider that, as surely as they persist in their one-sided attachment to the logical-physical, without bowing humbly beneath ``the thinking which in sharpness, certainty, etc. is a match for the logical-physical, and whose account of miracle, for example, is a match for the physical objections'': it will go ill with them; then the higher science, the science that has no ungodly knowledge from any tree of knowledge, inasmuch as, where natural science is concerned, it has no knowledge whatsoever, but only the \textit{credo} of godliness from the tree of life, will come upon them with punishment; then the royal principle of the sciences, the good will, in order to forestall ``anarchy and inner dissolution,'' will find itself compelled to cast them out of the community of the sciences. Cast out of the community of the sciences as out of a higher ethical region of light, namely the region
% --- p. 16 ---
\opage{16}\setcounter{footnote}{0}%
of love, they may then, with their ἀσθηνῆ καὶ πτωχὰ στοιχεια (weak and beggarly elements) belonging to a lower knowledge, be permitted to sit in the dark pit of inductions, researches and rummagings, of sublunary researches, of geognostic rummagings, in unscientific weeping and gnashing of teeth.

It is impossible for natural science to rise above the logical-physical; it is impossible for natural science to let the theologians' ``good will'' govern reality and determine truth. ``Should we not'' --- says H. C. Ørsted, otherwise so peaceable, in deep displeasure at the theological-episcopal business of subjectivity --- ``be ashamed in our inmost selves, if we caught ourselves wanting to have a truth other than the real one?'' Natural science cannot --- it is against its constitution --- say \textit{credo ut intelligam}; it cannot bend its neck beneath a believing sole ruler: it cannot, for its part, give way. But the monarchy's ``good will,'' for its part, as surely as it really is a principle, cannot give way either. ``The good will'' must, as a \emph{principle}, be consistent; it must, as ``\emph{the royal principle},'' lay claim to supreme dominion; it must, as the \emph{royal principle of the sciences}, lay claim to \emph{scientific} supreme dominion. Power here stands against power, ``sign against sign'': a classical tragedy is brewing.

The classically tragic is recognized precisely by this, that two ideal, and to that extent equally justified, powers collide in irreconcilable conflict, a conflict that ends with the weaker succumbing to the stronger. Now if it is to be assumed that the royal principle must in the absolute sense be the stronger, then things look black for the natural sciences. A storm of judgment stands on the horizon; it is the old tragedy of tragedies, the tragedy of ``Prometheus Bound,'' that begins all over again. Once more must Prometheus, ``who has taught men science and every art and called forth
% --- p. 17 ---
\opage{17}\setcounter{footnote}{0}%
cultural consciousness on earth, chained to the rock, atone for the act by which he broke through to his knowledge, because he obtained it by a way that was not legitimate'';
% Errata: p. 17, l. 4 from top: legitim; read: legitim`` (closing quote added in the English)
once more: O cruel fate! And we who, with the ``cultural consciousness'' of the nineteenth century, were already dreaming that we had seen Prometheus unbound, seen the chains fall from the Titan's mighty limbs, seen the gaze he fixed upon the dying kite, seen Heracles with his bow, radiant in so divine a light that night itself, the night of brutalities, had to flee and hide in the caves of the rocks. The dream is over; alas, it was only a dream, an imagined reality, a conjuring trick of the evil will. For, says Bishop Martensen: ``the evil will, which has stepped out of the dependence of love, must devise a theory for itself, strive tantalically to put an imagined reality in place of the objective one.''

In the lofty classical tragedy everything is symbolic. Prometheus and Io, Kratos and Bia, Hermes, Okeanos --- all the figures appearing in the tragedy are symbolic. Prometheus, the inventor of the arts, the father of the natural sciences, is cultural consciousness, the autonomous consciousness which, by excluding faith, has become ``too absolute in knowledge''; while Io, the unhappy beloved of Zeus, a subjectivity, if we may put it so, run wild, stung by the gadfly of paradox, is the consciousness which, out of misunderstood theonomy, has excluded knowledge and has become ``too absolute in faith.'' Kratos and Bia, who are sent out by Zeus to bind Prometheus, are indeed, by the words, \emph{Violence} and \emph{Power}; but in the symbol it is \emph{love and the good will} that are sent out by the royal principle to forestall ``anarchy and inner dissolution'' in the monarchy.

Only one change will take place in the old tragedy, a change which is then to that extent also symbolic, inasmuch as it bears witness that ``the system of laws
% --- p. 18 ---
\opage{18}\setcounter{footnote}{0}%
and forces which we call nature is engaged in a teleological development, a continued creation''; the change is characteristic: Prometheus, namely, will this time not be chained to a rock, and for the good reason that the royal principle does not own a single rock in the whole of natural reality; instead he will be bound to a tree, one of the two ``absolutely heterogeneous trees'' in Paradise, namely the tree whose fruits were forbidden to man.

The new penal law, however, will be just as harsh as the old. For although the Titan is bound to the tree of knowledge, that is, to the tree whose fruits taste best to him, ``the good will,'' the will which alone has ``the theory corresponding to the good,'' has of course seen to it that he is bound with his back to it, so that in his long term of punishment he will not grow fat on ``the forbidden eating,'' but must make do with ``the meager fare of propaedeutics,'' and will thus surely come to feel ``love.''

To be sure, natural science stands, according to the monarchy's order of rank, on a rather subordinate step among the sciences; for, says Bishop Martensen (occasional piece, p. 109): ``Natural science, as an empirical science, has no other grounds to rely on than induction and analogy. From a series of phenomena that, under the same conditions, continually recur as the same, one derives what one calls a law; but it is easily seen that unconditional universal validity and necessity cannot possibly come about by this route.'' But for natural scientists the matter looks otherwise. ``In my book, Aanden i Naturen,'' says H. C. Ørsted, ``I have with full conviction paid homage to the old doctrine of the immutability of the laws of nature.'' This ``old doctrine'' is so decisive for the researcher that the scientific character of research stands and falls with its validity.

We do not, of course, doubt that Bishop Martensen, with the help of the royal principle together with \textit{credo ut in\-%
% --- p. 19 ---
\opage{19}\setcounter{footnote}{0}%
telligam}, will be able to furnish a rigorous proof that the laws of nature are not immutable. ``In general,'' he says (p. 109), ``natural science, as merely empirical science, knows nothing whatsoever about what is possible or impossible in the \emph{absolute} sense. And with regard to reality and the facts, it moves only in a middle, within a circle of conditions that happen to be factually given, it knows not how, while beginning and end are shrouded in impenetrable darkness. It still owes us the answer to the old question: which came first, the egg or the hen? The acorn or the oak?'' We will now go along with this point of view; we will even concede that theology is the only science which, from its supernatural standpoint, does not owe the answer when it is seriously asked how matters really stand with creation and nature in the oak and the hen. Indeed, it shall not even be doubted that dogmatics, by determining the right relation between creationism and traducianism, will arrive at the incontestable result that originally we have the creationistically natural in the oak, the traducianistically creaturely in the acorn, the creature (κτισις) in the hen, nature (φυσις) in the egg. But all this presupposed, there still remains the question, fateful for the fate of natural science in the monarchy: how will theology go about convincing natural scientists that the doctrines of natural science belong to the ἀσθηνη καὶ πτωχα στοιχεια (weak and beggarly elements)?

Natural scientists are proud of their science. These men, little familiar with the dogmatic story of man's fall and ``the two absolutely heterogeneous trees,'' will surely never get it into their heads that the question of the immutability of the laws of nature is to be decided in theology. Call it a whim, an idée fixe, an inspiration of the Devil, an effect of the evil will that has stepped out of the dependence
% --- p. 20 ---
\opage{20}\setcounter{footnote}{0}%
of love, or whatever one likes; but certain it is that every one of the researchers is sure of the reality of his science. An astronomer, for example, holds himself so fully assured that his conception of the starry heavens is scientific that, if the royal principle should find occasion once more to raise doubts about the fundamental thought of the Copernican system, he will be bold enough to declare such a doubt unscientific. No one is more willing than the geologist to concede that, with respect to the developmental history of the earth, much, infinitely much, still remains to be investigated; but that geology should not be a science, a real and true science --- that you will never get him to concede; indeed, the geologist in his scientific zeal is even capable of venturing so far that, if someone asks him about the region in which Paradise actually lay, he will refer the questioner to the nursery. What has been said here of the astronomer and the geologist holds correspondingly of the botanist. Just set a botanist the ``higher scientific'' task of investigating, together with the geologist, in order to secure the ``\emph{objective} reality,'' the nature of the soil in which ``the two \emph{absolutely heterogeneous} trees'' must be assumed to have grown --- whether both grew, for example, in meadow soil or in arable land, in clayey or in sandy earth, or whether the soil from which the tree of life sprang up was perhaps \emph{absolutely heterogeneous} with that in which the tree of knowledge struck root: he will open his eyes wide. If one declares that his science is a ``weak and beggarly elementary teaching,'' since he cannot even determine for us the place of the two \emph{absolutely heterogeneous} trees in the system, but owes us the answer to the question whether both belong, for example, to the Gramineae, to the Rosaceae, or perhaps to the pink family: then the otherwise so good-natured botanist is even capable of becoming rude toward ``the good will'' and forgetting himself to the point of most humbly begging
% --- p. 21 ---
\opage{21}\setcounter{footnote}{0}%
``the royal principle,'' together with the whole dogmatic tree-school, to go to the Devil. That such an outburst is highly improper we must concede; but such things can indeed happen in despair. Natural scientists become so desperate when a dogmatician, from the supernatural height of theonomy, looks down with superiority upon the natural autonomy of their science.

It will be a hard fight for the ``royal principle'' before it gets ``naturalism'' thoroughly rooted out of natural science and the researchers convinced that there really can be a science whose fundamental cognitions and first presuppositions are in open conflict with what they must declare to be notorious facts, a science whose ``account of miracle is a match for the physical objections,'' a science which, on the ground of ``the good will,'' makes it objectively irrefutable that creation ends with a paradise, that before the fall there was no trace of corruptibility, pain and death on earth. What these empiricists, the geologists in particular, will not think up to give their ``weak and beggarly elementary teaching'' some kind of semblance of real and true science! Instead of beginning what they call the history of the globe with the right historical, namely the suprahistorical, beginning, with Paradise, ``the two absolutely heterogeneous trees'' and ``the forbidden eating,'' they begin --- just think! --- with the greywacke! that is, with ``devising a theory for themselves, striving tantalically to put an imagined reality in place of the objective one.''

By a series of uncanny pictures, little edifying for the dogmatician, the geologists of the present age seek to make vivid a reality, by no means objective but tantalically imagined, of an immeasurable primeval age, in that --- of course only at random, for their own ``and the ignorant's amusement'' --- they strive in naturalistic fashion to objectify for themselves animal life in the various periods of the earth, from the greywacke and the carboniferous limestone right up to the Chalk and Molasse periods, always under the surreptitious presupposition
% --- p. 22 ---
\opage{22}\setcounter{footnote}{0}%
that even the Chalk and Molasse periods precede the fall. To form a concept of the naturalism that can reveal itself in such tantalic objectifications, we need only cast a glance at one of these pictures, the picture, for example, that shows us ``the inhabitants of the Jurassic sea and its shores in full activity.'' How tantalically imagined is this activity! ``In the air,'' says the geologist,\footnote{Geologiske Billeder af Bernhard Cotta. Efter Originalens tredie Oplag ved C. Fogh. Copenhagen 1859. p. 243 ff.} ``one sees two great flying lizards of the genus Pterodactylus, perhaps two jealous rivals who in the heat of battle have lost sight of the prize, for the poor female pterodactyl is just falling prey to a long-necked, famished Plesiosaurus, which has got hold of one of her wings and will soon make an end of her in spite of her anxious cries. Millions of years ago her cries of terror died away in the air,'' etc.

That the whole conception of the Jurassic period in such a picture must rest on an enormous illusion of the senses, dogmatics can, in consequence of the royal principle, expressly prove. For it can be scientifically demonstrated, by fruits gathered not from the tree of knowledge but from the tree of life, that ``naturalism'' is here guilty of an enormous hysteron-proteron; for it can be proved that such a struggle among animals, which can take place only in a nature struck by the Lord's curse, cannot possibly have taken place before the fall. Here the battle is joined; here is precisely naturalism's Achilles' heel, one of the weak points from which the science of faith will be able victoriously to attack the elementary teaching of unbelief.

The proof is, by \textit{credo, ut intelligam}, irrefutable, and can be grasped by anyone capable of following its train of thought. Before the fall, all nature in its innocence was theonomous; only after the fall --- this can be seen from ``the Christian dogmatics'' --- did it become, by means of the serpent and the
% --- p. 23 ---
\opage{23}\setcounter{footnote}{0}%
cosmic principle, autonomous. But now look at whatever parts of that Jurassic picture you like: in all the scenes, on all sides, facts appear which testify that autonomy has already displaced theonomy, that what exists is being badly governed. See there, farther forward in the picture, a mighty crocodile lies! The mere presence of a crocodile forebodes how what exists is governed here. The crocodile lies, as one sees, under bananas, cycads and tree ferns; it too, I should think, will not make do with ``the meager fare of propaedeutics,'' but is evidently lying in wait for a tasty dessert. That such a dessert, however tasty it may otherwise be, must absolutely fall not \emph{before} the fall but \emph{after} the fall, since, looked at more closely, it is a sinful, namely an objectively-naturally sinful dessert, can be derived exegetically from the context. For from what immediately follows it is seen that a well-armored sea turtle has --- of course autonomously, not theonomously --- forestalled the crocodile and has seized a splendid belemnite creature, which it is on the point of dragging ashore. Everywhere traces of corruptibility, traces of sin, of curse! Farther along an Ichtyosaurus swallows --- and this is supposed to have happened before the fall? --- a whole cuttlefish, while another cuttlefish --- of course after the curse --- together with a short-tailed crab seeks safety in hasty flight. If, we say, all this is reality; if it, as our naturalists after all maintain, really once took place, and if it, as already shown, cannot possibly have taken place \emph{before} the fall: then it must necessarily have taken place \emph{after} the fall. \textit{Quod erat demonstrandum}: here is the proof.

From the supernatural standpoint of theonomy, theology is doubtless in the right; but how will the royal principle go about getting natural scientists lifted up to the standpoint of theonomy? It is, in earnest, a great and important step that Bishop Martensen has on this occasion
% --- p. 24 ---
\opage{24}\setcounter{footnote}{0}%
resolved upon, in that, by pronouncing the judgment and breaking the staff over the Promethean presuppositions of natural science, he has challenged ``cultural consciousness'' to a fight to the death. Whatever the outcome may be, it shall be granted him that by his appearance on the battlefield he has, if not immediately, then at least mediately, furnished a contribution that lays bare the innermost kernel of the problem.

We too, for our part, have long been aware of the ember smoldering beneath the ashes, and have waited with mounting suspense for the moment when it would burst into flame. We have therefore striven as best we could to get clear, not merely about the relation between faith and knowledge in speculative generality, but --- so that we should not miss concrete reality --- expressly about the relation between religion and natural science. We have not, it goes without saying, been able to soar as high as Bishop Martensen, and have come with our investigations to more meek results.

When, then, we asked ourselves: how is it possible that two principles as heterogeneous as faith and knowledge can be united in one consciousness, we were thinking here as much of the natural-scientific as of the philosophical consciousness. If the natural scientist cannot unite faith and knowledge in one consciousness, then neither can the philosopher; if the natural scientist can, then the philosopher can too. That theologians and clergymen are easily drilled into uniting faith and knowledge we can well understand; but with the natural scientist it is another matter. In developing the problem we have kept our eye fixed on the difficulties of natural science, and for good reasons. In the first place, we find in the natural scientist what one does not always find in theologians and clergymen: acquaintance with, and therefore respect for, a real nature with real facts and real laws. In the second place, we can with fair certainty presuppose of the natural scientist what one may not so readily presuppose of theologians and
% --- p. 25 ---
\opage{25}\setcounter{footnote}{0}%
clergymen, that if he really does occupy himself with thinking about religious problems, his thinking must have a personal originality of its own; it then rests not on official drill but springs from an irresistible inner need.

Looking back on the progress of the natural sciences and the course of culture in the last two centuries, we believed we had discovered, as regards religious natural scientists' conception of faith and knowledge, an unmistakable striving to let the separation between a religious and a scientific, a practical and a theoretical cognition come into force. We believed we found such a separation, not indeed theoretically developed, but nonetheless factually acknowledged, for example by Newton. Newton's mechanics bears as little trace of faith's supernatural ``good will'' as d'Alembert's or Laplace's; and yet Newton was, as is well known, a religious man, a man who wrote notes on the prophet Daniel, expounded the Apocalypse, and believed in miracles and prophecies.\footnote{For Newton the prophetic word has immediate divine authority. \textit{Auctoritas Imperatorum, Regum, Principum humana est: Conciliorum, Episcoporum auctoritas humana. Prophetarum auctoritas divina est \dots\ Isaaci Newtoni Opuscula Lausannæ et Genevæ MDCCVLIV. Observationes ad Danielis prophetæ vaticinia p. 291.} As regards the interpretation of the prophecies, Newton does indeed warn against going too far into particulars, but nevertheless adds, work cited, \textit{p. 292: Rerum futurarum prædictiones Ecclesiæ conditionem per omnia secula respiciunt.} How immediately he understood the apocalyptic visions is seen from the following statement: \textit{Post paratas lucernas, vidit Joannes apereri januam Templi ac voce, quasi buccinæ, evocabatur ad eam magni atrii portam, quæ Orientem spectat, ut illic visa videret. Opuscul. XXV. Observat. in Apocalyp. p. 451.} It is clear from these and similar statements that Newton must on the whole regard cognitions of the natural and cognitions of the supernatural as just as heterogeneous as the natural and the supernatural. Now if cognition of the natural falls under science, and if there cannot be absolutely heterogeneous sciences, then it follows that cognition of the supernatural must be a cognition for practical use, a practical cognition.} Had,
% --- p. 26 ---
\opage{26}\setcounter{footnote}{0}%
however, any physicist, while Newton was still in full vigor, ventured to attack the \textit{Principia philosophiæ naturalis mathematica} with arguments drawn from what God is supposed to have said in his Word, then the God-fearing Newton would surely not have hesitated to repel so foolish an attempt with words that could easily be construed to the same teaching as the serpent's words: ``It is not true, what God has said (namely, from the standpoint of science).''

If in our day a biblical theologian were to counter optics' explanation of the formation of the rainbow by citing God's word to Noah (Gen. 9:13): ``I will set my bow in the cloud, and it shall be a token of a covenant between me and the earth,'' and were to add that the rainbow is expressly due to a miracle, according to what God himself has said: would not even the most God-fearing of all scientific opticians give an answer that could be interpreted as the teaching: ``It is not true, what God has said (namely, from the standpoint of science)''?

Let a dogmatician who, with the royal principle, knows so excellently how to defend the possibility of the miracle, just try to explain to God-fearing zoologists the possibility of the miracle of Balaam's she-ass (Num. 22:21 ff.)! ``And Balaam rose up in the morning and saddled his she-ass and went with the princes of Moab. \dots\ But the angel of Jehovah stood in a hollow way between the vineyards. \dots\ And when the she-ass saw the angel of Jehovah, she lay down under Balaam; then Balaam was angry and struck the she-ass with the staff. Then Jehovah opened the mouth of the she-ass, and she said to Balaam:

% --- p. 27 ---
\opage{27}\setcounter{footnote}{0}%
what have I done unto thee, that thou hast smitten me now three times. \dots{} am I not thine ass, upon which thou hast ridden from the time thou wast, unto this day,'' etc. How, now, will the dogmatician with the royal principle go about presenting this miracle in such a way that the presentation ``outweighs the physical objections''? If he appeals to what the angel said, the zoologists answer as with one mouth: ``it is not true, what the angel said (namely, from the standpoint of science).'' If he appeals to what the ass said, they answer in chorus: ``it is not true, what the ass said (namely, from the standpoint of science).'' If he appeals to what God has said in his Word, they keep on: ``it is not true, what God has said (namely, from the standpoint of science).''

What the natural scientist objects against all the particular miracles in the whole of biblical history, he objects also against the fundamental miracles: creation, incarnation, and inspiration. ``It is not true that God in the beginning created the heaven and the earth (namely, from the standpoint of science); it is not true that God's only-begotten Son was conceived by the Holy Spirit and born of the Virgin Mary (namely, from the standpoint of science); it is not true that the Holy Spirit came at the feast of Pentecost in a storm and in flames of fire, and in a single instant taught the apostles the many foreign tongues (namely, from the standpoint of science).''

In order not to rush into any final, irrevocable judgment on these pronouncements, ``the good will'' will doubtless do best to repeat to itself what sort of standpoint it is from which one speaks here: that it is constantly from the standpoint of \emph{science}, in particular of \emph{natural science}. For on the standpoint of natural science it is agreed---something that is not always the custom and practice among theologians and clergymen---to keep strictly to the separation between \emph{what one knows} and \emph{what one does not know}. It lies in the nature of the case
% --- p. 28 ---
\opage{28}\setcounter{footnote}{0}%
that whoever knows something thoroughly will take note of the boundary between knowing and not-knowing. By reflecting on the limits of human knowledge, a natural scientist will readily arrive at the following result, stated in the philosophical Propædeutik (p. 67): ``By virtue of the inner, original connection of empirical research with the a priori, there is an essential knowledge; but in consequence of the heterogeneity of thinking and sensing, real cognition is subject to finite conditions, and all our knowledge is limited. The limit of knowledge is a double one: a \emph{relative} one, insofar as there is always something here that we do not yet know, but which we can nevertheless come to know in time; an \emph{absolute} one, insofar as a final separation of the content of consciousness and of actuality is grounded in our organization: here, then, is something that we shall never come to know in this existence.''

From the relation between the relative and the absolute in the limitation of our knowledge, the relation between faith and knowledge can be cleared up. Had human knowledge only relative limits, then of course a ``\textit{plus ultra}'' would always have to be possible; then all our cognitions without distinction would in the end have to be grounded in science. What science could not ground on these terms, what it found self-contradictory and objectionable, would then be not merely untrue from the standpoint of science, but by that very fact would have to be regarded as untrue in itself: in this sense it would quite rightly be an impossibility that something could be true in faith which was not true in knowledge. If, on the other hand, human knowledge has, besides its relative limits, an absolute limit, so that here there is something that we shall never come to know in this existence: then here is a mystery, a riddle of existence, whose solution in science must turn out untrue. To this mystery must be assigned the absolutely ideal presuppositions of personal life: that a holy, almighty will is the author of all existence; that man's essence, ethically regarded,
% --- p. 29 ---
\opage{29}\setcounter{footnote}{0}%
finds itself in an original conflict with its ideal; that every individual who opens his eyes in the kingdom of the spirit must acknowledge himself \emph{guilty}; that the heart can find peace only when the guilt has been perceived, confessed, and atoned, and certainty of an eternal life in the kingdom of personality and truth has been obtained. The more conscience sharpens these ideal presuppositions, the deeper is the religious need; the deeper, stronger, and more inward the religious need is, the more absolutely personal is the decision; and the more absolutely personal the decision is, the more clearly it is seen that the absolutely ideal presuppositions of personal life must be had in faith, not in knowledge.

But why, then, not \emph{both} in faith \emph{and} in knowledge, \emph{first} in faith \emph{and then} in knowledge, first \textit{credo} and then \textit{intelligam}? Yes, why not? As a rule it is easier to make the opposition between faith and knowledge evident to natural scientists and religious individualities than to theologians and speculative clergymen. The religious individuality \emph{feels} what faith is; the natural scientist \emph{knows} what knowledge is; but those whose faith has an aftertaste of knowledge do not know what \emph{faith} is, and those whose knowledge tastes of faith to an incredible degree do not know what \emph{knowledge} is.

For the natural scientist faith is absolutely heterogeneous with knowledge; in this he has both a temptation to give up faith and a possibility of preserving it. How he stands firm in this temptation, how he personally makes use of this possibility, for that he must render account to ethics in his conscience. In this respect we can subscribe to Bishop Martensen's words, that the will ethically will ``choose or reject the theoretical that presents itself to consciousness, reserve to itself the decisive Yes and No, Placet and Veto, since it must itself be responsible and render account, also for the whole inner household.'' But let us now consider in particular, first the temptation and then the possibility. The temptation is to make faith and knowledge somewhat homogeneous and to let the cognition
% --- p. 30 ---
\opage{30}\setcounter{footnote}{0}%
of the absolute limit of knowledge recede into the background. That faith is made homogeneous with knowledge is recognized by this: that consequences for faith are derived from the presuppositions of knowledge, and consequences for knowledge from the presuppositions of faith. As often as the natural scientist is caught in this temptation, he notices at once how unreasonable faith is, how little it can be reconciled with knowledge.

We choose as an example the relation between nature and creation. One will strive in vain to draw from the presuppositions of faith consequences for the benefit of science; for the miracle of creation abolishes the concept of nature held by science. One will strive in vain to develop from the presuppositions of knowledge consequences for the benefit of faith; for the concept of nature held by knowledge makes the miracle of creation into a fantasy. ``To apply the thought of creation to an understanding of what is created means: to see the whole world under the double point of view: \emph{nothing and yet really something}. The concept in which this contradiction is resolved is the concept of omnipotence; the created is in itself a nothing, but at the bidding of omnipotence it has become something.'' Herewith the presupposition of faith is posited, but the scientific concept of nature is abolished; for how is one here to think the transition from creation to natural production? ``A relation that has begun with a miracle must also be continued by a miracle: not merely for a few moments, not merely, if we may say so, for some eight to fourteen days after the creation, but for centuries, for millennia, for eternities of eternities, created things must continue to be in themselves what they originally are and have been, a nothing, a pure nothing for the Almighty.'' Insofar as water is created, it has not merely a miraculous origin but a miraculous being; at the bidding of omnipotence it is a real something, and yet for omnipotence itself, the almighty will, a nothing: \emph{there is no nature in the created water}. Insofar as water has a natural existence, insofar as there is,
% --- p. 31 ---
\opage{31}\setcounter{footnote}{0}%
in other words, nature in the water, it comes about through the chemical combination of oxygen and hydrogen, and thus has not only a natural existence but a natural origin; \emph{but the natural origin is, of course, no creation}.

As the existence is, so is the origin: the scientific concept of nature abolishes creation; as the origin is, so is the existence: creation abolishes the concept of nature. ``If Adam is created, then Adam is no nature, then there is no nature in Adam. If there is no nature in Adam and Eve, then the children of Adam and Eve cannot be \emph{born} either; they must all be \emph{created}.'' (Om Theologiens Naturbegreb med særligt Hensyn til Malebranche). If there is a real human nature, hence \emph{nature}, in the descendants of Adam and Eve, then there has also been a real human nature, hence \emph{nature}, in Adam and Eve; then Adam and Eve have been real natural individuals, have had a real natural existence. But if Adam and Eve have had a natural existence, then they have also had an origin just as natural as any of their descendants. From the miraculous the natural cannot be derived; from the natural the miraculous cannot be derived: the miraculous is absolutely heterogeneous with the natural.

From this point of view a believing individuality will easily be tempted to break the staff over natural science in order to preserve faith, while the natural scientist, conversely, feels tempted to give up faith in order to maintain science. The temptation is that faith and knowledge are so easily made a little homogeneous, although they are absolutely heterogeneous. If the absolute in the heterogeneity is stressed sharply enough, a glaring contradiction is heard: \emph{what is truth in faith is untruth in science; what is truth in science is untruth in faith}. But in the contradiction itself lies the possibility of its resolution.

It is the absolute limit of our knowledge, it is the riddle of
% --- p. 32 ---
\opage{32}\setcounter{footnote}{0}%
existence, on which the attention of the knower as well as of the believer must here be fixed. If all human knowledge has an absolute limit, if there really is here a mystery of the spirit, a riddle of existence, then there is also here something that must be \emph{believed}, because it cannot be \emph{known}. Now if what is hidden in the riddle, that which is solely and exclusively destined to be an object of faith, is nonetheless drawn in under knowledge and science, what then? Then knowledge oversteps its absolute limit. And when knowledge has overstepped its absolute limit, what then? Then science blithely sets about fathoming, grounding, and comprehending that which by its nature neither can nor shall be scientifically comprehended. But when science busies itself with comprehending the incomprehensible, what then? Then it necessarily arrives at contradictory concepts. What can only be expressed in contradictory concepts must be untrue; the mysteries of faith can only be expressed in contradictory concepts: therefore they are untrue. That something can be true in faith and untrue in science now follows directly from this, that \emph{science, by wanting to comprehend what lies outside its limit, itself becomes untrue}.

But how can something that is really true in science be untrue in faith? Answer: Because it is not only knowledge but also faith that has an absolute limit: that knowledge has faith as its limit necessarily presupposes that faith has knowledge as its limit. But that faith has knowledge as its limit means that a concept of faith cannot possibly be developed into a concept of knowledge. In relying on the miracle, faith presupposes a real nature with real laws. But how does nature relate to the miracle? Why, that is precisely what is the riddle of existence. If now faith, forgetting its absolute presupposition, namely the riddle, oversteps its limit and begins to want to overhaul the scientific concept of nature for the benefit of the miracle: then
% --- p. 33 ---
\opage{33}\setcounter{footnote}{0}%
the concept of nature held by natural science becomes not merely an untruth in faith, but \emph{faith, having engaged in something it is not equal to, itself becomes untrue}.

That faith and knowledge can be united in one consciousness, not \emph{although} but \emph{because} they are absolutely heterogeneous principles, means, in other words, that they can be united only in an \emph{absolute mystery}. The absolute mystery is at once an exclusive boundary and a negative middle. As exclusive boundary the mystery sets a separation between faith and knowledge as between two regions. Each region has its unconditioned, its center: the center of faith is freedom, the almighty, holy will; the center of knowledge is necessity, the rational necessity revealing itself in unchangeable laws. Now when the believer looks toward the center of faith, and the knower toward the center of knowledge, the extremes are stretched to the utmost: what is true in faith, namely that the unconditioned is the miracle, is untrue in knowledge; what is true in knowledge, namely that the unconditioned must be sought in the necessity of reason, is untrue in faith. As negative middle, on the other hand, the mystery unites the separated extremes and shows itself as the medium of passage between the regions: through the mystery goes the way from faith to knowledge, from knowledge to faith. When the believer and the knower both at once direct their gaze toward the mystery, the conflict is settled and the extremes are reconciled. Seen from the standpoint of knowledge, the miracle is now possible. For though it is indeed impossible that an arbitrary breach of the eternal laws of nature, identical with the necessity of reason, should ever be able to occur, yet, when the intermediate determinations hidden in the mystery are understood as implied, the revelation of the power of will, wisdom, and love in which the believer sees the miracle can very well be thought to be consistent with the unchangeability of the laws. Seen from the standpoint of faith, the necessity of reason, with the unchangeable laws of nature, can be reconciled with the presupposition of the miracle. For though it is indeed
% --- p. 34 ---
\opage{34}\setcounter{footnote}{0}%
impossible that necessity should be able to hamper freedom, and the laws of nature restrict God's unconditioned will, yet, when the intermediate determinations hidden in the mystery are understood as implied, it can well be thought that the laws of nature, far from being broken, are on the contrary fulfilled in the very highest sense every time a divine act of will is accomplished.

From this reconciliation it follows:

1) \emph{that a science of faith (theology) is unthinkable.}

Were a science of faith to be thinkable, then concepts of faith would have to be transformable into concepts of knowledge, and vice versa; but, owing to the absolutely heterogeneous characteristic mark of these concepts, that is impossible. The concept of knowledge is marked by necessity, the rational necessity in which the determination of freedom is---on account of the mystery---a vanishing moment; the concept of faith, on the other hand, is marked by freedom, the divine, overreaching act of will, in which the determination of necessity---on account of the mystery---incessantly vanishes. The concepts creation, paradise, the fall, and the like are concepts of faith, not concepts of knowledge; whereas the concepts graywacke, Jurassic period, Cretaceous period are concepts of knowledge, not concepts of faith. The kind of reality that must be ascribed to paradise and the fall is no more to be sought in the Jura and the chalk than the reality that must be ascribed to the Jura and the chalk is to be sought in paradise and the fall. Graywacke, Jura, and chalk belong to another region than creation, paradise, and the fall. From this point of view it can be judged how much weight can be laid on Bishop Martensen's words when he says (p. 133) of my philosophy: ``With him'' (Prof. N.) ``there stands a so-called modern consciousness,''---about graywacke---``there stands a worldly knowledge, wholly independent of faith and''---in the imagination of the person concerned---``contradicting faith''---about animal life
% --- p. 35 ---
\opage{35}\setcounter{footnote}{0}%
in the Jurassic period---``swaggering and flaunting in fullest bloom alongside faith.''

2) \emph{that no concordat can be thought between faith and knowledge.}

A concordat, that is, a finite, arbitrary agreement, can only be thought between relatively homogeneous powers which restrict one another in an external and finite way: restriction presupposes a limitation coming from outside. But faith and knowledge cannot possibly---on account of the mystery---restrict one another. There is indeed in faith a limit for knowledge, insofar as what can be believed cannot be known; but this limit is one with knowledge's own immanent limit. With its immanent limit, its inner limitation, knowledge is wholly unrestricted over against faith, as unrestricted as if there were no faith at all; likewise faith, with its immanent limitation by knowledge, is unrestricted over against knowledge, as unrestricted as if there were no knowledge at all. Just as little as knowledge is ever able to attack faith, just so little will faith ever be able to attack knowledge; just as little as knowledge is able to conform itself to faith, just as little is faith able to conform itself to knowledge. If, then, they can---which the absolute mystery absolutely prevents---neither contend against one another nor support one another, neither benefit nor harm one another, then surely a finite agreement between them is unthinkable.\footnote{Bishop Martensen's discovery of the concordat and his fantasizing about the many concordats will be dealt with in the third part.}

3) \emph{that the separation of faith and knowledge cannot halve human consciousness.}

Consciousness is halved by relatively homogeneous, but not by absolutely heterogeneous factors. When a consciousness is halved
% --- p. 36 ---
\opage{36}\setcounter{footnote}{0}%
in thinking, it gets half a thought instead of a whole one: what is lacking is thought; if it is halved in will, it gets half a will instead of a whole one; what is lacking is will. On the other hand, one cannot in a stricter sense say that a consciousness is halved by thinking and willing in such a way that the half of the thought becomes will, the half of the will thought; for thought can only in a highly improper sense be transformed into will, and will into thought. If, then, it stands firm that faith can in no way be transformed into knowledge, or knowledge into faith, then it also stands firm that faith can just as impossibly be halved by knowledge as knowledge by faith. If a consciousness cannot be halved by friendship and mathematics, by love and chemistry, then still less can it be halved by ethics and astronomy, by religion and geology. There are mathematical friends, but there is no friendly mathematics; there are enamored chemists, but no chemical enamorment, no enamored chemistry. There can be ethical and unethical astronomers, but astronomy is neither ethical nor unethical. Science, in other words, is neither ethical nor unethical, neither Christian nor un-Christian, neither believing nor unbelieving: love is practical, but science is theoretical. The bugbear quality in Bishop Martensen's pathos (p. 79): ``Thou shalt only practically love and acknowledge me, but theoretically thou shalt''---which is impossible---``deny me! Practically thou shalt hold to the faith which is delivered to my congregation; but theoretically thou shalt''---by virtue of a theological fancy---``hold to atheists and antichrists, shalt walk in the way of deniers and sit in the seat of the scornful,''---thus falls away of itself.

Otherwise, quite otherwise, will natural science, and with it the natural scientist, be treated in the universal monarchy. Here it is of no use that the natural scientist wants to have his religion in a practical cognition, and his science in a theoretical cognition independent of religion; for Bishop Martensen demands that theoretical
% --- p. 37 ---
\opage{37}\setcounter{footnote}{0}%
cognition shall coincide with practical, the practical with the theoretical. In the monarchy it is not enough that the astronomer himself is a believer; no, he must be a believer in consequence of a believing astronomy; not enough that the geologist personally wants to be a Christian; no, he must be a Christian in the practical on the ground of a Christian geology in the theoretical. For, says Bishop Martensen (p. 133): ``It is of no use at all that one, when one gets into difficulty, gives assurance that the practical, that life, nevertheless has primacy over science. For of such assurances it must be said that they have the appearance of life but deny its power, so long as a heathen, puffed-up knowledge, instead of being cast off, is to be conserved, developed and fostered, tended, protected and nourished alongside faith in one and the same consciousness.'' In the monarchy it is not enough that a conscientious natural scientist is willing to render account to religion and ethics for his personal relation to the impersonal truths of science; that he will no more introduce unbelief than belief, no more heathenism than Christianity, into his scientific knowledge; that the will, therefore, with which he honestly and uprightly devotes himself to the cultivation of science is ethically determined by the will with which he personally nourishes and develops faith in life. No, more is required in the monarchy. The will with which a natural scientist cultivates his science must be so believing a will that it can make science itself believing. If his will cannot bring it so far as to make his science believing: then, of course, he gets two wills, one in science and another in faith; he then wills ``not one thing but two, and has a double aim.'' But woe to the natural scientist who in the monarchy ``wills not one thing, but two, and has a double aim''! It will go as unhappily for him with his natural science as it has gone for us with ``our newest philosophy''; and it is, alas, now a well-known matter that it has gone very unhappily for us. For, says Bishop Martensen (p. 139): ``As
% --- p. 38 ---
\opage{38}\setcounter{footnote}{0}%
knowing, it'' (our newest philosophy) ``wills man's emancipation from Christianity, indeed from the living God. It thus has its home in the camp of unbelief, from which it proclaims that, over against modern science and the cultural consciousness built upon the natural sciences, Christianity has become a powerless and untenable matter. As \emph{believing}, it wills the relation of obedience and authority, and maintains the complete independence of faith from science, though only as an absolute absurdity---a religiosity which must be characterized as obscurantism, but which differs from other obscurantism in this, that from first to last it is \emph{theorized}. How the two wills go together into one has not been made comprehensible, and by the nature of the case never will be.''

That on this occasion we press so close upon natural science with our philosophy is by no means done for the sake of philosophy, for philosophy must defend itself; no, it is done only for our own sake. It is, of course, natural that we, who in our position toward the monarchy are so unhappy, should cherish a heartfelt sympathy for other, highly esteemed fellow sufferers. It is out of sympathy and concern that, in conclusion, we raise the question: What is to become of the unhappy natural scientists? We are not speaking here of all natural scientists, but only of the unhappy ones. Which, then, are the unhappy natural scientists? Unhappy we call those natural scientists who want to cultivate not only their science but also their God. Why are they unhappy? Because, like ``our newest philosophy,'' they will not one thing, but two, and thus have two wills. This ``the royal principle'' can in no way permit. That Christ in the ``Dogmatik'' is allowed to have two wills is another matter; that the principle can justify, for it is comprehensible. ``The opposition between the natural and the spiritual, the cosmic and the holy principle'' must of course ``necessarily be found in the second Adam as a duality in his will; since the second Adam cannot
% --- p. 39 ---
\opage{39}\setcounter{footnote}{0}%
be thought monothelitically (which would be to think him monophysitically), but must be thought dyothelitically: so there also stirs in the second Adam the principle that makes the fall of the first \emph{possible}.'' It is otherwise with an unhappy natural scientist. To be sure, there also stirs in the natural-scientific Adam ``the principle that made the sin of the first Adam possible''; but notwithstanding this, dogmatics must think every natural scientist, even the unhappy one, not dyothelitically but monothelitically, since it must absolutely think him not dyophysitically but monophysitically. Now if dogmatics, in accordance with a higher scientific necessity, must think the natural scientist monothelitically, then the natural scientist must also be monothelitic. But, we ask, how can a natural scientist with two wills be monothelitic? As soon as a natural scientist makes as if to want to have two wills and to behave dyothelitically, he must immediately take one of two courses: either out of the congregation or out of science. Let one be ruler, the good will! But two good wills cannot possibly be accommodated in one natural scientist; if the will to acknowledge the laws of nature is good, then the will to acknowledge the law of Christ is not good; if the will to acknowledge the law of Christ is good, then the will to study ``accursed'' nature is not good. Once more: either out of the congregation or out of science! Here there is no mercy. A natural scientist with two wills Bishop Martensen cannot tolerate in his monarchy, and therefore not in his congregation either; for congregation and monarchy, and monarchy and congregation, will not two things, but one. Were Bishop Martensen to tolerate a dyothelitic natural scientist in the congregation, then it would have to have been made comprehensible how two heterogeneous wills, one theoretical and one practical, went together into one; but, says Bishop Martensen, ``how the two wills go together into one has not been made comprehensible, and by the nature of the case never will be.''

% --- p. 40 ---
\opage{40}\setcounter{footnote}{0}%
\section*{II.\\ The Conditions of History in the Monarchy.}
\addcontentsline{toc}{section}{II. The Conditions of History in the Monarchy}
\markboth{The Conditions of History in the Monarchy}{The Conditions of History in the Monarchy}

The deep chasm between the object, method, and aim of the natural scientist and those of the dogmatician has at all times been so conspicuous that already at the beginning of the first part we could well foresee how natural science would fare once the theological universal monarchy was established. With history, now, it is another matter, if we may otherwise place confidence in the very appreciative testimony, highly honorable to history, of ``the Christian dogmatics.'' ``For,'' says Bishop Martensen (Dogmat.\ § 12), ``although it is indeed the Creator Spirit who speaks through nature to the created spirit, nature nevertheless speaks with its mute language only in an indirect and veiled manner of the Creator's eternal power and Godhead; the direct, unequivocal revelation can be found only in the world of spirit itself, in the world of the word, of conscience, of freedom, or in \emph{history}. Revelation and history are therefore inseparable; but if there were no other history than \emph{world history}, then God's revelation would still lack its adequate medium.\dots\ The holy voice of God does indeed sound through the voices of world history, God's work is indeed in the events of world history; but in the tumult of world history God's voice and the voices of men are mingled together for our ear, and the holy will of Providence, which we glimpse through human destinies, again conceals itself from our gaze beneath the restless stream of the world's course. If there is truly to be talk of a holy revelation of God, then there must be a history within history, then there must be within world history a \emph{holy history}, in which God reveals himself as God, a history in which \emph{the holy purpose of the world} reveals itself as such, where the word of God so encloses itself in the human word that the latter becomes
% --- p. 41 ---
\opage{41}\setcounter{footnote}{0}%
the pure organ of the former, and where God's work so encloses itself in man's work that the latter becomes perfectly transparent to the former.''

Now this is both apt and beautiful; appealing especially so long as one views the historical facts from a speculative distance of a couple of miles.\footnote{In my licentiate dissertation \textit{De speculativa historiæ sacræ tractandæ methodo}, translated by B. C. Bøggild 1842, one will find, e.g.\ in § 10 and § 11, similar fundamental thoughts expressed.} But when the facts, in their particulars and multiplicity, are moved closer to the observer; when the relation between historical faith (\textit{fides historica}) and religious faith (\textit{fides religiosa}) is to be precisely determined; when the right of criticism is to be recognized in its essence and extent: then come the difficulties, and the difficulties are many. Without scientific criticism history is no science; but without a critical standard there is no scientific criticism. Now what is it that secures a standard for criticism? The logical-physical! As a science, then, history is on this point just as autonomous as natural science. History, when seen from this side, can no more give up the logical-physical than natural science can; it can no more than natural science let itself be governed by ``the good will.''

When the monarchy demands of the historian that in the last instance he let everything depend on ``the good will,'' the meaning is by no means that he should outright break with the logical-physical, deal arbitrarily with facts, and let the end, in the Jesuitical sense, sanctify the means. No, the difficulty is of a subtler kind. For the good will cannot be separated from the idea of the good; but the idea of the good must, according to the royal principle, be absolutely one with the idea of the true: the good is true, the true is good.
% --- p. 42 ---
\opage{42}\setcounter{footnote}{0}%
For, says Bishop Martensen: ``Only when one has ascended from all subordinate ideas to the good, as to the royal principle \dots\ can there be talk of the living opposition between truth and lie, which is something quite other than the merely logical opposition between correct and incorrect, essence and appearance.'' Here is the crux. It is the relation between idea and actuality, idea and fact, on which the emphasis falls. If the idea is determined by the fact, then actuality rules; but in actuality one does not get beyond the opposition between correct and incorrect, essence and appearance. If the fact is determined by the idea, then the aim is directly at the opposition, the living opposition between truth and lie; but how then does it fare with the merely logical opposition between correct and incorrect, essence and appearance? The more weight is laid here on actuality, the more the idea is pushed back; the more weight is laid here on the idea, the more actuality is pushed back. But from this it follows in turn that the critical standard must change according as it is applied to \emph{world history}, to \emph{church history}, or to \emph{holy history}. In world history actuality with its facts has decisive preponderance over the idea: first the facts, afterwards we can speak of the idea. In holy history the idea, as the idea of the holy, of the absolute good, has unconditional preponderance over the external finite facts: first the idea, then we can speak of the historically actual. In church history criticism wavers between the priority of the idea and that of the fact.

What difficulties may pile up for historical criticism when the principle of the monarchy is to be satisfied is now easy to see. We begin with an example taken from church history, namely the criticism of the Pseudo-Isidorian Decretals. Had it here been solely a matter of evaluating the historical facts, criticism could in this case have made them
% --- p. 43 ---
\opage{43}\setcounter{footnote}{0}%
speak so clearly that not even a shadow of doubt would have remained. In the Isidorian collection there are, after all, letters in which a Bible translation that came into being only after the fourth century is cited from the third century, letters in which the Roman bishop Victor of the second century writes concerning the Easter feast to the Alexandrian bishop Theophilus in the fourth century; here, then, a Jesuitical elevation above the logical-physical is truly required when a historian is nevertheless to attempt to defend the genuineness of the collection. Catholic scholars had already repeatedly (a Petrus Comestor in the twelfth, a Marsilius in the fourteenth, a Nicolaus of Cusa in the fifteenth century) spoken out against its genuineness, when the Magdeburg Centuries at last furnished a detailed proof of its spuriousness, confirmed from several sides. Nonetheless the Jesuit Franz Turrianus ventured to enter the lists for Pseudo-Isidore. The Jesuit and his client were, at the beginning of the seventeenth century, sent packing with such
% Errata: p. 43, line 17 from bottom, "saadan" should read "saadant" (not applied in the transcription).
force by Blondell (\textit{Dav. Blondelli Pseudo-Isidorus et Turrianus vapulantes. Genev. 1628}) that even believing papists thought that now it might be enough. But no, still at the end of the eighteenth century a Marchetti (\textit{Saggio critico sopra la storia di C. Fleury}) could not convince himself of the spuriousness of the Decretals, and there will surely still be found many who cannot convince themselves of it. This obstinacy is more than simple defiance; it is an expression of the fact that the idea and ``the good will'' are in play.

Being itself a creation of the idea of the papal church, Pseudo-Isidore is a bearer of the idea. If now the idea of the pope's primacy is posited as identical with the highest good for the church, and this highest good as the ``royal principle'' of history, then the Jesuits' perseverance in defending Pseudo-Isidore is explicable. That the facts run counter to their endeavors may indeed be unpleasant enough for them; but they do not therefore lose heart.
% --- p. 44 ---
\opage{44}\setcounter{footnote}{0}%
Through great difficulties they help themselves by distinctions. Everyone who would criticize must be able to distinguish, and Jesuits can distinguish. But among all the distinctions belonging to historical criticism one will hardly be able to find any more fruitful than the one by which Bishop Martensen, certainly without any ulterior thought of possible abuses, has separated ``the merely logical opposition between correct and incorrect from the living opposition between truth and lie.'' Once the logical opposition between correct and incorrect is reduced to something subordinate, historical criticism is as if struck by a stroke. What does a historian moved by ``the good will'' care about a few contradictions between individual facts, when it is a matter of defending the idea of the good? Contradictions of this kind can in any case lead only to a difference between the correct and the incorrect; but to defend the idea is to defend the truth. When we see how Bishop Martensen (Leilighedsskr.\ p.~80) can be zealous to have Christ acknowledged as king over all the sciences, and sees the idea of the good satisfied only in this theoretical kingship, how then can it surprise us to see the Curialists zealous to defend a primacy which, according to their ``good will,'' necessarily belongs to Christ's vicar?

Already in the apprehension and determination of an idea in church history, it becomes evident what power religious faith is able to exercise over historical faith; but this power is still far more decisive when the talk is of \emph{holy history}. In holy history the fundamental idea must originally be determined by religious faith, solely and exclusively by religious faith. Therefore dogmatics, as we have seen, is so assured that holy history must absolutely be a history ``in which the word of God so encloses itself in the human word that the latter becomes the pure organ of the former, and where God's work so encloses itself in man's work that the latter becomes perfectly
% --- p. 45 ---
\opage{45}\setcounter{footnote}{0}%
transparent to the former.'' In accordance with this, Bishop Martensen has in his ``dogmatic elucidations'' (p.~55) expressed himself clearly on faith as the first condition of cognition and on the dogmatician's relation to the holy facts. ``I make my own,'' he says, ``the words with which Anselm introduces his Proslogion, and which in my thoughts could serve as the inscription for every dogmatics: `I do not venture, O Lord, to fathom thy depths; for in no wise is my cognition capable of this; I only desire in some measure to know thy truth, which my heart believes and loves. For I do not seek to know in order to believe, but I believe in order to know. For this too I believe, that unless I believed, I should not know either.' --- Since I find no occasion to depart from the track here indicated and followed in my Dogmatik, I continue still to regard the dogmatician's relation to revelation in a manner corresponding to the natural scientist's relation to nature. The genuine natural scientist does not give up nature because he does not find himself able to explain every one of its phenomena, but neither does he give up penetrating into the interior of nature, so far as he is able.''

It deserves the very greatest attention that Bishop Martensen still continues to regard the dogmatician's relation to revelation as the natural scientist's relation to nature; for from this one sees how sure he is of the holy facts. By what is the natural scientist's relation to nature determined? By unalterable laws, by facts that cannot be dismissed. ``The genuine natural scientist does not give up nature because he does not find himself able to explain every one of its phenomena.'' Right! But why does the natural scientist not give up nature? Because he cannot give it up, because in immediate intuitability, in sensible presence, it compels his acknowledgment. Now if revelation, with its facts and laws, is in a corresponding manner to compel
% --- p. 46 ---
\opage{46}\setcounter{footnote}{0}%
the dogmatician's acknowledgment, then it must of course be presupposed that the facts of revelation in their sphere are just as certain, just as intuitable and undeniable as those of nature, the laws of revelation in their sphere just as unalterably determined as the laws of nature. From this presupposition, in itself respectable, the old orthodox dogmatics quite rightly set out; but how has it fared with the presupposition? Historical-philological criticism has treated it harshly, yes, very harshly.

If, as Protestant theological orthodoxy originally presupposes, Holy Scripture is to count as a perfectly reliable, divine document of revelation: then it must be assumed to be raised above all scientific criticism. It was then impossible to conceive that there could be found in Holy Scripture errors either dogmatic, moral, or historical, chronological, topographical, onomastic.\footnote{\textit{Omnia et singula sunt verissima, quæcunque in illa (script.\ sacr.) traduntur: sive dogmatica illa sint sive moralia, sive Historica, Chronologica, Topographica, Onomastica, nullaque ignorantia, incogitantia aut oblivio, nullus memoriæ lapsus Spiritus Sanctus
% Errata: p. 46, note, line 7 from bottom, "Sanctus" should read "Sancti" (not applied in the transcription).
amanuensibus in consignandis sacris litteris tribui potest aut debet. Quenst.\ p.\ I. pg.\ 77.} Θέσις.} Even toward the end of the eighteenth century orthodox theologians asserted not only that the holy writers never erred, but even that they could not have erred.\footnote{``Ihre (der heiligen Schrift) Verfasser haben nicht allein nie geirrt, sondern auch nicht einmal irren können, und ihre Erzählungen sind daher nicht allein so gewiss, als jede andere historische Wahrheit, sondern auch untrüglich gewiss.'' Chr.\ W. F. Walch. Krit.\ Nachricht.\ v.\ d.\ Qvellen der Kirchengeschichte. Leipzig 1770. S.~70.} On so rock-firm a ground holy history could withstand the attacks of criticism. Holy portents, great wonders surround Christ's coming into the world, his conception, birth, and childhood; but no one can deny these portents, these
% --- p. 47 ---
\opage{47}\setcounter{footnote}{0}%
wonders, for they are, after all, facts. Miracle upon miracle meets us in Jesus' public life; they are, after all, facts. Who would not want to believe in Christ's resurrection, when the fact is so historical that no one can deny it without denying all history? What does it matter that the Ascension conflicts with the laws of gravity, when the fact itself is just as certain as the facts from which the laws of gravity have been abstracted? If now ``the genuine natural scientist'' may not deny a natural fact because he cannot at once explain it, then the genuine biblical scholar may still less deny a supernatural fact because he cannot comprehend it. From this standpoint Bishop Martensen is consistent when he ``continues still to regard the dogmatician's relation to revelation in a manner corresponding to the natural scientist's relation to nature.''

But must we not now, on Bishop Martensen's behalf, bitterly lament that this standpoint vanished some hundred years ago, just as we must all, of course, lament that the speculative paradise has vanished? And how has it vanished? It is again ``the story of the Fall and of \emph{the two absolutely heterogeneous trees}, of which the fruits of the one were `forbidden to theology.'\,''\footnote{That the fruits of criticism really were forbidden to theology can be seen, among other things, from a verse that Georg Pasor puts in the mouth of Sebastian Pfocken, known from the Purist controversies:\\
\hspace*{2em}\textit{Diligo philologum, qui eviscerat abdita verba}\\
\hspace*{2em}\textit{Jure novum Christi sed mage fedus amo}\\
\hspace*{2em}\textit{Ergo novo quisquis de federe tenuia decit,}\\
\hspace*{2em}\textit{En pupillam oculi vulnerat ille mei.}\\
\hspace*{2em}--- --- --- --- --- --- --- ---\\
\hspace*{2em}\textit{Falx critica in quosvis per me grassetur. At unis}\\
\hspace*{2em}\textit{A Bibliis fas est abstinuisse manus.}} Criticism, historical-scientific criticism, is here typologically signified by the tree of
% --- p. 48 ---
\opage{48}\setcounter{footnote}{0}%
knowledge; the serpent in the tree is sound common sense, which has always shown a sinful partiality for the logical-physical. Just as the serpent in paradise, that first serpent, the dogmatic-cosmic principle, which is older than the Devil, inasmuch as the Devil, according to our ``Christian dogmatics,'' is a son of the cosmic principle, seduced Eve with treacherous words that could be construed as the instruction: ``It is not true, what God has said (namely from the standpoint of science)'': so now sound common sense, a natural son of the Devil, came hidden in a motley skin of logical-physical texture, crept critically up into the tree of knowledge, stuck out its head with ``autonomy,'' and began to converse with theology in treacherous words that could be construed as the instruction: It is not true, what the Bible has said (namely from the standpoint of criticism). And the woman --- for in the symbol theology is a woman --- saw that the tree, the tree of criticism, was good to eat of, and that it was pleasant to the eyes and a tree to be desired to make one wise, and she took of its fruit, and she did eat, and she gave also a little to the Faculty, and it did eat. Then the eyes of them both were opened, and they saw --- the naked, they saw in criticism naked actuality; and they fled to the grove of aesthetics and sewed themselves fig leaves, and they cut strips from the back of philosophy and made themselves girdles. And they (theology and Faculty) recognized that they could not, without danger of ``a place of their own outside science,'' continue still to regard the dogmatician's relation to revelation in a manner corresponding to the natural scientist's relation to nature; they recognized that they could not continue still to speak loudly of the supernatural facts of history, since the supernatural facts were probably not historical, and the historical ones not supernatural.

Scriptural theology, so youthfully fresh, trusting, and proud of its Bible in its paradise, how it has, after the Fall through the logical-physical, gradually become weakened, bowed, faded,
% --- p. 49 ---
\opage{49}\setcounter{footnote}{0}%
humbled! If it must be said that the Devil in speculation --- and the deepest speculative is always the highest existent --- is a dogmatic son of the cosmic principle and only in this way can be thought as a genuine, objective, scientific Devil: then it must also be said that the theology that trifles with criticism is an undogmatic daughter of fallen orthodoxy. True enough, the fallen one thinks that she still has science and conscience. Yes, conscience! Theology may well reason about conscience, when conscience sits with half a speculation and sleeps in a paragraph of the dogmatics. Let conscience awaken, let theology for once hear the voice of the congregation!

``\emph{What have you}'' --- the congregation asks theology, sinful by its Fall --- ``\emph{done with my Christmas Gospel}?'' Alas, dear congregation, do not be angry. Criticism has unfortunately torn the Christmas Gospel completely to pieces. But do not lose heart! remember that it is knowledge that is to help faith along; remember what Bishop Martensen says with Anselm: ``I do not venture, O Lord, to fathom thy depths''; but I believe that knowledge is to help faith along. Do not weep! The Gospel that criticism has torn to pieces, apologetics can of course put together again; for synoptic apologetics is surely scientifically within its good right when it carefully gathers up the scraps and shows how they must quite certainly have fitted together.

You look at me so darkly. How can I help it that science is objective? Who, now, would ever have thought that the historical facts in the Christmas Gospel were so wondrously brittle, the natural components just as brittle as the supernatural. I will not speak at all now of the supernatural, of the angels who sang for the shepherds, of the star that went before the wise men --- for as to ``the mythicizing garb'' that comes in ``to fill out and shape'' the historically actual, even apologetics admits that it is at its wits' end --- I will, as far as the natural is
% --- p. 50 ---
\opage{50}\setcounter{footnote}{0}%
concerned, merely let you know in all confidence that it is quite wrong with Quirinius.

``\emph{What have you}'' --- the congregation asks theology, sinful by its Fall --- ``\emph{done with my Easter Gospel}?'' Alas, dear congregation, the narratives of Christ's death and resurrection really present not a few difficulties for science. We will not dwell on the chronological-synoptic; for as regards the chronological-synoptic, the congregation cannot sufficiently thank apologetics, inexhaustible in ingenuity. But who does not feel the weight that must press upon science when he thinks of the events at Jesus' death. ``And, behold, the veil of the temple was rent in twain from the top to the bottom; and the earth did quake, and the rocks rent; and the graves were opened; and many bodies of the saints which slept arose, and came out of the graves after his resurrection, and went into the holy city, and appeared unto many.''

In order to come off honorably with this miracle, synoptic apologetics is itself close to performing a miracle, a genuine apologetic miracle. So that you may nevertheless get a little notion of how theological exegetics is able to reconcile faith with knowledge, I communicate here the following interpretation, taken from our newest textbooks: ``The Lord's resurrection had taken place when the owners of the family tombs inspected the damage that the earthquake had wrought in these rock-tomb chambers. Through the cracks in the rocks and the stone doors, which had been loosened and thrown out by the shaking of the earth, remains of their dead were to be seen, as they rested on the floors or on the stone benches. At the sight of these remains the owners of the family tombs received visionary revelations''; and then people deny that here is a thinking ``whose presentation of the miracle outweighs the physical objections.''

``Where,'' you ask, ``has faith gone?'' Faith is
% --- p. 51 ---
\opage{51}\setcounter{footnote}{0}%
of course objectively in the visionary revelations that the owners of the family tombs received; subjectively in the faith with which the interpreter himself believes in their faith, who received the visionary revelations. ``And knowledge?'' Yes, of knowledge there is surely enough here. Here there is, after all, when we properly inspect the damage done at this resurrection, much more knowledge in the interpreter than there is in the evangelist. ``But now the reconciliation of faith and knowledge?'' Dear congregation! If you cannot see the reconciliation of faith and knowledge in the hermeneutic reconciliation of the religious and scientific moments of the interpretation: then it must be because your cognition of faith is only practical and subjective, not at the same time theoretical-objective. True, no proper explanation of the miracle is given here --- ``theology does not occupy itself with explanations of miracles'' --- but a scientific explanation has nevertheless really been given here. Remember, dear congregation, that with our explanations we walk by faith! Yes, we walk by faith, not by sight; but for that reason we may well walk scientifically, if only we take care that during this our earthly pilgrimage no more comes into our practice than there is in our theory. We can then say with Bishop Martensen: ``I do not venture, O Lord, to fathom thy depths; for in no wise is my cognition capable of this''; but this I do venture to assert, that I can very well arrive at a higher scientific cognition of thy miracles without being able to explain thy miracles. Otherwise, after all, I should not be able in \emph{any measure} to know thy truth, which my heart believes and loves. For the faith of my heart is not merely a subjective and practical faith, my love not merely a subjective and practical love; my faith is at the same time an objective and theoretical faith, my love an objective and theoretical love. With the heart we believe, and with the heart we love: as the heart is, so are faith and love; as faith and love are, so is the heart. Remember, dear congregation, we have not merely received a subjective and
% --- p. 52 ---
\opage{52}\setcounter{footnote}{0}%
practical heart; we have in the spirit at the same time received a theoretical-objective heart. Yes, I know that my heart is theoretical-objective; that is to say: I know it in faith. For this too I believe, that unless I believed it, I should not know it either.

``What, then,'' asks the congregation, ``have you, with your theoretical heart, done with the Gospel history, with the Transfiguration on the mount, with the feeding in the wilderness, with the walking on the sea?'' etc.\ etc. We need not continue these negotiations to get an impression that the theology which is in agreement and yet in disagreement with ``the consciousness of culture,'' which knows much and yet cognizes little, has a bad conscience when, in the midst of its knowledge, it is to render the congregation an account of its faith. Whoever wants real insight into the disproportion between faith and knowledge brought about by theology must not let himself be beguiled by dogmatic-speculative phrases, but must acquaint himself with the progressive theological interpretation of the New Testament, in particular of the four Gospels, and come to know the way that has historically and consistently led to a mythical treatment of the life of Jesus.

To be sure, Bishop Martensen says (Leilighedssk.\ p.~92) that it is only ``Schelling's cast-off clothes with which Strauß and others have adorned themselves, in order in this lion's skin to attack Christianity''; but this assumption belongs, according to the logical opposition between the correct and the incorrect, manifestly to the half-correct, not to say the at bottom incorrect. One need only cast a glance at the Straußian ``Introduction'' to the Life of Jesus to convince oneself of the real state of the matter. ``After Semler (Introd.\ § 8) had already spoken of a kind of Jewish mythology and had outright called the narratives of Samson and Esther myths; after Eichhorn had led the way in the manner mentioned (§ 6), the concept of myth was soon, by Gabler, Schelling,
% --- p. 53 ---
\opage{53}\setcounter{footnote}{0}%
and others, with greater right set up as a concept common to all ancient history, sacred as well as profane, according to Heyne's principle: \textit{a mythis omnis priscorum hominum cum historia, tum philosophia procedit}, and Bauer even ventured to come forward with the Mythology of the Old and New Testament (1802). \dots\ Just as (§ 9) Eichhorn, in the domain of the Old Testament, was driven by the force of the matter, in the case of the story of the Fall, from his earlier natural mode of explanation over to the mythical interpretation, so Usteri in the New Testament domain with regard to the story of the Temptation. In an earlier work he had, following Schleiermacher, taken it as a parable delivered by Jesus but misunderstood by the disciples, but he soon saw the difficulties of this interpretation, and since he had left both the supranaturalist and the natural view of the narratives, in their various shadings, still far further behind him, nothing else remained for him but to take up the mythical standpoint, which he then did in a later treatise, and in a very forceful manner''; etc. The theologians, as one sees from this, arrive at their mythical mode of explanation with a necessity that does not hang on Schelling's ``cast-off clothes.'' But Bishop Martensen, who has cast a covetous glance at Schelling's \emph{new} clothes, involuntarily comes to think of the \emph{old} clothes every time he sees a mythical lion's skin; for this much is certain: out of Schelling's new clothes no lion's skin can be made; but a sheepskin can well be sewn from them that ``sparkles with novelty.''

Where the disproportion between the defense and what one wishes to defend is as colossal as, e.g., in the apologetic interpretation of the synoptic Gospels; where the ``concordat'' between Gospel faith and theology is so unfortunate, so painful for both, that theology contradicts faith in every other line, and faith in turn, wherever it sees its chance, theology;

% --- p. 54 ---
\opage{54}\setcounter{footnote}{0}%
where the ``concordat'' breeds discord to such a degree that the contracting parties, despite the caution with which the synoptic apologist constantly marks out the boundary, incessantly jostle and bicker on the middle way, and that notwithstanding that ``the middle way is broad''; where such a phenomenon, we say, sets the tone and takes on the stamp of the normal: there the foundation itself must conceal an enormous misunderstanding.

For a number of years we have steadily kept an eye not only on the speculatively jubilant, ever triumphant dogmatism, but also on the very active, prudently self-moderating, thought-halving doctrine of the middle way. The result of our investigations is that it is by no means scientific criticism, but Christian apologetics, that undermines Christianity in the name of the cultural consciousness. To confirm this view of ours it will suffice, for our present purpose, to bring out the following main points. Gospel-faith and scientific criticism disagree to such a degree concerning the first presuppositions for a judgment of the sacred history that the disagreement can be removed only by referring faith and knowledge to absolutely heterogeneous principles. The gospel history contains such a union of natural and supernatural facts that the judgment on each individual gospel necessarily becomes dependent on the judgment on the original principle of gospel-faith itself. According to the principle of faith, the miracle, the supernatural, belongs so absolutely to the history of revelation, and has for the believing consciousness so incontestable and unconditional a validity, that faith in revelation is essentially a faith in the supernatural. According to the principle of knowledge, the miracle, the supernatural, is absolutely inadmissible: science knows \textit{a priori} that every supernatural event is mythical. It is, when we look more closely, by no means something about which the evangelists do not agree; on the contrary, it is that on which they all agree that is fundamentally an offense
% --- p. 55 ---
\opage{55}\setcounter{footnote}{0}%
to criticism, namely the supernatural. What creates difficulties in earnest is by no means the circumstance that, e.g., one evangelist has two angels at the grave of the resurrection while the other has only one; no, it is much more the circumstance that here there is a resurrection, that here there is talk of angels, that resurrection and angels are placed under the historically real instead of being placed under mythology.

Now where scientific criticism is allowed to develop freely and consistently, as in Strauß and Renan, there the result is clear at once. Even if not a single discrepancy, not a single historical inaccuracy, could be discovered among all the gospel narratives, knowledge would nonetheless drive out faith; for the miracle, scientifically seen, is a myth. It is otherwise with the business of the concordat in Christian apologetics: here faith and knowledge can lie there fooling each other for as long as need be. Pressed by the ``concordat,'' faith dares not take heart and say it outright, that the supernatural is just as indispensable to the principle of faith as the natural is to the principle of knowledge; hampered by the ``concordat,'' criticism for its part will not speak out either. Since in this way the real and essential ground of the disagreement has been concealed, the misunderstanding is a smoldering ember---and so the wrangling over particulars, over the individual narratives, can never come to an end: now criticism acts as if it would gladly believe the gospel of the angels' song of praise, if only there had not been the historical blunder with Quirinius; then apologetics replies and proves that the inaccuracy can very well be excused, and that, without taking the appearances of the angels literally, one can very well assume ``a miraculous fundamental fact''; now criticism acts as if it would gladly grant the great miracle of ``the rending of the veil of the temple,'' if only it could ``believe that some consideration should have held the apostles back from mentioning so significant
% --- p. 56 ---
\opage{56}\setcounter{footnote}{0}%
an occurrence''; then apologetics in return has the childish satisfaction of being able to defend the disputed occurrence with the remark that the miracle of the veil was, after all, not an occurrence of ``such significance that silence about it could count as evidence against the historical fact''; etc., etc. If this mischief, as much scientific as religious, is to be stopped, then faith and knowledge must agree to reject an apologetics that is equally intolerable to both. If gospel-faith is then to be preserved without criticism surrendering its right, this can happen only on the condition that faith and knowledge are separated as two absolutely heterogeneous principles.

When faith is allowed to develop its cognitions from its own principle, the principle of faith itself becomes critical; there forms, as regards the gospel, a criticism of faith in sharp opposition to the criticism of knowledge. How sharp the opposition is can be seen from the fact that what is essential according to the criticism of knowledge is inessential according to the criticism of faith, and vice versa. For the criticism of knowledge it is essential, e.g., to find out the historical presuppositions and conditions under which the representations of the supernatural, the unhistorical, must be assumed to have formed; to that end the criticism of knowledge strives with untiring diligence to investigate the origin of the narrative, its first elements, its legendary growth, in short, its authenticity. It is otherwise with the criticism of faith. It proceeds in principle from an incomprehensible union of the supernatural with the natural, and can in no way permit any attempt to dismember the gospel by separating out the heterogeneous elements; on the other hand, it looks sharply to whether the narrative before it has all the marks corresponding to the essence of revelation; for it does not accept the fundamental miracle on account of the event narrated, but it appropriates the event narrated insofar as this, in all its individualizing features, makes the fundamental miracle intuitable. Of the criticism of knowledge as of the criticism of faith it holds that each of them, according to its
% --- p. 57 ---
\opage{57}\setcounter{footnote}{0}%
principle, must work to bring about a reconciliation between content and form, between the matter and its clothing; but while knowledge with its criticism consistently comes to rest in the myth, faith with its criticism comes to rest in the word of revelation, the word of life that lives in the congregation, the word of light that in the power of the spirit goes forth from the gospel and illuminates the gospel.

In the feeling that things do not, after all, quite hang together as regards the religious yield of scientific research, ``Christian apologetics'' is indeed also, on every suitable occasion, intent on an aesthetic apprehension and appropriation of the sacred facts. ``For the religious-aesthetic significance of these narratives''---it is said of the birth- and childhood-history of John the Baptist and of Jesus\footnote{\textit{Dr.} H. N. Clausen. Fortolkning af de synoptiske Evangelier. Første Deel. Kjøbenhavn 1850. p. 7.}---``in their whole form, the testimony is given in the power with which Christian art in its most beautiful period of flowering felt itself especially drawn to this series of scenes, which---itself a cycle of images in which the wondrous and mysterious emerges in its clearest objective intuitability, in union with festive loftiness and the richest grace of simplicity---has, through pictorial representations and lyrical adaptations, furnished rich means to awaken and nourish the life of ecclesiastical devotion.''

Now even if, from the standpoint of humane culture, we must find what is quoted here to be both soundly and aptly said, it is nonetheless just as certain that the whole conception evidently suffers from the weakness rampant in the cultural consciousness, of preferring a gentle fusion of the religious and the aesthetic to a separation in principle. If the aesthetic view is already heterogeneous with the
% --- p. 58 ---
\opage{58}\setcounter{footnote}{0}%
scientific-critical one for the reason that, whereas the former permits no final conflict between the content and the individualizing form, the latter constantly aims at bringing such a conflict into consciousness. But if the aesthetic apprehension is relatively heterogeneous with the scientific-critical one, it is absolutely heterogeneous with the religious one. The religious consciousness lets the union of the natural and the supernatural events take place in a higher reality, a reality that faith recognizes in a living way, a reality of which knowledge knows nothing.

Before we go on to treat the demands that must be made on ``the good will'' if it is to succeed in assigning biblical criticism its proper place in the monarchy, we shall indicate our own view in a few strokes. That the relation between faith and knowledge as two absolutely heterogeneous principles must show its influence on the treatment of the sacred history is self-evident. In the first part we apprehended the relation chiefly with regard to the \emph{mystery}; we shall now apprehend it with particular regard to \emph{consciousness}. When faith and knowledge are united in one consciousness, it is necessary to apply criticism in order to prevent mixtures. The criticism is twofold: a criticism that sees to it that the two spheres of cognition are kept apart under all reciprocal determinations, boundary propositions and reflected illuminations, and a criticism that sees to it that within each of the separated spheres a distinction is made between the essential and the inessential, the necessary and the contingent, the valid and the invalid. Under this two-sided application of criticism there then arises a \emph{criticism of faith} that stands in the same relation to the \emph{criticism of knowledge} as faith to knowledge: if faith is absolutely heterogeneous with knowledge, then the criticism of faith is absolutely heterogeneous with the criticism of knowledge.

To illustrate the opposition between the criticism of faith and the criticism of knowledge we cite some examples from the Gospels, directing attention particularly to the determination of the relation between idea and reality, between the natural
% --- p. 59 ---
\opage{59}\setcounter{footnote}{0}%
and the supernatural, between myth and revelation. It will then appear: 1) that the criticism of faith, by reducing the outwardly historical to a moment of the inward (\textit{fides religiosa}), builds on the idea, whereas the criticism of knowledge, by making the outwardly historical (\textit{fides historica}) the main thing, builds on the fact; 2) that the criticism of faith, with the almighty will as its principle and the miracle as its absolute presupposition, judges the interaction of the natural and the supernatural from the standpoint of the religious, whereas the criticism of knowledge, with a tacit denial of the miracle, works at getting the natural components separated from the supernatural ones; 3) that the criticism of faith evaluates the historical character-portrait according to its significance for a \emph{revelation}, whereas the criticism of knowledge evaluates it according to its significance as \emph{myth}.

In the gospel of the birth of Christ, e.g., the opposition described will assert itself in every respect. For the criticism of faith the question of what in the historical sense is called authenticity is an altogether subordinate question. This criticism stands on the standpoint of the congregation; it does not doubt the historical authenticity at all, and on the whole has little or no interest in testing the facts in their particulars; on the other hand, it has a sharp eye for the religious authenticity; it aims at developing and making clear the religious opposition between the canonical and the apocryphal. But what if it now learned how inaccurate the narrative is concerning Quirinius and the census? It never learns it, it can never learn it; for the moment it learns it, it is not criticism of faith but criticism of knowledge. What for the criticism of faith is something indifferent, something subordinate and inessential, is on the contrary something decisive for the criticism of knowledge. The criticism of knowledge cannot and will not close its eyes to manifest discrepancies; it neither can nor will make any exception in this respect for the sacred history. So much for authenticity.

% --- p. 60 ---
\opage{60}\setcounter{footnote}{0}%
As regards, next, the relation between the natural and the supernatural elements of the narrative, it will be evident that the criticism of faith cannot possibly aim at separating what the holy will has joined together for the real and effective presentation of the religious idea. For the criticism of faith the angels with their song of praise have just as immediate an existence as the shepherds with their flock. The difference is not that the shepherds are real, the angels unreal, or that the shepherds have a proper, the angels only an improper reality; no, the difference is on the contrary that the angels have a heavenly and blessed reality, the shepherds only an earthly and perishable one. But what if doubt now arose about the relation of the supernatural to the natural, e.g., doubt as to how the star in the sky could literally keep going before the wise men until it stood still over the manger in Bethlehem? Such a doubt can never arise; for the moment such a doubt arises, the moment the aim is to separate out the supernatural as a suspicious admixture, at that moment the criticism of faith has given way to a criticism of knowledge. But, it will perhaps be objected, if no reflections are made on the possibility of what is supernatural in the phenomena, then there is surely no criticism here. We answer: there is a criticism here, insofar as a sharp separation must be made here, with regard to the religious itself, between the essential and the inessential. Indeed, in an otherwise canonical narrative there can be apocryphal components that obscure the religious; the criticism of faith does not go by an indeterminate feeling, it goes by a principle.

As regards, finally, the apprehension of a gospel narrative as a historical character-portrait, the opposition between the criticism of faith and the criticism of knowledge will surely be conspicuous. For the criticism of faith it is a task of decisive significance to make a divorce between the ideas of the imagination and the absolute idea of will; only by reflecting this separation through can the believer make the principle of faith clear to himself and, in consequence of
% --- p. 61 ---
\opage{61}\setcounter{footnote}{0}%
this principle, determine the opposition between revelation and myth. For the criticism of knowledge there is no revelation; what we call revelation is only a particular form of the mythical: that Christianity is the most perfect religion means that Christianity is the most perfect mythology. For the criticism of knowledge, with the myth as its goal, the perfection of the historical character-portrait must consist in its making the idea of the imagination intuitable, in its satisfying aesthetically. ``For the religious-aesthetic significance of these narratives''---it is said of % Errata: "hedder det om" (p. 61, line 10 from top), read: "hedder det, som ovenfor anført, om" (i.e. "it is said, as quoted above, of"). Printed text kept.
the birth- and childhood-history of John the Baptist and of Jesus\footnote{\textit{Dr.} H. N. Clausen. Fortolkning af de synoptiske Evangelier. Første Deel. Kjøbenhavn 1850. p. 3.% Errata: "S. 3" (p. 61, note, line 1 from bottom), read: "S. 7." Printed "3" kept.
}---``in their whole form, the testimony is given in the power with which Christian art in its most beautiful period of flowering felt itself especially drawn to this series of scenes, which---itself a cycle of images in which the wondrous and mysterious emerges in its clearest objective intuition, in union with festive loftiness and the rich grace of simplicity---has, through pictorial representations and lyrical adaptations, furnished rich means to awaken and nourish the life of ecclesiastical devotion.'' This is certainly both aptly and beautifully said---from a religious-aesthetic standpoint, that is. But, we ask, is the principle of the ecclesiastical life of devotion an aesthetic principle? By merely religious-aesthetic merits revelation will never be distinguishable from myth; after all, aesthetic-religious merits are something a Strauß too will grant our canonical synoptic Gospels. No, the opposition, the qualitative opposition between the criticism of knowledge and the criticism of faith shows itself precisely in this, that the criticism of faith posits a difference of sphere between aesthetic contemplation and ``ecclesiastical devotion.'' That devotion can have the aesthetic as a moment is not to be denied; but in principle the aesthetic
% --- p. 62 ---
\opage{62}\setcounter{footnote}{0}%
(the beauty of images, poetry) is something contingent for devotion itself, the religious-ethical the one thing needful, the absolutely essential.

The criticism of faith, as criticism, is not edifying; it moves in reflections and conceptual analyses just as much as the criticism of knowledge does; but since the whole movement of thought takes place in the sphere of faith, since the concepts that are analyzed are concepts of faith, it will be one of the great tasks to develop, and by the interpretation of the individual gospel narratives to make clear, that relation between the historical and the suprahistorical, representing and thinking, total view and impression of the individual, which essentially satisfies devotion, while the concept of devotion is developed at the same time.

If we ask, in conclusion, whether in a consciousness that is to know and recognize both the principle of faith and that of knowledge there is not then a particular psychological medium in which the ``two absolutely heterogeneous principles'' can meet and part and in a manner enter into negotiation with each other, we answer: Yes! there is quite rightly such a medium here. But what shall we call this medium? Knowledge we surely cannot call it; for knowledge belongs to science and is not neutral territory. Nor can we call it faith, for the cognition of faith excludes science, and so does not offer anything neutral either. We must then call it \emph{opinion}. Opinion is the right name, the categorical designation for the medium sought. The mighty opposition between the ancient and the modern view of heaven and earth, an opposition that concerns not merely the phenomena in time and space, but the causes, conditions and essentiality of the phenomena, an opposition that from the standpoint of the cultural consciousness must seem to undermine the divine authority of the Old and the New Testament and therewith of revelation, will not disappear, but will allow itself to be evened out by continued negotiations on the neutral ground of opinion.

% --- p. 63 ---
\opage{63}\setcounter{footnote}{0}%
The essential difference between \emph{knowledge} and \emph{opinion} became clear already to the Greeks; but the difference in principle between faith and opinion seems not yet to have become clear to the theologians. Just as the difference between knowledge and opining can be cleared up only by a criticism of knowledge, so the difference between faith and opinion can be cleared up only in a criticism of faith. An example. What is told in the first book of Moses about Paradise, about the tree of life, the tree of knowledge, etc., is certainly told in the opinion that it all belongs to outward reality, that the tree of life and the tree of knowledge had their roots in the garden's sensible, material soil just as surely as either the vine or the fig tree. In this opinion our Lutheran-orthodox fathers, as recently as a hundred years ago, certainly understood the narrative of Paradise. But in this opinion we, at the present stage of the cultural consciousness, cannot possibly understand it. Now if there were no essential difference between faith and opinion, then a reconciliation of religion and culture would certainly be impossible; but if there is a difference here, a difference in principle, then the reconciliation is possible.

But that there really is a difference in principle here between faith and opinion is to be recognized by this, that we can preserve the faith of the fathers and reject their opinion. It belongs to the indispensable presupposition of faith that the narrative of Paradise, with all the characteristic features contained in it, is ascribed essential and true reality, essential and true actuality. But what reality? An outer or an inner? Well, about that there are different opinions. But then it is surely only a myth; or rather: whether it is myth or revelation that we have before us in this legend---about that there may surely be differing opinion? No, the legend of Paradise, as surely as it reveals to the eye of faith in pure, clear features one of the deepest mysteries of the will, cannot possibly be myth. But how then is one to represent to oneself the relation between the inner
% --- p. 64 ---
\opage{64}\setcounter{footnote}{0}%
and the outer, between what happens bodily and what takes place in the world of the spirit? Well, about that there are different opinions.

That, when real analyses are undertaken, it must come, in the domain of opinions, to manifold conflicts and negotiations between faith and knowledge will, it is hoped, be clear from what has been said here. We have for a long time already been working out these presuppositions, and in applying them to the interpretation of the facts in the sacred history we have glimpsed a prospect that the old disputes about the credibility of Holy Scripture might, not indeed end, but yet enter a new phase. From this standpoint it was then so far from being the case that anyone in the name of the congregation needed to take offense at either the natural or the mythical interpretation, that he ought much rather to appreciate the help and service that science, through its thoroughgoing criticism, has indirectly rendered faith. Had keen-witted men such as Paulus, Eichhorn and Strauß not, by an open, consistent and resolute procedure, displayed the stages characteristic of the way of knowledge, gospel-faith would still have found it difficult enough to find its way back to its original point of departure and to assure itself of its distinctive principle.

These biblical studies, these preparations, alas, these beautiful hopes must vanish like vapor, now that Bishop Martensen goes and sets up his theological universal monarchy. For that matter, it will cost a good deal of work before such a monarchy gets set up. If Bishop Martensen really feels the call and gift of grace of the spirit---and he must surely feel it, since he comes forward with so great an undertaking---to accomplish what ``Scripture clearly says''---with express regard to science?---``that all things, both the visible and the invisible, shall''---scientifically?---``be gathered together under Christ as the head (ἀνακεφαλαιώσασθαι τὰ πάντα ἐν χριστᾷ Eph. 1), that all things are system''---
% --- p. 65 ---
\opage{65}\setcounter{footnote}{0}%
scientific system!---``in him (πάντα ἐν αὐτῷ συνέστηκε Coll. 1)'' (Leilighedsskr. p. 80): then the Bishop must in truth be prepared for a vigorous campaign against the many idols of knowledge: another Elijah, a ``philotheoretical'' Tishbite, he must from his speculative Carmel let fire rain down upon ``autonomy,'' the ungodly autonomy that beguiles philosophy, tempts theology, whores with ``heteronomy,'' and sacrifices ``theonomy'' to the Moloch of criticism.

Here there must be a clearing-up on behalf of history, and a thorough one, not only outside but also inside the camp of theology. For it is not enough that ``antichrists'' like Strauß and Renan are put under the ban; no, theology's own land must for the time being be laid under interdict. For here are theologians, great theologians, men such as Semler, Eichhorn, Paulus of Heidelberg, Schleiermacher, de Wette, Daub, Marheineke and others, of whom it can be proved that they have all had a hand in tasting the fruit of knowledge, have all shown a decided predilection for the logical-physical, have all bowed the knee to rational necessity and sacrificed to the Moloch of criticism. The theologians still living may perhaps be converted; but the dead, after all, cannot. By an antinomy in the Dogmatik, express objection has been made against the conversion of the dead. Surely, for the solution of the antinomy, a scientific exception would never be made for the theologians, so that theologians who while they lived had neglected to change standpoint in due time were allowed to change when they were dead? If this will not do, it must fall especially hard on the dead. Here, through the excommunication of so many and such significant theologians, there will be a gap which, since the excommunicated by their reputation and their number represent a whole age, we may perhaps call the gap of the apostasy, of the great apostasy, in the history of theology. It will, of this we are convinced, grieve Bishop Martensen for many of these men of honor, not for Eichhorn, not for Paulus, but for Daub, for Marheineke
% --- p. 66 ---
\opage{66}\setcounter{footnote}{0}%
and above all for Schleiermacher. Theology would gladly have preserved Schleiermacher as a witness to the truth, a genuine witness to the truth; but, by the royal principle! it is impossible. A theologian whose criticism does not even allow the Gospels to be genuine---how should he himself be genuine?

What in all the world can it help to scold the knowledge that sublimates naturalism and rationalism into speculative theology, since the naturalism is theology's own naturalism, the rationalism theology's own rationalism, since it still sounds down to us from the theological summit that the highest being is at the same time the deepest speculative? How, if we may ask, have the gentlemen rationalists all at once got beyond their rationalism? Is it by saying ``historical'' and ``something historical'' and ``a certain fundamental fact,'' or by declaiming about ``a history that precisely as historical reality has become the turning point in the religious history of the world''? Or has perhaps a new amalgam of orthodoxy and rationalism been discovered? What can it help to want to conceal what can after all be seen from all sides, that Strauß's ``Jesu Liv'' is a theological mirror in which the spirit of the age can contemplate its own form? ``If there is,'' says Baur,\footnote{Cf. \textit{Dr.} H. N. Clausen: Fortolkning af de synoptiske Evangel. 1ste Deel p. 74.} ``any work that, as far as possible, is built solely on the labors of its predecessors, seeks solely to put together the results of a series of investigations conducted long ago by many others, seeks solely to draw the last consequence from premises on which agreement had been reached beforehand, it is this one. When one blames the Straussian criticism for the negativity of all its results, then let one say what other result than a merely negative one could be drawn from the investigations then existing into the origin of the Gospels and their mutual relation!

% --- p. 67 ---
\opage{67}\setcounter{footnote}{0}%
As opinion stood against opinion, and all the opinions together contradicted and canceled one another, they all fell into an indifference in which one had nothing fixed and certain to hold on to.'' If Strauß is to be excommunicated, then there are many who must be excommunicated; if Strauß is ``an Antichrist,'' then there are many ``Antichrists'' in the camp. If the kingdom is to be established, then judgments of punishment must go before it; then Elijah must come with flames of fire, in the storm of wrath, with a glowing chariot --- dogmatically glowing horses.

To stay with the types, it is not even enough for the cleansing of the sanctuary that Bishop Martensen sits dogmatically in a glowing chariot and is Elijah; if the whole undertaking is not to end in ``a literary scandal (\textit{parturiunt montes}),'' he must necessarily also be a Judas Maccabaeus: only a Judas Maccabaeus can cleanse the sanctuary and celebrate the festival of the dedication of the temple which is to be celebrated here. But in order to cleanse the sanctuary and celebrate the great festival there is one point, admittedly only one single point, but an important point, that he absolutely must clear up. \emph{How is the logical-physical related to the higher scientific?} To this question a clean and clear answer must be given. Without a clear answer to this question we really know neither up nor down; and when, in the midst of all our dogmatic knowledge, we know neither up nor down, what good can it do that we exert ourselves to be either Elijah or Maccabaeus? Once more: \emph{How is the logical-physical related to the higher scientific?} where is the boundary?

Bishop Martensen's pronouncements concerning this important point are no doubt, each taken by itself, very definite; but since the one pronouncement goes so definitely in one direction and the other in another, the ``total view'' in the end becomes indefinite out of sheer definiteness.

Bishop Martensen must --- he says this very definitely ---
% --- p. 68 ---
\opage{68}\setcounter{footnote}{0}%
continue ``henceforth to regard the dogmatician's relation to revelation in a manner corresponding to that of the natural scientist to nature.'' But Bishop Martensen must --- he says this just as definitely (Leilighedsskr. p. 93) --- acknowledge ``that with the ethical concept of God there enters a concept of knowledge different from that which obtains where there is question only of a merely logical and physical knowledge. For if God is the personal, willing God, then we can know neither more nor less of him than he himself will reveal to us.'' Now, in order most definitely to strengthen this latter proposition into irreconcilable conflict with the former, it is added: ``Construct him (the ethical God) a priori we cannot, for only the logically necessary, but not the free spirit, can be constructed a priori.\dots\ \emph{Compel} the ethical God to reveal himself to us, as we can compel nature by physical experiments to betray her secrets, we also cannot''; in other words: ``regard the dogmatician's relation to revelation in a manner corresponding to that of the natural scientist to nature'' we also cannot. ``Quite freely must the God of the good communicate himself to us, and we should not even be able to seek him if this seeking did not stem from himself.'' Thus: just as the natural scientist's relation to nature rests on the logical-physical, insofar as it rests on a priori-empirical necessity of reason, so the dogmatician's relation to revelation does \emph{not} rest on the logical-physical, insofar as it does not rest on a priori-empirical necessity; to be sure there is a ``just as'' here, but a ``just as'' of this nature is strongly reminiscent of the well-known one: just as the lion is a ravenous beast, so shall we also walk in newness of life.

Phrases and total views are of no use. If we are to enter the universal monarchy and by \textit{credo, ut intelligam} make ``the good will'' master over science, then the good will must really have the goodness to issue its orders in definite form, to give us clear instructions, and not to fob us off with total views.

% --- p. 69 ---
\opage{69}\setcounter{footnote}{0}%
I remark this, however, not for my own sake, for I do not intend to enter the universal monarchy, but for the sake of those who either are already inside it or might wish and desire to get in; in particular I think again with concern of the Faculty's synoptic apologetics, for the synoptic apologetics is really --- in the middle of the middle way --- in embarrassment over the boundary of the logical-physical.

It is not learning, not acuteness, not critical tact and refinement that the synoptic apologetics lacks; it lacks only one thing: a principle. In an unprincipled manner it has, by the ``\emph{concordat}'' between faith and knowledge (to repeat one of the Bishop's favorite expressions), entered into negotiations between the physical and the hyperphysical which can do nothing but halve its consciousness and inflict manifold logical inconveniences on the halved consciousness.

Let one cast a glance, for example, at one of the Gospels, that of the feeding of the five thousand in the wilderness. For exegetical consideration --- be it said for the comfort of faith --- this miracle has ``as complete an external warrant as any of the Gospel miracles.'' But when the exegete now --- in order to satisfy the demand of knowledge --- is, according to Bishop Martensen's dogmatic direction, to relate himself to this sacred fact ``in a manner corresponding to that of the natural scientist to nature,'' what does he then see? He sees ``an act of power, in absolute incomprehensibility, of the creative power of the Deity, a τέρας which closes itself off against every exertion of thought'': now consciousness is halved. The halving of consciousness is here so striking that one can show with almost anatomical precision how the one half holds on to faith, the other to knowledge.

The interpreter ought --- so runs the demand from the believing half --- to raise himself above the logical-physical, ought in earnest to hold fast to ``the ethical concept of God,'' and to comfort himself that, since God is ``the personal, willing God, we can know neither
% --- p. 70 ---
\opage{70}\setcounter{footnote}{0}%
more nor less of him than he himself will reveal to us''; now ``the personal, willing God in the Gospel before us has willed to reveal to us merely that, but not \emph{how}, he can feed five thousand with five loaves and two fishes.'' We are, in other words, to keep to the ethical without asking about the physical. Herewith the Bishop's one direction has been complied with.

But --- remarkably enough, how everything fits in theology --- just as the halved consciousness needs a halved direction, one half for faith and one half for knowledge, so Bishop Martensen has also given apologetics a whole direction falling into two halves: one half which, in accordance with ``the ethical concept of God,'' suits faith, and one half which, on the analogy of the natural scientist's relation to nature, suits knowledge. From the latter half, then, comes the demand that the interpreter, as a theological natural scientist, strive to penetrate into the interior of the event. ``The genuine natural scientist does not give up nature because he does not find himself able to explain every one of its phenomena, but neither does he give up penetrating into the interior of nature so far as he is able.'' The genuine theologian does not give up the miracle because he does not find himself able to explain it --- who, indeed, busies himself with explanations of miracles? But he does seek, so far as he is able, to penetrate into the interior of the miracle. The genuine theologian cannot refrain from asking about the physical: if he overlooked the physical in the hyperphysical and contented himself with the ethical, he would be no genuine theologian at all; as a genuine theologian he must then ask ``whether the increase of the provisions is to be thought of as having taken place in the hands of Jesus or in those of the disciples, and how it took place: whether by growth (but the talk here is not of natural bodies which still stand under an organic law of development) or by multiplication (Orig. \textit{in Matth.} 1, 11: τῷ λόγῳ καὶ εὐλογίᾳ αὔξων καὶ πληθύνων. Clem. Strom.
% --- p. 71 ---
\opage{71}\setcounter{footnote}{0}%
VI, 11): in the case of the loaves a growing, in the case of the fishes a resurrection, etc.''\footnote{\textit{Dr.} N. H. Clausen, Fortolkn. af synopt. Evang., 1st part, p. 347, note.}

Now when the genuine theologian, by asking scientifically, has assured himself that no single reasonable answer can be given here to his many reasonable questions, indeed, has assured himself that an intrinsically unreasonable miracle, ``a τέρας which closes itself off against every exertion of thought, rejects every question as unanswerable, stands before us'': what does he do then? Does he perhaps give up such a miracle, or demand that it be counted among the apocryphal elements that have crept into the New Testament? By no means! That would be flatly contrary to the ``concordat,'' an indefensible encroachment by knowledge. Only in a Strauß and those like-minded with him do we find such encroachments, such glaring blunders! Only a Strauß will say it straight out, that ``this miracle belongs to those which can seem in any way credible only so long as one knows how to keep it in the half-darkness of an indefinite notion: as soon as one draws it forth into the light and will examine all its parts closely, it dissolves into a figure of mist. Bread which swells up in the hands of those distributing it like sponges that have been moistened; fried fish whose broken-off parts grow back, like the claw one tears off a living crayfish, obviously do not belong in the realm of the actual, but in a quite different one.''\footnote{Jesu Liv § 102.}

But for that very reason Strauß, as one sees from this extremity, is no genuine theologian; he lacks something: he has a knowledge, but he lacks faith. To lack faith is, in the theological sense, to lack the one, and no doubt the better, half of consciousness; to lack faith is to misjudge the ``concordat.''

% --- p. 72 ---
\opage{72}\setcounter{footnote}{0}%
It would be a breach, a sinful breach of the ``concordat,'' if any theologian gave up faith (\textit{fides historica}) in the actuality of the miraculous feeding, ``since for exegetical consideration this miracle has as complete an external warrant as any of the Gospel miracles.'' We must constantly remember that it is \emph{historical faith} which, according to the ``concordat,'' furnishes the right, exegetically reliable foundation for religious faith. Were the genuine theologian. to give up faith in the historical element in the miracle of the feeding, he would indeed have to give up faith in the historical element in the whole of sacred history. But, we ask, what then in this story of the feeding is the properly historical? To this the genuine theologian must give us the definite answer: \emph{nobody knows}!

If only the genuine theologians could stop at a halving of consciousness, that could pass; they could then believe in the miracle of the feeding with the one half and know, with the other, that the miracle is incredible. But that will by no means do; Bishop Martensen will not permit it. Every reasonable person must agree with the Bishop that a ``dualism between the theoretical and the practical'' such as the one described here cannot be resolved in ``a self-contradictory concordat.'' But how will Bishop Martensen now in actual fact --- for general phrases accomplish nothing --- go about helping apologetics, stuck deep in the ``concordat'' and ceaselessly wobbling between faith and knowledge, out of the halving?

Indeed, if there is in any scientific consciousness a sign of dualistic halving, it must be in the synoptic-apologetic one. Let one imagine a conflict between faith and knowledge such as this: Faith demands that in the miracle of the feeding I acknowledge the historically actual; for if I deny it, I deny the Gospel history. Knowledge shows me that I cannot possibly either think or picture to myself the historically actual in the miracle of the feeding. I am thus to believe in
% --- p. 73 ---
\opage{73}\setcounter{footnote}{0}%
the actuality of that whose actuality I can neither think nor picture to myself. What faith demands is impossible for knowledge; what knowledge demands is impossible for faith: if this cannot be called dualism, what then can be called dualism?

Now since nothing is so dangerous to the universal monarchy as dualism and self-contradictory concordats, Bishop Martensen, who wants to establish the monarchy, must surely know a way to abolish the synoptic-apologetic dualism, the synoptic-apologetically self-contradictory concordat.

But how is the dualism to be abolished? To make a critically subtle separation between the matter in itself and its clothing, to give up the Gospel account as ``a mythicizing'' conception and to appeal to an unknown and unknowable ``basic fact'' as to a Kantian Ding an sich, will by no means do. In the first place it conflicts with the respect we owe to the Evangelists' eye for the historically actual. For if the event narrated stands on the same level of credibility as any other Gospel miracle whatsoever, how will one be able to call the historical character of the account into doubt without calling into doubt the credibility of the whole Gospel history? Doubt about the Evangelists' capacity for historical narration is \textit{eo ipso} doubt about the history narrated by them. Next, the separation between a basic fact unknowable in itself and a phenomenon of actuality heterogeneous with this basic fact, a phenomenon of actuality which, into the bargain, is just as unknowable as the corresponding basic fact, is utterly contrary to the essence of history, of sacred history. Let us above all not forget what ``the Christian dogmatics'' has taught us about sacred history, that it is a history ``in which God's word so encloses itself in the human word that the latter becomes the pure organ of the former,'' a history ``in which God's work
% --- p. 74 ---
\opage{74}\setcounter{footnote}{0}%
so encloses itself in man's work that the latter becomes perfectly transparent to the former''; in other words: a history in which every basic fact can be known in its phenomenon, since the phenomenon becomes transparent, ``perfectly transparent,'' to its basic fact.

But how is the dualism to be abolished? To appeal to tact and sense, a peculiar union of religious sense and critical tact, can surely not lead to anything either. Of course, it can no doubt lead to a synoptic-exegetic-apologetic tactic; but this tactic is of a very dubious kind. The tactic consists in making a distinction among the Gospel miracles, not with regard to the credibility of the historical --- for the ``historical,'' the ``basic-historical,'' must be held fast --- but with regard to their, if we may say so, religious fruitfulness. In this respect one divides the miracles, not on principle but by way of suggestion, into two main classes with a suitable intermediate class. To the first main class belong miracles of religious fruitfulness, such as are properly accessible to dogmatizing, ethicizing thought; we can call them the \emph{reasonable, almost rational miracles}. To the second main class belong the miracles that are unfruitful in the religious sense (the miracle, for example, of the swine of the Gadarenes); we can designate them as the \emph{unreasonable, almost irrational miracles} for sound human understanding.

Can the dualism now perhaps be abolished by introducing an intermediate class and applying gradations? By no means! For it is presupposed that all miracles against whose historical credibility criticism, apologetic criticism, has nothing to object are to be believed without distinction. But faith in ``an unreasonable'' miracle, offensive to the understanding, is a faith in the absurd, a faith which stands in sharp opposition to knowledge, whereas faith in a reasonable miracle, accessible to religious and scientific thought, is
% --- p. 75 ---
\opage{75}\setcounter{footnote}{0}%
a rational faith, reconciled with theological knowledge. By this tactic, then, faith becomes absolutely heterogeneous with faith; the dualism enters into faith as the demons of old entered into the swine.

But how is the dualism to be abolished? The solution of the riddle cannot possibly consist in faith and knowledge, each from its side, ceasing to be principles. Certainly on these terms we should be spared being plagued with two absolutely heterogeneous principles; indeed, we should even be spared being plagued with any principle at all: the solution of the riddle would then be that the whole is dissolved into a complete absence of principle. But, by the royal principle of the monarchy! a complete absence of principle is the complete ruin of theology.

How then is the dualism to be abolished? Bishop Martensen must know; I do not know.

\section*{III.\\ The Position of Philosophy in the Monarchy.}
\addcontentsline{toc}{section}{III. The Position of Philosophy in the Monarchy}
\markboth{The Position of Philosophy in the Monarchy}{The Position of Philosophy in the Monarchy}

An assessment of the position that philosophy will come to occupy in the great theological monarchy can perhaps best be introduced by directing attention to what is, for that matter, only of passing and psychological interest: that Bishop Martensen has announced his monarchical principle and begun the new conquests by writing a whole occasional piece, full of many different and interesting pieces of information about many and different philosophical concordats and theories, against my philosophy. ``Whereas Professor Nielsen,'' says the Bishop (Leilighedssk. pp. 8--9), ``has set himself the great and comprehensive task of lodging a protest against all theology, I have in the present work set myself a far more subordinate task, namely, to lodge a protest, not against all philosophy, but only against the R.
% --- p. 76 ---
\opage{76}\setcounter{footnote}{0}%
Nielsenian.'' We shall in what follows try to show how Bishop Martensen has come to solve his task, no doubt subordinate in and for itself, in so superordinate a manner that he can nevertheless really --- perhaps without his knowledge --- be said to have lodged a protest against all philosophy.

As we are to look up, from a more subordinate, merely logical-physical standpoint, toward what is superordinate in Bishop Martensen's supra-logical, supra-physical philosophy, we shall, for the sake of order, once again, perhaps for the last time, begin by examining how matters stand with philosophy, in particular logic, in our ``Christian dogmatics.'' Since the standpoints are so different, we shall begin with a glance at the difficulties flowing from this.

From the logical-physical point of view of philosophy it is naturally always a difficult matter to enter into scientific negotiations with a theologian; the more distinguished the theologian is, the greater the difficulty. Now if Bishop Martensen belongs --- as we may well presuppose as given --- among the most distinguished theologians of the present day, and if the presupposition we have set up is correct, then it follows that philosophical negotiations with Bishop Martensen must be among the very most difficult. Great difficulties, however, according to the old saying \textit{per ardua ad astra}, have something tempting about them; herein let one seek the reason why I have so often, perhaps all too often, been tempted. Eighteen years ago, when the temptation first visited me, so that in what seemed to me a fairly objective review I ventured to set \textit{Dr.} Martensen's Dogmatik alongside S. Kierkegaard's Johannes Climacus, I was struck by the words of judgment: ``Blessed are they who will not force their wisdom upon others!'' which being interpreted is: Blessed are they who will not criticize ``my Dogmatik''! And I, unblessed one, who now, as punishment for having disobeyed the judgment, have got so much into the swing of criticizing ``my Dogmatik'' that I can scarcely stop again: let this be a warning
% --- p. 77 ---
\opage{77}\setcounter{footnote}{0}%
example to younger fellow philosophers, that they think well before they touch a certain something called ``my Dogmatik''; for to touch ``the Dogmatik'' is in philosophy what it is in the fairy tale to touch ``the golden goose.''

Let us hear Bishop Martensen's declaration, well considered after seventeen years' consideration (Leilighedsskrift p. 7): ``But when I for my part now come forward against Prof. R. Nielsen, let no one by any means expect that I will in any way undertake to defend my Dogmatik against the attacks he has directed at it. When a scientific protest has been lodged against `all' theology (preface to Grundideernes Logik II), when, therefore, the very possibility of theology in every form is denied, it can certainly not occur to me to come forward with a defense of my own attempt at a presentation of the Christian doctrines of faith.''

Is this now not a fearful warning? In my guilelessness I proceed from the presupposition that Bishop Martensen's presentation of the Christian doctrines of faith is just about the best scientific presentation that can be given in our time,\footnote{This has been aptly expressed by one of the younger participants in these negotiations: ``When Professor Nielsen in his attacks on theology has constantly aimed at Martensen's Dogmatik, this has happened for the same reason for which, in a bombardment, one aims at the church towers.'' Rudolf Schmidt in Dansk Kirketidende. 1866. No. 43, p. 679.} and that therefore the attack which essentially strikes this theology must strike all theology; and behold! Bishop Martensen then has the unforeseeable modesty to reduce his presentation to an attempt, perhaps even a subordinate one, and that without so much as giving me a hint as to where I might then find a more significant attempt, a worthier object for my criticism.

% --- p. 78 ---
\opage{78}\setcounter{footnote}{0}%
But why --- let us hear it once more --- can it not occur to the Bishop to come forward with a defense of his own Dogmatik, attacked by me with so many definite objections, while after so many years of silence it has nevertheless finally been able to occur to him to send out an occasional piece with a protest against my philosophy? The reason, he says, is that ``the very possibility of theology in every form is denied'' in my attacks. This argument is typically theological; for it has the remarkable property that when one says \textit{credo, ut intelligam} and sees it in a total view, it grows to the size and weight of a mountain; if, on the other hand, one does not say \textit{credo, ut intelligam}, but subjects it to a logical analysis, it shrinks and vanishes.

Thus, with \textit{credo, ut intelligam}: No one can attack theology with emphasis without attacking it in a certain definite form; theology receives its definite form in the individual theologians, in the individual theological ``attempts.'' The individual ``attempt'' is not theology; theology as theology must be sought in all ``actual and possible attempts'' of all actual and possible theologians. In this way it is an impossibility to attack theology. Every time a theology has been attacked, even if the attack was so striking that the theologians could not say a word against it, one still has the comfort that it is only a theology, not theology, that has been attacked. Even if a philosophizing critic had so superhuman a productivity that he could dispose of a single theology every day over a period of fifty years, he would still be far from finished with all actual theologies; and even if he finished with all the actual ones, he would still have all the possible ones left. If, into the bargain, each individual theologian, following Bishop Martensen's example, can ignore every attack on theology in his own theology so long as the attacker has not
% --- p. 79 ---
\opage{79}\setcounter{footnote}{0}%
finished with all other possible and actual theologies, then it is evident from this that the theologians need trouble themselves as little about attacks on the individual theologies as about attacks on theology. No science can occupy a stronger position --- that is to say, when it holds fast to \textit{credo, ut intelligam}.

But without \textit{credo, ut intelligam}, what then becomes of the total view? The individual theologian must be enormously modest if he is really to consider his theological attempt so wretched that an attack can well strike it without in any way affecting theology itself. Such modesty is unnatural for a theologian, unnatural for a theological bishop, altogether unnatural for a Bishop Martensen. The more significant the individual theological attempt is, the greater the significance that not only an attack on, but also a defense of, such an attempt must be able to have for theology as theology.

By being unwilling to rebut an even fairly detailed and persistent attack on his own theology on the ground that the attack coincides with a protest against all theology, Bishop Martensen has hardly done either himself or theology as a whole any great service.

Had the Bishop --- which must indeed have been an easy matter for him with his ``royal principle'' --- struck a blow in defense of his ``attempt,'' then the victory won would obviously have cast luster not only on his theology but on theology as a whole. Non-theologians, such as have already begun to waver in the faith, would then have said, every man to his neighbor: ``There you can see what a false alarm it is, this constant crying that theology must go; they wanted to attack all theology, and see now how pitifully they have come off by attacking Bishop Martensen's theology; yes, yes, theology --- we thought as much --- is, upon my faith, a science after all.'' Now, on the other hand, since Bishop Martensen withdraws from the defense on
% --- p. 80 ---
\opage{80}\setcounter{footnote}{0}%
the pretext that it is all theology that has been attacked, a mistrust easily awakens in those who are not firm in the faith; they will say: ``What a pity that Bishop Martensen, the only one here at home who can strike the blow, has let so good an opportunity for a brilliant victory slip through his fingers: surely things can't be amiss with the old theology?'' That is naturally unbelief, we admit it, sheer unbelief; but such are people: the world lieth in wickedness.

Not only for the sake of theology but also for his own sake Bishop Martensen ought, it seems to me, to have undertaken to defend his ``attempt at a presentation of the Christian doctrines of faith.'' Suppose that on the part of the opponent it really was a rash act, this lodging of a protest against all theology: would it follow without further ado that it must also be a rash act to lodge a protest against the scientific character of Bishop Martensen's Dogmatik? It takes a strong \textit{credo} for anyone to hold fast to a conclusion such as this: to criticize Bishop Martensen's Dogmatik is to lodge a protest against all theology; to lodge a protest against all theology is unblessed: blessed are they who will not lodge a protest against all theology; blessed are they who will not criticize Bishop Martensen's Dogmatik. Such a \textit{credo} would surely be uncanny for the dogmatician himself and would conflict with his modesty.

But if such a \textit{credo} may not be laid at the foundation of our consideration of the matter, how then can it occur to the Bishop to use the circumstance that in the preface to Grundideernes Logik II a protest has been lodged against all theology as a pretext for withdrawing from defending himself against the protest that has been lodged against his own theology? Something must lie behind this.

I have it. The reason why Bishop Martensen will not defend himself against my attack is, as I

% --- p. 81 ---
\opage{81}\setcounter{footnote}{0}%
well have supposed, not at all the reason that I have lodged a protest against all theology; no, the reason is quite a different one. In the occasional piece (p. 106) the reason is stated; the reason is that one fundamentally cannot have dealings with me. ``It belongs to the basic character of Prof. Nielsen's polemic that he constantly attributes to his opponents intentions which they do not acknowledge, and which he is not entitled to impose upon them; and it is this character of his polemic that makes any discussion with him impossible.''

Here, then, we have a proper reason, and, if the matter stands as Bishop Martensen says, a sufficient reason, a reason that may well explain to the reading public why the Bishop does not wish to attempt any defense. But then he could have told us so straightaway. In a scientific discussion of faith and knowledge it does look quite peculiar, almost ladylike, first to fob us off with a reason that is not a reason, only afterward, when we least expect it, to surprise us with the right and proper reason. So then --- \textit{hic Rhodus hic salta!} --- the character of my polemic makes any discussion with me impossible; that matter must be investigated.

Fortunately the investigation can take place quite objectively, since under the present circumstances it will not be difficult to lay out so many data that the thinking reader can himself follow along and see where we stand.

It is quite correct, what Bishop Martensen states, that in my reply to \textit{Dr.} Zeuthen I called theology a science of miracles, a science that explains miracles. What, now, has the Bishop to object to this? His rebuttal is most peculiar: ``Theology,'' he says, ``is a science of revelation, and to revelation belongs the divine wonder. But theology is by no means a contingent collection of miracles. Nor does it occupy itself with explanations of miracles.'' The notion that seems especially to arouse
% --- p. 82 ---
\opage{82}\setcounter{footnote}{0}%
the Bishop's displeasure is the notion that explanations of miracles would absolutely require that theology become a science of ``a contingent collection of miracles,'' and that it would then be a matter of ``explaining the divine mode of operation itself, or the How of the wonder.'' But in this he attributes to me ``presuppositions and intentions'' which I do not acknowledge, and which he is not entitled to impose upon me. For I have really nowhere put forward the unreasonable demand that an explanation of miracles should consist in explaining ``the divine mode of operation itself''; still less could it occur to me that theology would become a science of miracles only by becoming a science of ``a contingent collection of miracles.'' I have simply and plainly asserted that theology, in order to be science, must be a science of miracles, a science that explains the miracles, that is, gets something scientific out of the miracles.

What, then, would Bishop Martensen have stand here in place of ``science of miracles''? Science of revelation! ``Theology is a science of revelation, and to revelation belongs the divine wonder.'' How characteristic! When the theologian defines his science, the definition becomes as of itself an ambiguity, a riddle; for a proposition such as this: to revelation belongs the divine wonder, is surely as ambiguous as a dark saying, as enigmatic as an oracle's answer. To revelation belongs the divine wonder: what does that mean? Might the meaning perhaps be that revelation is one thing and the wonder (by the wonder we may surely understand the miracle) another; why, that would be quite clever: in this way theology could become a science of revelation without becoming a science of the miracle, a science of revelation that was not a science of miracles.

But how are we unfortunate people, who would now so gladly understand what Bishop Martensen understands by theology, to manage to conceive of a revelation that is not a miracle, and a miracle, namely a ``divine''
% --- p. 83 ---
\opage{83}\setcounter{footnote}{0}%
miracle, that does indeed belong to revelation but does not coincide with revelation? For, I will honestly confess it: in my untheological simplicity I had hitherto imagined that when there was talk of wonder, of the wonder, particularly and emphatically of ``the divine wonder,'' then the wonder was revelation, and revelation wonder. But that will not do: to revelation belongs the divine wonder, that is to say: to revelation belongs revelation; no, by Hermeneutics, that will not do. There is then nothing for it but to understand it according to the letter: theology is a science of revelation, but no science of miracles, that is: theology cannot get anything scientific out of the miracles.

But when theology has at last resolved to let the miracles, indeed the miracle itself, remain outside science, can it then, perhaps, get revelation into science? No, that cannot be the meaning either. What, then, is the meaning? I cannot find the meaning; I explain and explain, but the longer I explain, the more inexplicable the whole becomes. I may begin from whichever end I like, it always ends in the same way; it ends with Rahbek's refrain: ``Whether you take it lengthwise or across, / It is as foolish as Lüttikau's verse.''

But, thinks the fair-minded reader, one ought not to fasten on a single word, a single proposition, but should nicely explain the particular from the whole and earnestly endeavor to understand Bishop Martensen from his ``total view.'' Very well! we shall try it with the total view. ``I continue'' --- this, after all, is Bishop Martensen's total view (Dogm. Oplysning. p. 55) --- ``to regard the dogmatician's relation to revelation in a manner corresponding to that of the natural scientist to nature.'' Thus: just as little as the natural scientist can separate nature from the facts of nature, just so little can the dogmatician separate
% --- p. 84 ---
\opage{84}\setcounter{footnote}{0}%
revelation from the facts of revelation. But to the facts of revelation belong, surely, the miracles. Thus: just as the natural scientist's explanation of nature coincides with the explanation of its facts, so also the dogmatician's explanation of revelation coincides with his explanation of the facts of revelation, i.e.\ of the miracles. The total view, as one sees, speaks expressly for the claim that theology must be a science of miracles, a science that does, in one way or another, occupy itself with explaining miracles.

But against the total view stands the special view: ``nor does it (theology) occupy itself with explanations of miracles,'' most decidedly.

The only way out that remains here is to investigate, not the various meanings of the words ``explain'' and ``explanation,'' for that leads nowhere, but the treatment of the miracles that our dogmatician has himself employed, in that he has shown us in practice what is to be understood by the enigmatic words: ``Theology is a science of revelation, and to revelation belongs the divine wonder.''

In order not to attribute to my opponents ``presuppositions which they do not acknowledge,'' I referred expressly, in my reply to \textit{Dr.} Zeuthen, to ``the Christian dogmatics.'' My words (p. 6) are these: What, then, has our ``Christian dogmatics'' brought forth toward the understanding of the three fundamental wonders of revelation: \emph{creation}, the \emph{incarnation}, and \emph{inspiration}? In a few strokes and in outlines as sharp as possible I there sought to show what dogmatics has achieved with its explanations. The objection that one cannot take account of a characterization so brief that it covers three fundamental wonders, and thus in effect the whole of dogmatics, in four octavo pages, does not hold in science; for science must aim just as fully at reducing as at deducing.

% --- p. 85 ---
\opage{85}\setcounter{footnote}{0}%
Had Bishop Martensen, as regards the reduction referred to, really had something to complain of, something that could in any way entitle him to so personally offensive a remark as that I attribute to my opponents ``presuppositions and intentions which they do not acknowledge,'' and that it is this character of my polemic that ``makes any discussion'' with me ``impossible'': then he had here the best opportunity to catch me red-handed. I, who had the presumption to wish, by a compressing reduction, to crush ``the Christian dogmatics'' and with it ``all theology,'' would then myself have been crushed. For I shall not deny that if it could really be proved that the reduction contained either distortions or essential misunderstandings, the theologians might well begin, little by little, to put up the illuminations.

Why has Bishop Martensen not taken hold of the reduction? Why has he contented himself with bare accusations, where it would surely have been easy for him to procure a little proof? Instead of venturing out with a foreign logic upon the open sea of ideas, among the maelstroms of crossing reflections, where a little fishing boat can so easily capsize, it would surely have been far safer to keep to the well-known waters, to stay by the dogmatic coastland and see to guarding his own a little. That Bishop Martensen has as yet had occasion to read Grundideernes Logik we may certainly not presuppose, judging from what we can gather from his citations in the occasional piece; on the other hand, we may well presuppose that he has had occasion to read his own dogmatics, just as we do not believe we are making excessive demands when we assume that, alongside going through his dogmatics, he could have found occasion to go through the four pages with the reduction in my reply to \textit{Dr.} Zeuthen.

It would have greatly advanced the discussion if Bishop Martensen had in one way or another, by some
% --- p. 86 ---
\opage{86}\setcounter{footnote}{0}%
striking example or other, shown me wherein my misjudgment of what I must call his dogmatic explanations of miracles actually consists.

Is there perhaps nothing here in ``the Christian dogmatics'' that might point to the ``science of revelation'' nonetheless, where it sees its chance, dabbling a little in explanations of miracles? To begin, for a change, with the third fundamental miracle, here is the miracle of Pentecost. Dare ``the Christian dogmatics'' now raise its three fingers and swear that it knows itself free, that as regards the miracle of Pentecost it has not occupied itself with anything that might resemble an explanation of a miracle? Dare it say ``by God'' to that? I do not believe it. For it really does stand in the dogmatics --- if Bishop Martensen does not remember it so exactly, he can look it up himself, p. 347 § 186: ``That speaking in tongues was thus the expression of the first enthusiasm, still breaking forth without rule and as it were overflowing, in which the apostles felt their whole natural existence thrilled through by the Spirit of God as by a lightning flash, felt themselves overwhelmed by the exceeding power,'' etc. Here, without circumlocution, is a scientific explanation of a miracle, and moreover an explanation which, when the words ``breaking forth without rule and as it were overflowing'' are duly accented, is not far from explaining ``the divine mode of operation itself, or the How of the wonder.'' Indeed, how seriously the How of the wonder is meant can be seen from the fact that dogmatics, in order to be able to explain ``the divine mode of operation itself, or the How of the wonder,'' does not hesitate to turn the whole miracle upside down; for it says in plain words (p. 137 § 186): ``When, on the other hand, the others present heard them speak of the great works of God, each in his own tongue, then it is'' --- for that Luke, Acts 2:3, says expressly that the apostles began to speak in \emph{other tongues} (ἑτέραις γλώσσαις) does not come practically into consideration here, where the question concerns the higher theoretical \emph{How} ---
% --- p. 87 ---
\opage{87}\setcounter{footnote}{0}%
``not necessary here to think of different languages, but of this, that each for himself felt addressed in his own innermost \emph{peculiarity}, that the inner man was so'' --- say, then, that there is no \emph{How} here --- ``addressed by this tongue of rapture that the soul'' --- here, surely, one sees ``the divine mode of operation itself'' --- ``felt itself released from its accustomed, its natural limitation and in a wondrous manner became manifest to itself.''

The reader will see from this that Bishop Martensen, in wishing to disagree with me, comes to disagree with his own dogmatics. Does theology occupy itself with explanations of miracles? Bishop Martensen answers with a decided No! But dogmatics and I answer with an equally decided Yes! Against me Bishop Martensen can easily hold his own; but can he now also hold his own against his own dogmatics?

Perhaps, in order to mediate a settlement between the Bishop and his dogmatics, one will lay quite extraordinary weight on the circumstance that the \emph{How} of the explanation is not properly metaphysical but merely psychological; concerning the psychological explanations one will perhaps apply to ``the Christian dogmatics'' what an ingenuous examination candidate, who was asked why his reader had \textit{mon amie} and not \textit{ma amie}, is once supposed to have said about French grammar: ``the French do not take it so strictly''; but that is an insult to dogmatics. As proof that dogmatics takes it strictly even with the metaphysical How of the miracle, what was said about the second fundamental miracle, the incarnation, in the reply to \textit{Dr.} Zeuthen is quoted verbatim.

``\,`When one,' says dogmatics, `has found it self-contradictory that the eternal Logos becomes man ``because the Eternal and Omnipresent cannot let himself be born in the midst of time,'' then it surely cannot be the task to think that the eternal Logos should, with the incarnation, cease to exist in his general revelation in the world, or that the Logos as a self-conscious personal
% --- p. 88 ---
\opage{88}\setcounter{footnote}{0}%
being should be enclosed in the mother's womb, be born as a child, grow in knowledge, etc., for this notion cancels the very concept of birth. Birth in time necessarily entails progress from the unconscious to the conscious, from possibility to actuality, from seed-grain and germ to full-grown organization. That the divine Logos lets himself be born means, therefore, that he sinks himself into the womb of humanity as possibility, as holy seed, in order to be able to rise up in humanity in mediating and saving human revelation.' Here, then, we have the explanation. But should we not easily agree that this explanation itself needs an explanation? For what is it, really, that the divine Logos has `sunk' into `the womb of humanity,' himself or not \emph{himself}? If it is not his proper self, then neither is it the Logos himself that is incarnated. The being whose proper self cannot be `sunk' as possibility into the mother's womb cannot be incarnated. But if it really is himself, his possible, essential self, yet not the self-conscious personal self from eternity, that is spoken of here: then the divine Logos must indeed have a double self, one that can be incarnated and one that cannot be incarnated. The Logos that can be incarnated is thus in himself, in his proper self, another than the eternal personal Logos himself: in other words: according to the explanation of `the Christian dogmatics,' the eternal personal Logos himself has not been incarnated at all.''

Since Bishop Martensen, as regards the quotation from his dogmatics (p. 269), has not demonstrated, and will never be able to demonstrate, any distortion whatsoever, either voluntary or involuntary, no dispute can naturally arise here about the text itself, but only about its interpretation. From what standpoint will the dogmatician have his dogmatic explanation understood here? Seen objectively, i.e.\ from the point of view of the logical-physical, the explanation contains the
% --- p. 89 ---
\opage{89}\setcounter{footnote}{0}%
glaring contradiction that the self which in the man Jesus was born in time of the Virgin Mary, and the eternal Logos-self, are to be both the same and yet not the same self. This contradiction can be resolved in two altogether different ways, according as it is resolved either in knowledge or in faith.

If the contradiction is to be resolved in knowledge, then all ambiguities must go; that is the first requirement. Among the conspicuous ambiguities are propositions such as this: ``That the divine Logos lets himself be born means that he sinks himself down into the womb of humanity as possibility, as holy seed,'' etc. But when all ambiguities are removed, what then becomes of the fundamental miracle? A fundamental contradiction, which under scientific analysis develops into a swarm of insoluble contradictions! If the many contradictions intolerable in science are to be removed, then the fundamental contradiction itself must be removed; but seen in the light of knowledge the fundamental contradiction is one with the fundamental miracle. In getting rid of the contradiction, science gets rid of the miracle; the explanation of miracles ends in the denial of miracles: this is the scientific solution of the knot.

If the contradiction is to be resolved in faith, then it by no means becomes the task to determine how the metaphysics of the incarnation is either to be thought or not to be thought, for the principle of faith rises above the logical-physical only by excluding all objective-\emph{logical-physical} and thereby at the same time all \emph{metaphysical} thinking. In the person of Christ the unity is presupposed, the original unity, given in the conception, of divine and human nature, divine and human will. In the light of this unity religious thinking apprehends and determines everything the Gospel teaches us about the life of Jesus, but the unity itself it does not think through: the unity is and remains an \emph{absolute mystery}. The first thing religious thinking has to do is to get rid of all dogmatic problems of the incarnation. For the believing
% --- p. 90 ---
\opage{90}\setcounter{footnote}{0}%
consciousness the contradiction of the fundamental miracle is quite certainly resolved; but not by its itself \emph{thinking} the resolution; no, by its \emph{believing} that the resolution is thought in the mystery.

If we might take Bishop Martensen at his word when he rejects all explanations of miracles, then we might also count on his assistance when we earnestly set about rejecting the half-scientific quackeries of ``the Christian dogmatics'' with the metaphysics of the problem of the incarnation. The problem of the incarnation is and remains a miracle problem; a miracle problem cannot be solved, not even to a certain degree, unless one, at least to a certain degree, occupies oneself with an explanation of miracles. Now if it is really Bishop Martensen's earnest intent, and so earnest a matter must of course be taken earnestly, to lodge a protest against all explanations of miracles and thereby against all science of miracles: then we are entirely agreed, and if we two are only agreed on this one point, we shall, I hope, surely get the better of dogmatics. For what does dogmatics want with its explanations of miracles? It wants --- let us just say it, in order to prove it --- it wants, to put it plainly, to \emph{talk drivel}. While we now, relying on an episcopal word: no explanations of miracles! intend to press ``dogmatics'' hard and charge it outright with \emph{drivel}, we must nevertheless note at once that we by no means forget that there is a difference between dogmatic drivel and other drivel. In order to distinguish dogmatic drivel from other drivel, we shall not use the word drivel of the dogmatic kind, but rather introduce a new category: the \emph{phlyaric}, or, as we might also call it, the phlyaretic, but not the phlyaristic (from φλυαρέω); and then ask: \emph{in what can the phlyaric or phlyaretic of the ``Christian dogmatics'' as such now properly be said to consist}? It can be said to consist in a seemingly profound, but fundamentally self-contradictory, insane and meaningless union of faith and knowledge. The phlyaric in the union
% --- p. 91 ---
\opage{91}\setcounter{footnote}{0}%
comes forward and manifests itself in a great number of \emph{explanations of miracles} that are unhealthy, extremely unhealthy, for science as well as for religion. Among examples of the many explanations, the explanation of the miracle of the incarnation deserves a special place. For closer understanding, the phlyaric is posited first in its \emph{negative} and then in its \emph{positive} moment; the mediation of the negative and positive moments then finally gives us the speculative unity.

1) \emph{The phlyaric in its negative moment.} Faith is presupposed, and then off it goes at knowledge: \textit{credo, ut intelligam}. The peculiarity of dogmatic knowledge must naturally be determined by the peculiarity of dogmatic thinking. Dogmatic thinking is very peculiar; it characterizes itself from first to last by this, that after having rejected the thinkable it ends by thinking the unthinkable. The metaphysics of the incarnation affords an excellent paradigm. The incarnation, namely, must not be thought in such a way that the eternal Logos has sunk himself, his eternal, unchangeable, ever-conscious self, into the womb of humanity: ``That the Son of God is in the mother's womb not as a self-conscious divine I, but as an unformed fetus, Scripture also indicates,'' etc. (Christl. Dogm. § 132). But neither must the incarnation be thought in such a way that the self, the self that begins in the mother's womb as an unformed fetus, should not be the eternal Logos-self. ``Although Christ testifies: I and the Father are one, he never says, I and the Logos are one. For he is the human \emph{self}-revelation of the divine Logos.'' (Dogm. § 132). That the phlyaric is here posited in its negative moment is recognizable from the result, for the result is completely meaningless. The last glimmer of possible meaning vanished with the remark that Christ, who testifies: I and the Father are one, never says: I and the Logos are one. Had Christ but one single time said, in the Johannine language: I and the Logos are one, then
% --- p. 92 ---
\opage{92}\setcounter{footnote}{0}%
the metaphysics of the incarnation would at least have had a distant prospect of a kind of meaning. One could then at least have tried to conceive an identity of the self that is conceived and formed in the mother's womb with the eternal Logos-self, corresponding to the identity of the Father with the Son, and in this way squeeze half a thought into the words of dogmatics: ``But the Holy One who is born of Mary, and who in the course of his growth becomes conscious of himself as a human I, becomes in the same measure conscious of his deity; knows himself as a divine I, because the fullness of the Godhead is the \emph{ground of life} of his human living; knows himself not merely as partaking in the divine Logos, but as the divine-human continuation of the eternal divine life from the beginning.'' (Dogm. § 132). This last way out is now completely blocked. Christ, after all, cannot say: I and the Logos are one. Can he really not? No, the great dogmatics knows that he cannot! And why not? Because Christ's human I, the I that has a beginning in time, has developed from an unconscious to a conscious --- and the self-conscious divine I, the eternal Logos-I, the I that has no beginning in time, has not developed from an unconscious to a conscious --- are not two identical I's, but literally one and the same I. Surely no thinking can go further --- in the phlyaretically phlyaric.

2) \emph{The phlyaric in its positive moment.} The peculiarity of dogmatic thinking is that it thinks scientifically in faith: \textit{credo, ut intelligam}. That there is scientific thinking in faith means that there is no thinking, but representing, intuiting, fantasizing in a metaphysical manner, for the elucidation of the mysteries, for the benefit of revelation, for the joy of faith. Dogmatics, which objectively builds on the word of Scripture, sees in the Scripture passages something quite other and deeper than either Christ himself or his apostles ever saw: it sees the metaphysical. ``He would,'' it says (Dogm. p. 133) of Christ's self-abasement, ``not merely
% --- p. 93 ---
\opage{93}\setcounter{footnote}{0}%
live in the pure glory of the Godhead, in metaphysical majesty; but he emptied out his fullness in the lowly form of a servant in order to be able to become the head of the kingdom of love, the reconciler and redeemer. But this self-abasement must at the same time'' --- here comes the proper positive explanation of a miracle --- ``be seen as his self-perfection,'' the self-perfection, namely, of the Logos who is in himself eternally perfect. In that he, the one and same Logos, ``lives in a human manner, and as the Son of Man possesses his deity only in the limitation of human individuality, in the bounded forms of human consciousness: he must surely be said to live in abasement and poverty, since he has renounced the majestic glory in which, as the omnipresent Logos, he shines through all creation.'' (Dogm. § 133). If we now note the two expressions: ``he has renounced'' and ``he shines through,'' and remember that it is the same, the selfsame Logos, the one and selfsame Logos-I, that has here renounced the majestic glory which he has nevertheless not for a single moment renounced, how then could we doubt that it is precisely the higher phlyaric, the deeper phlyaretic, that is by this explanation posited in its positive moment?

3) \emph{The phlyaric in speculative completion.} The mediation of faith and knowledge can naturally be carried through only on the condition that both parties approach each other with all possible diplomatic courtesy, that is, in all sincerity and disinterestedness, animated only by the warmest wishes for peace and good understanding. During the mutual rapprochement a happy change takes place in them both: faith becomes genial, knowledge good-natured. That faith becomes genial means that it gives up its old-fashioned orthodox touchiness, is ashamed of what is to the Jews a stumbling block and to the Greeks foolishness, smiles at philosophy, and falls in love with speculation. It can naturally not
% --- p. 94 ---
\opage{94}\setcounter{footnote}{0}%
be any ordinary faith, not the popular, simple asylum-faith, that carries on so charmingly; no, it is a cultivated, well-bred faith, a faith that knows its mystics and has heard lectures at the Institute; only such a faith, a subjective faith with an objective heart, a faith that would not love practically if it did not love theoretically, can open its arms, its deep, dogmatic arms, to the higher knowledge. That the knowledge which is to be able to give itself without nausea to so cultivated, so Institute-cultivated a love must absolutely be \emph{good-natured} is surely self-evident. For there is here a difference, a very considerable difference, between good-natured and not good-natured knowledge. The not good-natured knowledge we may with reason call the teasing kind, namely the dogmatically teasing kind, because it always takes its pleasure in teasing dogmatics. The teasing consists quite simply in this, that with merciless consistency it pursues every contradiction in which dogmatics in its ingenuousness has become involved. Out of a couple of innocent ``must-not-be-thoughts'' it conjures up, by means of the logical-physical, a demonic swarm of ``must-not-be-thoughts''; by critical-dialectical witchcraft it gets the dogmatic contradictions, if we may say so, to breed. Indeed, what is most outrageous: precisely the deepest positive, that of which dogmatics expressly says that \emph{it must be thought}, it transforms into a restless negative, something thoroughly unthinkable, a ``cannot-be-thought'': that is malicious. The malicious knowledge, the knowledge that will not let go of the logical-physical, is a sinful knowledge, a knowledge that cannot say \textit{credo, ut intelligam}, an \emph{unbelieving} knowledge. It is otherwise with the good-natured knowledge: blessed are they who will not criticize the good-natured knowledge. The good-natured knowledge does sometimes look somewhat thoughtless; but, by Dogmatics, it is not thoughtless, for it is lost in thought. Why is it lost in thought? It is in love, subjectively in love with the objective heart of faith: it is, after all, a believing knowledge. It

% --- p. 95 ---
\opage{95}\setcounter{footnote}{0}%
believing knowledge is a knowledge that longs --- out beyond the logical into the supralogical: miracle-metaphysics is the goal of its longing. If it does not reach the goal of its longing, it dies; but if it reaches the goal of its longing, it lives, lives on the positive of faith, scoffs at the negative of criticism, and lulls all contradictions to sleep in absolute joy over the relative mysteries. Should one or another contradiction be so unruly as to scream out from the nest of the mysteries, what does that matter? ``The French do not take it so precisely.'' Believing knowledge is good-natured: good-naturedly it sinks to the bosom of dogmatics, good-naturedly it rests at the heart of faith. Practically, good-naturedness is love; theoretically, love is good-naturedness. If in the practical sphere it must be said that love covers the multitude of sins, then here, where the deepest practical is at the same time the highest theoretical, it must be said that good-naturedness, the incredible good-naturedness of believing knowledge, covers the multitude of contradictions. On these terms the union of faith and knowledge is accomplished; the phlyaric is set in speculative completion.

In this result I may well venture to think that I have Bishop Martensen with me; I rely on his word that dogmatics must cease with its explanations of miracles; I am certain that he will acknowledge all the consequences flowing from his presuppositions, as certain as if I already had them secured by episcopal confirmation. Only one scruple still remains, namely this: that I have not understood his words altogether correctly. For it is indeed conceivable that Bishop Martensen does not want his rejection of the theological explanations of miracles understood to mean that theology straightforwardly renounces the explanations of miracles which it actually gives up, but to mean that in its abasement and poverty it ``renounces'' the explanations of miracles with which, in its ``metaphysical glory,'' it nonetheless continues still to ``illuminate'' dogmatics. If that is the meaning, then it is another matter; then
% --- p. 96 ---
\opage{96}\setcounter{footnote}{0}%
we dare expect no confirmation, but must, whatever the outcome may be, venture out on our own. We now pass on to a treatment of our last, and the system's first, great explanation of a miracle: the explanation of \emph{the miracle of creation}.

Bishop Martensen expresses himself in his occasional piece (pp.~59--60) on my conception of the concept of creation in the following manner. ``We shall not dwell on the author's presentation of the theological concept of creation, which is very incomplete and incorrect. We shall only remark in general that creation unconditionally presupposes the ethical concept of God, and that without this it is altogether absurd to speak of creation. The concept of creation is therefore presented altogether wrongly when one, as Professor Nielsen in many places likes to present it, determines it as a pure act of arbitrariness on the part of the divine omnipotence, whereby one remains standing at a merely physical determination, whereas in the concept of creation there is and can be question only of the omnipotence determined by love and wisdom. If one wishes, therefore, to criticize the concept of creation, then let one first and foremost criticize the ethical concept of God itself and show that it is impossible for omnipotence to be ethically determined. If one has proved this --- but hitherto the proof has failed to appear --- one can spare oneself the trouble with the concept of creation as well as with all the other dogmas. If one attacks the dogma of creation by polemicizing against naked omnipotence and a mere act of arbitrariness, one certainly makes the matter very easy for oneself, but one also merely fences in the air, to one's own amusement and that of the ignorant.''

To judge by this testimony, my case stands but poorly. My hope, my beautiful sanguine hope, that Bishop Martensen himself would have helped me against dogmatics, has now as good as vanished; only a glimmer is left; it glimmers in a note (Leilighedsskr. p.~60):
% --- p. 97 ---
\opage{97}\setcounter{footnote}{0}%
``Theology, as well as Christian philosophy and theosophy (J. Böhme, Fr. Baader), has at all times recognized in creation an unfathomable mystery, and has never claimed to be able to fathom the How of creation.'' A very pretty little note, quite as if plucked out of Grundideernes Logik; I lean upon it and continue to hold that dogmatics ought to have a spanking when it is so naughty as to want to fathom the unfathomable, to claim to be able --- naturally only to a certain degree --- ``to fathom the How of creation.'' Only listen to the lofty measure in which dogmatics begins its §~59 on creation: ``That God creates implies that he brings forth what is not God, that whose essence is different from his own, a free finitude, which he wills to fill with his fullness.'' If it were a concept of faith absolutely heterogeneous with the concept of knowledge that was aimed at here, then from our point of view we should have nothing to object; but how should ``the highest science,'' with its ``good will,'' ever be able to stoop to separating concepts of faith from concepts of knowledge? If dogmatics conceded the validity of a development of concepts of faith absolutely heterogeneous with any development whatsoever of concepts of knowledge, then the dogmatic Christ, although he himself is both Theos and Logos, could not become Theologus, hence not Summus Theologus either, hence not a king of the sciences either. No, by the royal principle dogmatics will know how to assert and preserve its lofty aim.

That the dogmatic concept of creation is taken without further ado as a well-grounded concept of science can be seen from categorical determinations such as: ``Just because God is love, he cannot close himself up within himself as the God of the ideas, but must determine himself as spirit,'' etc.; further: ``If it can be said that God creates the world in order to satisfy a want in himself, then this must, in accordance with the concept of love, be understood in such a way that this
% --- p. 98 ---
\opage{98}\setcounter{footnote}{0}%
lack is just as much a \emph{superabundance}. For this lack in God is not, as in the God of pantheism, the blind hunger and thirst for existence, but is one with the exceeding richness of freedom, which can do nothing other than will to reveal itself.'' These few sentences alone, in which logical and illogical, physical and hyperphysical, necessity and love, lack and superabundance are stirred together in speculative ambiguities, are sufficient to show that science has become believing, and faith scientific.

In lumping the concept of creation together, to the best of its ability, with the concept of divine love, ``Christian dogmatics'' behaves like a well-studied coquette. It will not say straight out that the act of creation is an act of necessity; for it is coquetting with faith. It will not say straightforwardly that the act of creation is an unconditioned act of freedom; for it is coquetting with knowledge. What does it say, then? It says that the act of creation is a divine act of love. Into love, faith can now put as much freedom as it likes; for love is indeed free: its want is satisfaction, its lack superabundance. ``If it can be said that God creates the world in order to satisfy a want in himself, then this must, in accordance with the concept of love, be understood in such a way that this lack is just as much a \emph{superabundance}.'' Into love, knowledge can put as much necessity as it likes; for love is moved by holy necessity: its satisfaction is want, its superabundance lack. ``Just because God is love he cannot'' --- mark it well: he \emph{cannot} --- ``close himself up as the God of the ideas, but must'' --- mark it well: he \emph{must} --- ``determine himself,'' etc. ``Christian dogmatics'' dares not say straightforwardly that the world is created out of God's essence\footnote{A man like Jacob Böhme dares to: ``Many writers have indeed written that heaven and earth were created out of nothing; but I wonder that among so many excellent men there has not been a single one who could describe the true ground, since after all the same God who now is has been from eternity. Where there is nothing, nothing comes to be either; every thing must have a root, otherwise nothing grows; had not the seven spirits of the divine nature been from eternity, then there would have come to be no angel, nor heaven, nor any earth either. --- When God created both this world and everything else, he had no other material out of which he could create it than his own essence, himself.'' Den christel. Troesl. by D. F. Strauß, trans. by H. Brøchner. Vol.~1, p.~543.}, for thereby it would
% --- p. 99 ---
\opage{99}\setcounter{footnote}{0}%
offend faith; but it is indeed coquetting with faith. What does it say, then? It says that the world is in a certain manner created out of ``the non-existent,'' out of the ``eternal possibilities'' of the divine will, hence out of the lack that is the superabundance of love. And now faith can put all the nothingness it needs for a \textit{creatio ex nihilo} into this non-existent, into these possibilities, into this lack. But ``Christian dogmatics'' dares not say straightforwardly either that the world is created out of nothing; for thereby it would offend knowledge; but it is indeed coquetting with knowledge. What does it say, then? Then it says, so as not to compromise ``naked omnipotence,'' that the nothing out of which the world is created is naturally not $=0$, ``against which representation the rule holds: \textit{ex nihilo nihil fit}''; but that it is a \emph{positive} nothing, a nothing rich in content and full, a nothing of ``magical realities,'' ``plastic archetypes,'' ``a heavenly host of visions,'' the lack in the superabundance of love with which alone the creating God can satisfy ``a want in himself.''

If any reader should doubt that the mediation described here is in categorical agreement with the dogmatic text, we ask him to compare what has here been rendered under various paraphrases with what stands literally in §~57, §~59, §~60, §~61 of the Dogmatik. One
% --- p. 100 ---
\opage{100}\setcounter{footnote}{0}%
will not be able to say that the heterogeneous elements of the concept of creation are worked together either logically or ontologically; they are merely stirred together in all ambiguity\footnote{Since Bishop Martensen (Leilighedssk. p.~114), on the occasion of my attacks on the dogmatic explanations of miracles, refers me to Strauß, who puts forward his objections ``with greater clarity and less passion,'' I will in this connection take the opportunity to recall what testimonial Strauß has given to the mediation which we here designate as a stirring together of heterogeneous components. ``Not everyone,'' says Strauß (Den christelige Troeslære, translated by H. Brøchner, Vol.~1, p.~61), ``possesses the apparatus or the perseverance with which Schleiermacher, in order to effect a mixture, pulverized Christianity and Spinozism so finely that it takes a sharp eye to distinguish the mixed components from one another; with many the Christian and the modern element are no better united than oil and water shaken together, which appear united only so long as one keeps on shaking them; while yet other dogmaticians, and some of them not undistinguished, even resemble a sausage-meat (the image is no more ignoble than the thing itself), in which the orthodox doctrine of the Church makes up the meat, the Schleiermacherian theology the fat, the Hegelian philosophy the spice. It is these mixtures, in which what is spoiled is supposed to be made palatable again by all sorts of additives, that Lessing already found so disgusting, so repulsive, so nauseating; it is the so-called rational Christianity, of which he confessed that he knew neither where its reason nor where its Christianity lay; the modern theology, to which he preferred the old orthodox one for the reason that the latter openly conflicted with sound human understanding, while the former would rather bribe it.''}. One can hardly even say that the mediation consists in a kind of agreement or concordat; for the ambiguities are in far too fluid a state to be able to preserve even so much as an outline of anything concordat-like; rather one might say, when one abstracts from the quasi-aesthetic form of the presentation, that what is
% --- p. 101 ---
\opage{101}\setcounter{footnote}{0}%
stirred together is mediated into a half-clear gruel, an ethical-hyperphysical gruel of love.

With such a mediation it must surely become evident that the dogmatic presentation, despite its exertions to give itself a scientific appearance, is heterogeneous with science; for gruel and science are heterogeneous. The concept of love, the highest ethical concept of God, is a concept of faith which dogmatics maltreats within science. In order to make the maltreatment recognizable, we have, in attacking the dogma of creation, polemicized not against ``naked omnipotence'' but against painted impotence. One need only lift a corner of the veil of ambiguity to discover the swarm of contradictions that scream behind the veil. To lull these contradictions to sleep by throwing the cloak of love over them does not help, any more than it helps to give them gruel. There is only one thing that helps, and that is: to clear out the nest. It may pain the fatherly heart --- it pains me too, after all ---, but away with untimely leniency! Dogmatics is effeminate; dogmatics is spoiled; dogmatics must go to school; dogmatics must have the rod until it breaks the habit of speculating in ambiguities, of presuming upon philosophy, and of mediating in gruel.

Bishop Martensen will not dwell on my presentation of the theological concept of creation, which is ``very incomplete and incorrect.'' Why, then, has the Bishop not been willing to show me wherein the incomplete and incorrect consists? When no one better informed will take charge of my dogmatic upbringing, it is no wonder, after all, if I shift and shift without hitting upon the right thing.

I am reproached with conceiving of creation as an act of arbitrariness, a work of ``naked omnipotence''; but this reproach does not strike me: it strikes the old teachers of the Church, it strikes the Scholastics of the Middle Ages, it strikes all our Protestant-orthodox dogmaticians. If
% --- p. 102 ---
\opage{102}\setcounter{footnote}{0}%
the omnipotence that brings forth a world out of nothing is to be called ``naked omnipotence,'' then I can call the Fathers --- all the way from Irenaeus\footnote{Read Irenaeus \textit{adv. hæreses II, 10, 4: homines quidem de nihilo non possunt aliquid facere, sed de materia subjacente. Deus autem materiam fabricationis suæ, cum ante non esset, ipse adinvenit.} Read Augustine. \textit{Confess. XII, 7: Fecisti coelum et terram non de te; nam esset æquale unigenito tuo --- et aliud præter te non erat, unde faceres: et ideo de nihilo fecisti coelum et terram.}} to Quenstedt, and among them one who will surely make an impression on Bishop Martensen, namely Saint Thomas\footnote{Compare \textit{Summa I, 45, 2: Antiqui philosophi non consideraverunt nisi emanationem effectuum particularium a causis particularibus, quas necesse est præsupponere aliquid in sua actione. Et secundum hoc erat communis eorum opinio, ex nihilo nihil fieri. Sed tamen hoc locum non habet in prima emanatione ab universali rerum principio \dots} with what follows: \textit{Si ergo Deus non ageret nisi ex aliquo præsupposito, sequeretur quod illud præsuppositum non esset causatum ab ipso. Ostensum est autem supra, quod nihil potest esse in entibus, quod non sit a Deo, qui est causa universalis totius esse. Unde necesse est dicere quod Deus ex nihilo res in esse produxit;} one will then not be able to doubt that the \textit{universale rerum principium} of which the Scholastic speaks is precisely the \textit{liberrimæ voluntatis imperium} of which Quenstedt speaks.} --- as witnesses that it is not my concept of creation, but Bishop Martensen's, that is ``very incomplete and incorrect.'' When a Quenstedt, as I have in fact cited by way of example (Grundideernes Logik II, p.~150), defines creation as: \textit{Actio Dei triuni externa, qua is res omnes visibiles ex nihilo \dots\ solo liberrimæ voluntatis suæ imperio omnipotenter et sapienter produxit}, does he not lay the decisive weight on creation's being an unconditioned act of freedom, a work of omnipotence? That love and wisdom must be assumed to co-operate in everything
% --- p. 103 ---
\opage{103}\setcounter{footnote}{0}%
that omnipotence does is self-evident; but it stands no less firm for all that, that it is omnipotence, not love, the unconditionally free act of will, not wisdom, on which the emphasis falls when the question is of a determination of the concept of creation. It is, as any candidate in theology can know, altogether backwards to want to derive the concept of creation from the ethical concept of God instead of developing it from the concept of omnipotence.

When Bishop Martensen wishes to press the problem of creation into the question of the possibility ``that omnipotence can be ethically determined,'' and that with a bang-effect such as this: ``If one attacks the dogma of creation by polemicizing against naked omnipotence and a mere act of arbitrariness, one certainly makes the matter very easy for oneself, but one also merely fences in the air, to one's own amusement and that of the ignorant,'' then even the less initiated will smell a rat and get an inkling that the dogmatician means to use his ``ethical concept of God'' as a shield in order to hack his way through and get away from the dangerous problem.

In Grundideernes Logik II, pp.~149--166, the concept of creation is treated, certainly exclusively from the standpoint of logic, but nevertheless with express regard both to the orthodox doctrine and to the Martensenian speculation. It is shown there, through a series of analyses, how double-tongued, self-contradictory and rickety the dogmatic-speculative treatment of the dogma of creation is from first to last. There are points enough in these analyses to take hold of. If Bishop Martensen has found in this whole section either distortions or the passing over of essential points or other blunders: why, then, has he not expressly objected to it? If Bishop Martensen really has such confidence in his ethical concept of God that it can solve the problem of creation for us in dogmatics, and that it is only logical blindness that has led me ``to fence in the air'': why, then, has he not in a definite manner disarmed the definite objections of logic by an altogether definite application of the ethical concept of God?

% --- p. 104 ---
\opage{104}\setcounter{footnote}{0}%
I will not weary the reader with either copying out or paraphrasing the objections made in the Logik, but will leave it to those concerned to look up the cited passages themselves, in Bishop Martensen's Dogmatik as well as in my Logik, and will only, in conclusion, call attention to what is typically theological in wanting, by general reflections on the sum and the result, to evade the actual reckoning and to justify such evasions with the high-flown word: total view.

After these preparations we now pass on to a consideration of the position that philosophy will come to occupy in the monarchy. Since the word philosophy is ambiguous, we must remark at once that it is philosophy as \emph{science}, exclusively as science, that is spoken of here. As science, philosophy is, with respect to its absolute content, the doctrine of ideas, and as the doctrine of ideas a distinctive science, different from the particular sciences, the so-called special sciences.

But with all their distinctiveness and mutual difference, all sciences without distinction must nevertheless rest on a common scientific foundation. The common scientific foundation is determined by rational necessity and does not reach beyond the logical-physical. On account of the foundation, it is impossible to conceive of absolutely heterogeneous sciences. Just as little as there can be a philosophy absolutely heterogeneous with the exact and positive sciences, just so little can there be two kinds of philosophy, of which the one should be absolutely heterogeneous with the other. But when philosophy cannot possibly be split into two absolutely heterogeneous philosophies, when it cannot possibly develop from two absolutely heterogeneous principles, then it can already be seen from this that its position in the monarchy cannot become enviable.

By referring us to Paradise, to ``\emph{the two absolutely heterogeneous trees}, of which the fruits of the one were forbidden to man,'' Bishop Martensen refers us to
% --- p. 105 ---
\opage{105}\setcounter{footnote}{0}%
two absolutely heterogeneous principles of knowledge and of science, the development of one of which must absolutely be forbidden to us human beings. Here we have again the Elijah-like, or if one will, the Maccabaean; we have ``the good will.''

What is at stake is a cleansing of the sanctuary of philosophy. First of all, the old Greek philosophy must be swept out; for it is certain that neither Plato nor Aristotle ever thought of going beyond the logical-physical; it is certain that neither Plato nor Aristotle would acknowledge a knowledge that is determined not with rational necessity but by ``the good will''; it is certain that Greek knowledge is autonomous, not theonomous; it is certain that Greek philosophy will never ever be able to be transformed so as either itself to become a philosophy of faith or a second time to bow under theology as a handmaid under her mistress, as a Hagar under a Sarah. It was only in the Middle Ages that Plato was ``der große Pfaff''; it was only in the Middle Ages that Aristotle passed for a \textit{christianus ante Christum}; it was only in the Middle Ages that theology was Scholasticism, and that Scholasticism screwed together the absolutely heterogeneous. Shall we then call back the Middle Ages? No, classical philosophy must get out of the monarchy.

But modern philosophy from Descartes to Hegel must, by Maccabaeus!, also get out of the monarchy. Cartesian philosophy indeed begins precisely with a principle that stands in absolute conflict with \textit{credo, ut intelligam}, namely: \textit{de omnibus dubitandum est}. A philosophy that begins with doubt and a philosophy that begins with faith are incompatible. Had it been possible to unite a philosophy of faith with a philosophy that has doubt as its methodical point of departure, then we would indeed by now have had, in our own literature, the first foundation for such a union in J. L. Heiberg's development of the Hegelian philosophy. ``The Hegelian philosophy,'' says Heiberg (Perseus No.~1, 1837,
% --- p. 106 ---
\opage{106}\setcounter{footnote}{0}%
p.~40), ``rests not merely on God's revelation in general (for all the other philosophies do the same), but it rests on the Christian revelation in particular, and it has this presupposition with a consciousness that comes from the end to the beginning, and precisely for that reason can absolutely ignore itself without losing or forgetting itself.''

Heiberg, precisely in the treatise cited, makes a well-founded objection against the Martensenian project of a distinct philosophical principle of faith. He says: ``The demand which Mr. Martensen makes of a philosophy of faith seems already to be fulfilled in the Hegelian one. For that this one nonetheless, at the beginning of its system, makes itself ignorant of what it knows, and a doubter, indeed even a denier, of the God whom it recognizes, happens only in order to fulfil the simultaneous demand: a presuppositionless beginning, and can hardly be otherwise with any philosophy that is really to answer to its name, and thus is not to be religion. \dots\ Should Mr. Martensen object to this that, even if the spirit of God is the final result of the Hegelian philosophy, this is presented only logically, not with the whole fullness of the Christian revelation for the heart or feeling, then I must reply to this that precisely therein consists the eternal difference between philosophy and religion, for if this difference were not eternal, philosophy would indeed become everything, and religion, as well as art, superfluous. But before I enter upon this, it will be expedient to await the closer development which Mr. Martensen has given us hope that he will bestow on us.''

The promised development has now been bestowed on us; we have its ripe fruits in ``our Christian dogmatics.'' The Martensenian Dogmatik is an example of what comes of passing Descartes smoothly by and attaching oneself to Anselm, of rejecting the philosopher's \textit{de omnibus dubitandum} and holding fast to the Scholastic's \textit{credo, ut intelligam}. It is nothing other than wanting to go beyond the logical-physical,
% --- p. 107 ---
\opage{107}\setcounter{footnote}{0}%
wanting to get away from rational necessity, wanting to erect a scientific doctrinal edifice outside the bounds of knowledge, hence outside the scientific foundation. But scientific doctrinal edifices without a scientific foundation are called castles in the air.

Why, then, does philosophy begin with a doubting of everything? For the sake of investigation, testing and assurance, so that rational necessity may come to displace doubt! Science knows no certainty, esteems no certainty, but that which is forced upon consciousness in a logical-physical manner. A proposition that I can both accept and not accept is without scientific value; only the proposition which I cannot possibly reject without rejecting science itself is scientific. ``But,'' asks Bishop Martensen (Leilighedssk. p.~102), ``if only that is to be called science which can be forced upon one with irresistible necessity, then science sinks down to something very empty and worthless, which it is not worth the trouble to make a great fuss about. Only formal logic and mathematics then become sciences in the stricter sense.'' This declaration betrays so strangely great a misjudgment of what gives science and scientific character, in all their forms, value and significance, that among the cultivators of the most diverse sciences there will hardly be anyone who does not immediately and at first glance see the misunderstanding in its whole extent. For what is it that a historian must strive after just as eagerly as a mathematician, a botanist just as eagerly as a mechanician, an anthropologist just as eagerly as an astronomer? It is \emph{certainty, incontestable certainty}! How and on what conditions is this certainty reached? It is reached in the various sciences by various paths and on various terms; but one thing is common to all scientific certainty, this one thing: that it rests and must rest on universally valid,
% --- p. 108 ---
\opage{108}\setcounter{footnote}{0}%
incontestable grounds, that it must ``force itself upon one with irresistible necessity.''

In support of his declaration the Bishop remarks: ``Already a speculative logic goes beyond the merely compelling; for no one will, e.g., be able to grasp what is great in Hegel's Logic or in a work such as Kant's Critique of Reason without architectonic imagination.'' This remark says nothing. For neither will there be anyone able to grasp what is great in Newton's \textit{Principia} or in a work such as Laplace's \textit{Mécanique céleste} ``without architectonic imagination.'' The assurance that works such as Hegel's Logic and Kant's Critique of Reason can be disputed about endlessly, and that endless misunderstandings are possible here, is, fortunately for science, plucked out of thin air. If experts could dispute endlessly about Hegel's Logic and Kant's Critique of Reason without secure, incontestable results being attainable: what would that prove? It would prove that Hegel's Logic and Kant's Critique of Reason were without scientific value.

If it only stands firm that no philosophy --- philosophy as science is what is spoken of here --- has ever attempted or will ever attempt to go outside the boundary of the logical-physical or to raise itself above the rationally necessary, and this main proposition surely must stand firm: then it is now, as every thinking reader can tell himself, exceedingly easy for me to prove that \emph{the protest which Bishop Martensen has lodged in his occasional piece against my philosophy is a protest against all philosophy}.

By gazing toward Paradise and constantly keeping the two \emph{absolutely heterogeneous trees} before his eyes, Bishop Martensen has found himself able to hold two absolutely heterogeneous principles of knowledge up against each other. There is, then, according to his way of seeing, an unbelieving and a believing knowledge, an unbelieving and a believing science: unbelieving knowledge has ``the evil will,'' believing knowledge ``the good will,'' as its principle. Herein is seen

% --- p. 109 ---
\opage{109}\setcounter{footnote}{0}%
the original. It is the setting up of the two absolutely heterogeneous trees, together with the two absolutely heterogeneous wills, within the horizon of objective knowledge, by which the Martensenian philosophy becomes epoch-making in the history of the sciences.

It is, to be sure, not so much regard for philosophy as for theology that brings about the breakthrough here; yet we hope that the revolution proceeding from the trees and the wills shall turn out to the good fortune and happiness, if not of theology, then at least of philosophy.

As a science of spirit, philosophy has, even in more recent times, had an inclination to assimilate the presuppositions of faith and become a little theological. Heiberg's words, that the Hegelian philosophy rests ``not merely on God's revelation in general but on the Christian revelation in particular,'' are characteristic in this respect. That is a cause for concern. Although the closer attachment of philosophy to theology has, no doubt, on the whole had the consequence that theology (as one sees in dogmaticians such as Daub and Marheineke) came to walk in philosophy's leading-strings, not the reverse, the theological has nonetheless been able to infect philosophy itself to such a degree that Hegelians such as Weiße have carried the doctrine of the Trinity to the point of caricature. In well-founded indignation at such a caricature, Strauß exclaims: ``So is it for this, good Hegel, that you found your profound categories: becoming-other and going-out-of-oneself (Entäußerung), negation and the negation of negation, the sublation of the moments in something higher, that the crude imagination's crudest monstrosities, of which the most insane Gnostic would not have been ashamed, should let themselves be framed in them? Three persons, expressly determined as distinct centers of consciousness, are nevertheless to be absolute; two such persons are to be united in a third person into a unity that is more than a mere agreement; a person, and an ab%
% --- p. 110 ---
\opage{110}\setcounter{footnote}{0}%
solute one at that, is for a time to be lost to itself, to become impersonal, to be split up into an infinite multitude of finite personalities, like a potato plant into tubers, but after a certain time has elapsed, to regain its personality, whereupon at the same time the creaturely personalities that arose in the intervening time are, without ceasing to be many and separate from one another, to flow together into the one personality of the divine spirit. --- Where is the \textit{Symbolum Quicunque?} Hand it to me! I would sooner swear to it ten times than once call our philosopher's propositions anything other than insanity.''\footnote{Troeslære, trans. by H. Brøchner. Vol.~1, p. 435.}
This temptation to go along with faith and explicate itself theologically will, as far as philosophy is concerned, naturally fall away now that philosophy is being driven out of the monarchy. And that will happen categorically. Philosophy cannot possibly split into two absolutely heterogeneous philosophies and still hang together in such a way that the one is subordinated to the other. If a philosophy with \textit{credo, ut intelligam} is to hold good in the monarchy, then philosophy as a science must get out of the monarchy.

By acknowledging theology---as a science that needed a thorough reform, to be sure, but still as a science---philosophy, even as philosophy of nature, was led astray into passing over the actual facts in the same ingenious manner as the mystics and the speculative theologians have always done. In Hegel's philosophy of nature---of Schelling's fantasy-intuitions we will not speak at all---there are to be found conceptions, explanations, and surveys that are by no means unworthy of a Bishop Martensen. Hegel too assumed ``that the natural sciences, although they belong to the sphere of necessity, can in no way be called science in the stricter sense''; Hegel too could hold that natural science was a subordinate science, because it ``had no other grounds
% --- p. 111 ---
\opage{111}\setcounter{footnote}{0}%
to rest upon than induction and analogy.'' This foolish conceit nearly cost philosophy dear. The Schelling-Hegelian philosophy came---not exactly by believing in miracles, but by believing in its own absolutely adequate knowledge of the Absolute---so far away from experience and reality that philosophical teleology was on the point of becoming just as hollow and fantastic as the theological. With the revolution proceeding from the trees, the temptation is over: philosophy as a science goes out of the monarchy.

For a number of years it has been a goal of my endeavors to secure for philosophical inquiries an actual foundation and to get the relation between philosophy and the special sciences somewhat clarified. I must on this occasion once more recall the standpoint maintained in the introduction to Grundideernes Logik I: ``If a philosophy of `the concept' is to be able to secure the reality of the concept, it must surely above all attend to how a real concept is formed; but to what degree the Hegelian philosophy, precisely with respect to this essential point, has miscarried in the formation of concepts becomes conspicuous when one casts a glance at the philosophy of nature. That Hegel did not form his first, elementary concepts of nature under the guidance of any mathematical physics, but in a purely logical-speculative manner, is evident at every point. This is seen not merely from the fact that he expressly combats the designations and notions of mathematical physics; it is seen chiefly from the fact that he has not even found it necessary to secure for himself an authentic understanding of the categories and ways of looking at things that he wishes to combat in the mathematicians.\footnote{Cf. lectures on ``Phil. Propædeut. from the academic year 1861--62,'' pp. 49--75.}
The shortcomings pointed out here do not concern mistakes in this or that particular;
% --- p. 112 ---
\opage{112}\setcounter{footnote}{0}%
such mistakes are always more or less inessential in a philosophical exposition of essential substance;---to make much of the inessential must (cf. Begrebet: Nøiagtighed, pp. 3--11) naturally be left to pedants and pettifoggers---; no, the shortcomings pointed out here concern the method itself, and betray a weakness that is, so to speak, innate to the system. It is the manner in which the categories of nature are from the outset conceived and treated that has miscarried here. The concepts of nature are concepts of reality; but for the concepts of reality to be formed, the various special sciences, each with its observations, analyses, and calculations, must go ahead and show the way. There is in the whole realm of nature not a single conceptual sphere of which it holds that the philosophical thinker, without particular consultation with the special science concerned, should himself be able to show the origin of the categories and find their roots in the soil of reality.''

In order to determine the relation between philosophy and the special sciences, I have in my propædeutic inquiries not been able to confine myself to general discussions of the empirical and the a priori, but have, so far as my powers allowed, had to pursue encyclopedic studies of a not altogether slight extent. That much still remains to be done here, I grant; but I did believe I had brought matters so far that anyone who kept up would think twice before once again serving the reading public a rehash of the Schelling-Hegelian wisdom, such as: ``that the natural sciences, although they belong to the sphere of necessity, can by no means be called science in the stricter sense.''

That Bishop Martensen, when once in an occasional piece he had occasion to lodge a protest against my philosophy, would give me a proper going-over in propædeutics does not come to me unexpectedly. If I have given up on dogmatics altogether, then naturally I ought to be examined sharply in propædeutics. One thing, however, I cannot quite understand. Perhaps I am mistaken;
% --- p. 113 ---
\opage{113}\setcounter{footnote}{0}%
but it seems to me as if the lofty air with which Bishop Martensen means to function as examiner stands out of proportion to the slight degree of preparation with which he has gone to the examination. ``One is compelled,'' says the Bishop (Leilighedssk. p. 97), ``to ask whether'' my ``occupation with dogmatics'' has not entailed for me ``no small neglect of philosophy''; and I am compelled on this occasion to ask whether His Reverence's much occupation with ``concordats'' has not entailed for himself no small neglect of propædeutics, in particular of § 13, § 14, § 15, and § 16. For the episcopal examination in philosophical propædeutics before which I am placed is laid out for an answer to the question: \emph{How shall we introduce so thorough a subordination into the realm of ideas that theology's ``good will'' can be acknowledged as autocrat over all the sciences?} This question I cannot answer. I fail, but in failing I take with me the consolation that I fall as a sacrifice to ``the good will.''

If it goes thus with me in propædeutics, then the reader can tell himself how it must go with me in logic. It is a severe critique that Bishop Martensen has issued against Grundideernes Logik, a critique that is truly worthy of ``the royal principle.'' For we must note that the monarchical philosophy springing from the two absolutely heterogeneous trees has transformed for us not only logic but also critique. Just as in Paradise there were two absolutely heterogeneous trees, so there are now in the theological monarchy two absolutely heterogeneous logics, a theological one and a subordinated one, two absolutely heterogeneous critiques, a theological one and a subordinated one.

How oppressive philosophy's position in the monarchy must become when a theologically superordinate critique is to deal with a subordinated logic is easy to see. The good will has
% --- p. 114 ---
\opage{114}\setcounter{footnote}{0}%
the theory corresponding to the good; the evil will has the theory corresponding to evil: what theory, then, does logic have? This question can be answered only from the superordinate standpoint of theological critique. For on the standpoint of a universally human knowledge one does not know whether one should preferably begin with the theory and say: if the theory is good, then the will is good, for a good theory can proceed only from a good will; or, conversely, begin with the will and say: if the will is good, then the theory is good, for a good will must necessarily produce a good theory. The difficulty increases when the talk is of logic as theory, of the good and the evil will in logic. What will is there now to discover in an Aristotelian logic, a good or an evil one? We answer: the will in the Aristotelian logic is neither good nor evil; it is merely logical. And what will is there, pray, in a Hegelian logic, a good or an evil one? We answer: the will in a Hegelian logic is neither a good nor an evil one, but logical. What will, then, ought there to be here in a Grundideernes Logik, a good or an evil one? We answer: the will in a Grundideernes Logik ought to be neither good nor evil, but logical.

This is the viewpoint from which one has hitherto proceeded in judging the various logics in the various philosophies. But this old, threadbare view will now, with the establishment of the monarchy, have to give way to one ``that sparkles with novelty.'' A beginning has been made with the revision to which Bishop Martensen in his occasional piece has subjected Grundideernes Logik. As a sample we cite a couple of the most forceful passages.

The Bishop says (p. 32): ``Prof. Nielsen simply cannot get rid of theology. He has it not merely outside himself, in that he must constantly combat Christian dogmatics without ever being able to find the end of these attacks, because, as will be shown in what follows, he has not yet been able to find the beginning; but he has it at the same time inside
% --- p. 115 ---
\opage{115}\setcounter{footnote}{0}%
himself, and must now produce a speculative theology, set up a doctrine of God, of the Logos, of the Trinity, of creation and the last things, etc., certainly as opposed to Christian dogmatics as possible, but purely philosophical, purely rational-dogmatic.''

Here is the superordinate! I, who from my subordinate standpoint fancied that I had bent both ends of ``Christian dogmatics'' together and swung the rod of critique, now stand here---with rod raised---and cannot ``find the end.'' How annoying!

But how does it actually come about that theology gets up into me without my knowing it myself? It comes about in a superordinary way; for it comes about in this manner, that the superordinate critique, which, so far as possible, contents itself with criticizing my prefaces and can seldom find time to read a preface all the way through, prefers to seize upon a single piece and throw it in my face.

In the preface to Grundideernes Logik II, Bishop Martensen unfortunately could not bear to read further than p. IX; had he only taken in p. X and p. XI as well, then, before he once again conjured theology into me, he would have had a little problem to solve. In this connection it must suffice to direct the reader's attention to the following statement on my part, which, as it seems to me, is unambiguous (p. XI): ``Instead of engaging with immediate representations of God and bungling with argumentative proofs of God's existence, philosophy begins by presupposing what is not something beyond or absent, but something essentially present for universal knowledge, namely existence itself with its infinitely rich content. This content of knowledge is then determined, according to its content-rich unity with itself, as the Absolute, and the task becomes to develop and clarify the forms of knowledge adequate to the absolute content of being. This task is scien%
% --- p. 116 ---
\opage{116}\setcounter{footnote}{0}%
tific, for such forms can be developed, revised, and corrected with a precision that in dialectics does not fall short of the exactness of mathematics.\footnote{Cf. Mathematik og Dialektik, pp. 1, 2.}
For the development of the forms, the method is so laid out that the content flows in of itself, as by a high pressure of the idea: if only the form is true, then every doubt about the reality of the content is thereby removed. As the perfect form for the absolute content, the concept of God is then developed; in the unity of the developed concept with its content, the idea of God is cognized; in the idea of God, God is cognized; but this cognition is not theological.''

It is, as every thinking reader must tell himself, naturally not the \emph{ethical} but the \emph{logical} concept of God that is to be developed in logic. In order, however, to forestall misunderstanding in this direction, I have in the preface under discussion expressly recalled the relation of logic to the philosophy of religion (p. VIII): ``Although a Grundideernes Logik cannot in any part coincide immediately with the philosophy of religion, it is nevertheless altogether necessary that it, unless it should be able to make do with an atheistic principle, must develop the logical form both of the idea of God and of the idea of the world, must logically determine the relation between these ideas, and thus furnish the foundation for a corresponding philosophy of religion.'' This intimation has in the present case not done me any good.

Assured that nothing here stands in the way of theology, Bishop Martensen has felt prompted to help a little with philosophy's decayed affairs by lodging a protest against my philosophy. I have set up only an ideal of knowledge, but the Bishop demands an ideal of wisdom and appeals to Plato. ``What is highest and deepest in Plato,'' he says, p. 105, ``is most often found in the places \dots\ where he speaks
% --- p. 117 ---
\opage{117}\setcounter{footnote}{0}%
mythically, in that he as it were visionarily touches upon what cannot be reached by the merely dianoetic way, what certainly also cannot be brought into a closed theory, but nonetheless belongs to philosophy, indeed to its highest, since philosophy is essentially striving after wisdom, and therefore not a closed theory. But precisely hereby we are led back to how desirable it would be that Prof. Nielsen had bestowed somewhat greater attention on philosophy, whereby, as said, his attack on theology would have gained in thoroughness and weight, and that he in particular had undertaken to examine the difference between an ideal of knowledge and an ideal of wisdom.''

So it is by an ideal of wisdom, and by drawing attention to the difference between an ideal of knowledge and an ideal of wisdom, that Bishop Martensen now means to take philosophy under his care. On philosophy's behalf we are much obliged to him; only we must remark that he who would take philosophy under his care must first of all see to making himself a little acquainted with its encyclopedia. Without knowledge of the encyclopedia one exposes oneself to getting lost among the disciplines and, for example, to searching for ideals of wisdom where, by the nature of the case, one cannot expect to find anything but ideals of knowledge. In his zeal to hasten to philosophy's aid, the Bishop has taken the encyclopedia in one leap, and, in leaping straight into the ideal of wisdom at such great speed, has come to mix logic, ethics, and philosophy of religion together in a most original way. Through the mixture a complete confusion has entered into the ``total view.''

How thoroughgoing the confusion is can be seen from the most peculiar objections against the theism in Grundideernes Logik with which Bishop Martensen has in part filled his occasional piece. These ethical, thetical, aesthetical, pathetical objections are, to be sure, only variations in which one and the same misunderstanding is smeared out; but they can nonetheless,
% --- p. 118 ---
\opage{118}\setcounter{footnote}{0}%
when one tries to bring some order into them, be brought under the following headings.

1) \emph{The theism in Grundideernes Logik is only a paraphrased pantheism.}

``The author must,'' says Bishop Martensen (Leilighedsskr. p. 33), ``either have taken theism in a wholly different meaning from that in which it is generally taken, or he must here have had a temptation in the direction of mediation. For in reality the knowledge he wishes to introduce is pantheism, or, if one prefers, sheer rationalism and naturalism, which also agrees with his other statements, since he never tires of repeating that there are no other causes than mere causes of reason and natural causes, which must then also hold of the highest cause, and assures us into the bargain that every `supernatural principle' ought to be excluded from science.'' In such a manner Bishop Martensen lodges a protest against Grundideernes Logik.

The main objection is, as one sees, this: that logic's concept of God is only logical, not ethical. By the protest, then, the Bishop evidently aims at introducing not ``a thorough subordination'' but a thorough confusion into the realm of ideas. For there is in this matter something that every reasonable person can understand at once, and it is the following: either the idea of God may not be taken up at all into a logic of ideas, or, if the idea may be taken up into a logic, then the concept of God must be developed logically. A development of the logical concept of God is so far from resting on any denial of the ethical that, on the contrary, it contains all the moments necessary for the ethical concept of God. Logic is, to be sure, not ethics, and the logical concepts are not ethical; but ethics must nonetheless have its logic, and the ethical concepts reflect themselves in the logical idea. The knowledge I wish to introduce is, so the Bishop asserts, ``in reality pantheism''; why? Because logic is supposed to be ethics: a rather interesting confusion.

% --- p. 119 ---
\opage{119}\setcounter{footnote}{0}%
``What'' (Leilighedsskr. p. 33) ``can produce the semblance of theism is this, that the principle which according to this philosophy bears existence, and with it also the philosophical system, is a principle of subjectivity. `If it is true,' says Prof. N., `that subjectivity must be principle not merely for itself but for all objectivity, and can be the former only by virtue of absolute knowledge, the latter by virtue of absolute power, and both as a unity-of-opposition of knowledge and power: then it will already be clear from this that the logical determination of the basic relation between God and the idea of the world must become theistic.' (Grundideernes Logik II. Preface, p. IX). But this inference is far too hasty and can by no means be granted without further ado.''

Why can the inference not be granted? ``One would surely,'' says the Bishop, ``have to be an utter stranger to the history of philosophy not to know that pantheism too has its principle of subjectivity, which has been determined now as a principle of knowledge, now as a principle of power, and now as a unity of both.'' A rather interesting confusion.

The principle of subjectivity that is developed in Grundideernes Logik is, in the course of the development, step by step, indeed point by point, distinguished from and justified by its opposition to the prominent principles of subjectivity in the history of philosophy. Only a critique absolutely heterogeneous to the scientific one can think it accomplishes anything here by an assertion such as this (Leilighedsskr. p. 48), that I ``combine Hegel's logical subjectivity with an eternal principle of power''; only a critique absolutely heterogeneous to the scientific one can suppose that a thinking reader would not discover by himself so thinly veiled a contradiction as this: first it is asserted that the subjectivity in my logic is Hegel's logical subjectivity; then it is conceded in the same breath that I have a quite different concept of logical subjectivity than Hegel. One would scarcely believe it if one could not see it with
% --- p. 120 ---
\opage{120}\setcounter{footnote}{0}%
one's own eyes. But the confusion is a fact, indeed ``a fundamental fact.''

The confused assertion, with the equally confused concession, is to be read in the occasional piece, pp. 48--49, and in its context runs as follows: ``The real thought-content in the Hegelian doctrine of `the concept,' the eternal νοῦς, is the one that Heiberg in his exposition of Hegel's logic sets forth in the following manner: `Subjectivity is expressed by I; but this is not yet consciousness; it is only the subjectivity common to the forms of nature and of spirit (hence logical), and thus \emph{the root from which consciousness develops, though not yet consciousness itself}.' (Pros. Skrifter vol.~1, pp. 276--77). This is clear and intelligible speech. But it is not theistic and does not give itself out to be so. When now Prof. Nielsen combines Hegel's logical subjectivity with an eternal principle of power, one fails to see what has been gained in theistic respect, with respect to the principle of subjectivity itself. Hegel, now, lets the infinite subjectivity, this principle of knowledge that is the root of consciousness but not consciousness itself, come to consciousness in man by the well-known proposition that man's knowledge of God is God's knowledge of himself. This proposition Prof. Nielsen rejects, and with every justification, on account of the incontestable limitation of our psychological consciousness, the dualism between knowledge and power. Absolute knowledge must then be found in ontological subjectivity itself, independently of human consciousness, which only participates in it. But when we then ask how this is to be thought, we certainly get a multiplicity of dialectical determinations of reflection concerning the reciprocal relation between the subjective and the objective, but must especially heed the admonition that we may not measure divine knowledge by our own `psychological yardstick.' But what does this mean other than that we may not personify here,
% --- p. 121 ---
\opage{121}\setcounter{footnote}{0}%
and are dealing only with an idea, a \textit{res cogitans}, which is without actual self-consciousness?''

These are the words. We pray for the assistance of the Spirit: that we may grasp these words! that we may be rightly instructed by ``the royal principle''! that Fichte and Hegel, Heiberg and Kierkegaard may smile upon us from the heaven of logic! that it may be granted us to enjoy the treasure, the rich treasure of logical, psychological, ethical, pathetical confusions with which Bishop Martensen on this occasion has so generously regaled logic!

\emph{First confusion}. Here it is said of Heiberg's speech, which is certainly clear and intelligible, that it neither is theistic nor gives itself out to be so. This peculiar remark has its origin in an equally peculiar confusion of logical and theological theism. In the name of the idea, naturally both Hegel and Heiberg most forcefully defended themselves against having their logical theism confounded with the vulgar representational theism that figures in dogmatic textbooks. In this respect they both had an equally sharp eye for the theologians' well-known confusion of thing and image, concept and representation. But if anyone lets himself be misled by this into the assumption that ``the speculative logic'' was not designed to give us a ``logical theism,'' then he is remarkably mistaken. That Bishop Martensen, with all his dogmatic speculation, was particularly well versed in either Platonic or Hegelian logic, I have to be sure never supposed. As regards Plato, the Bishop has naturally rejected the ``meager fare'' of propædeutics with which a dialogue like the ``Parmenides'' regales one; he has naturally, in order to enlarge his collection of philosophical reading-fruits, picked out for himself ``what is highest and deepest,'' namely the places where Plato ``speaks mystically, in that he as it were visionarily touches upon what cannot be reached by the merely dianoetic way.'' As regards Hegel, it seems to me as if I had once heard that ``the great Logic,'' over which so many in their
% --- p. 122 ---
\opage{122}\setcounter{footnote}{0}%
day got headaches, was declared a kind of barren formalism; that a dogmatic genius could well content himself with the little abstract in the Encyclopedia and trust to his eye for the positive. But that Bishop Martensen should to such a degree have effaced every memory of J. L. Heiberg's logical-dialectical activity, and should to such a degree stand outside the understanding of the logical in ``the speculative logic,'' that he can really think he is saying something to the point when he says that Heiberg's speech about logical subjectivity neither is nor means to be theistic---that has indeed surprised me.

Anyone who has ever read and understood Heiberg's philosophical writings can bear witness with me that Heiberg neither can nor will separate his logical concept of God from the Christian concept of God, although he both can and will acknowledge ``the eternal difference between philosophy and religion.'' The Christian concept of God is theistic; the theistic concept of God is a concretion of the logical concept of subjectivity. ``Since Hegel,'' says Heiberg (Perseus No. 1, pp. 39--40), ``has reduced the philosophical presupposition to so abstract a determination of God that in systematic respect it is no presupposition, and hence also no hindrance to the absolute beginning of the system, it cannot have been taken up contingently from inner experience, as thought or substance or the monad were. Now when one knows that this philosophy ends with the cognition of God, that it presents all development as a cycle in which the end point coincides with the starting point .... then one will see that the system could not begin with God's \emph{most abstract} determination if it did not end with the \emph{most concrete}.'' And yet Heiberg's clear and intelligible speech is supposed to be neither theistic nor to give itself out to be so!

\emph{Second confusion}. When the theism in Grundideernes Logik is designated as \emph{pronounced} theism, this is done in order to distinguish it from a theism that is not pronounced but \emph{obscure}
% --- p. 123 ---
\opage{123}\setcounter{footnote}{0}%
and ambiguous, namely the Hegelian. The difference between obscure and pronounced theism must fundamentally be sought in a different determination of the ontological essence of subjectivity. In this respect Bishop Martensen now mixes Hegelian, non-Hegelian, logical, ethical, selfhood, person, and personification together in a peculiar manner. It is said outright that I combine ``Hegel's logical subjectivity with an eternal principle of power''; the subjectivity in my logic is thus supposed not to differ from the subjectivity in the Hegelian. But then it is said in turn that there is an eternity's difference between my principle of subjectivity and the Hegelian: for in Hegel this principle can first come to consciousness in man, whereas in me absolute knowledge must be found in ontological subjectivity, independently of human consciousness, which only participates in it. And then it is said in turn that this difference is no difference; and why not? Because we may not measure divine knowledge by our own psychological yardstick. When, then, according to the Bishop's instruction, may we measure divine knowledge by our own psychological yardstick? When we mix the ethical into the logical in a suitable way, and introduce ``love'' and ``the good will''! Instead of the knowing self we need only put the willing self, and ontological subjectivity at once comes to consciousness; for in the will, after all, is love. And love? Yes, love! Love in logic and love in ethics and love in dogmatics, love in the three flames, divine love---yes, that is another matter: divine love we may naturally well measure by our own psychological yardstick!

So then: how do we know that the God who \emph{wills} in dogmatics is not a personification, a \textit{res volens}, but a self-conscious personal self? We \emph{must} presumably let Bishop Martensen answer: from this, that the divine will is so perfectly like the human that it can be measured by our own
% --- p. 124 ---
\opage{124}\setcounter{footnote}{0}%
psychological yardstick. And how do we know that a God who \emph{knows} in logic is a personification, a \textit{res cogitans}, not a self-conscious personal self? We \emph{must} presumably let Bishop Martensen answer: from this, that the divine knowledge is so perfectly unlike the human that we may not measure it by our own psychological yardstick. And how, then, do we know that the divine will in dogmatics and the divine knowledge in logic must absolutely be so unlike each other that the will can be measured, but knowledge cannot be measured, by our own psychological yardstick? How do we know that a personification of a \textit{res cogitans} can indeed be thought, but not of a \textit{res volens}? We \emph{must} presumably let Bishop Martensen answer: from this, that here a confusion has been brought about, a thorough confusion in the realm of ideas!

\emph{Third confusion}. There are in our psychological consciousness two heterogeneous elements, a human and a divine, a perishable and an eternal. If now the two elements are mixed together in such a way that one now takes up both where one should take up only the one, now again excludes both where one should exclude only the one: then the confusion is in full swing. It is such a confusion that sets in with what we may here call Bishop Martensen's third confusion. Just as in the first confusion he has mixed up conceptual theism and representational theism, and with the second confusion goes astray among personified and non-personified subjectivity, so in the third confusion he comes to mix the imperishable with the perishable, the divine with the human, in psychological consciousness. The thing is that the author of the occasional piece has been so greedy to find his \emph{ideal of wisdom} in logic that he has not had occasion to assimilate the logical ideal of knowledge.

As for the principle according to which Grundideernes Logik employs psychological consciousness in the development of the ideal of knowledge

% --- p. 125 ---
\opage{125}\setcounter{footnote}{0}%
and the knowing subjectivity, is treated in detail in the first section of the first part. We single out the following (p. 230): ``That in logic, as in every other science, everything essentially depends on trustworthy points of departure and reliable criteria, if investigations are to be undertaken at all, is self-evident. It is the psychological subjectivity, with its presuppositions and conditions, from which we must constantly proceed in Grundideernes Logik, and to which we must return again. In the logical analyses everything depends on how the points of departure given as facts are used: to remain in finitude does not satisfy; to abstract from everything finite and let the thoughts float in the empty infinite does not satisfy either. Here, in the treatment of the relation between perfect knowledge and the psychological factors of knowledge in logic, there is something that recalls the treatment of the so-called `indeterminate forms' in mathematics; none of the finite presuppositions can simply be retained as fact, but every fact can and must be used as a moment: the task is like calculating with infinite magnitudes. Immediately and simply, no one can say how much 0.\,$\infty$ amounts to,\footnote{Cf. Mathematik og Dialektik. Copenhagen 1859. pp. 65--90.} or what value one must assign to $\infty-\infty$; likewise no one can immediately and simply decide how a real concept is absorbed into its existence, and what movements the acts of thought must call forth in the pure idea. Not by a leap, but by a transition corresponding to the peculiar character of the functions, must the comprehending subjectivity in each case move from the finite to the infinite, if it is to succeed in finding the way back from the infinite to the finite.''

Had Bishop Martensen, from his monarchical height, been able to condescend to consider what lies in this and what follows from it, he would surely have spared us the sort of crudities
% --- p. 126 ---
\opage{126}\setcounter{footnote}{0}%
with which, in the occasional piece, he so unctuously holds forth on the difference between personified and personal subjectivity. It is not by saying ``love'' and ``the ethical'' and ``three flames'' that one assures oneself in science of the absolute validity of subjectivity; what is required here is a many-sided, thorough, logical-dialectical analysis of the reciprocal determination between the psychological and the ontological moments of subjectivity. But by such an analysis one will then also be able to assure oneself that an ontological personification can no more replace for us the \emph{knowing} than the \emph{willing} Deity, that we thus --- from the standpoint of knowledge --- have the same, that is, just as great, certainty of a divine self-consciousness in logic as in ethics.

But, the thinking reader will perhaps object, in logic everything is after all so abstract; ``the fullness,'' the true concrete, comes only when we come to dogmatics and to ethics. For the clearing up of misunderstandings belonging here, the author of the occasional piece will perhaps best be able to guide us, since, for the benefit of philosophy, he finds himself in constant progress and is now on the point of entering upon his:

\emph{Fourth Confusion.} In an occasional piece such as the present Martensenian one, everything is of course system --- ``for,'' says Bishop Martensen p. 142, ``an ἀσύστατον, something not cohering in itself, whose multiplicity lacked unity, would be in contradiction with itself'' ---; but when everything in a work is systematic, and the same systematic work is rather rich in confusions, then the confusions must of course also be systematic. Were the systematics to be fully developed here, one would naturally have to take account not only of the mutual connection of the members but also of their number. As regards the number, it is nearly going with us and the confusions in the occasional piece as it probably goes with the astronomers and the stars in the Milky Way: the longer we continue our observations, the more there become. We had, to be sure,
% --- p. 127 ---
\opage{127}\setcounter{footnote}{0}%
thought of reducing the whole number speculatively, by strict economy, to twelve --- the number twelve is much to be recommended in speculative systems, and for various reasons: in the first place it fits a typical tripartition just as well as a predominant quadripartition, and thus offers two great sides for ``a total view''; in the second place the number twelve has a many-sided symbolic significance: twelve signs of the zodiac, twelve tribes in Israel, twelve parish churches in Han-Herred\footnote{If there are more, the superfluous ones must be removed by an abstraction; if there are not so many, the missing ones must be added in a total view; for the total view is the most important thing.} etc. In the third place .... --- But it exceeds our powers; it becomes too lengthy.

So, to the connection! What is fundamental in the confusions rests on an incessant confusion of concept with immediate representation. Uncertainty about what is concept and what is representation must be especially painful for one who wants to help philosophy along. From uncertainty about concept followed the confusion of logical and theological theism: that was the first confusion. From uncertainty about concept followed uncertainty about personal and personified subjectivity, with excessive confidence in dogmatic ``love in three flames'': that was the second confusion. From uncertainty about concept followed uncertainty about the divine and the human in the psychological subjectivity, combined with a characteristic distaste for everything called reflection and determination of reflection, and thus with a characteristic distaste for the logical in logic: that was the third confusion. From these confusions, by no means subordinated but mutually coordinated, one will now easily be able to explain to oneself the fourth and last that we will name here: the confusion with the abstract and the concrete.

Since Grundideernes Logik aims at developing subjectivity as \emph{knowing}, it will naturally not
% --- p. 128 ---
\opage{128}\setcounter{footnote}{0}%
overlook that a knowing subjectivity must essentially be a \emph{willing} one. But, of course, as logic it can bring out only the logical, not the ethical, moments of self-determination, and must above all strive to give us these moments in the form of the concept, not in the form of the image. ``All determination of subjectivity is essentially self-determination, but a self-determination that presupposes something other and is conditioned by something other. That the complete determination of subjectivity assumes a threefold form has its ground precisely in the fact that the determination of the other is to be absorbed into the determination of the self. Now on the level of relativity and finitude this is an impossibility. The determination of the other stands, in the psychological sense, as a barrier, a hindrance to self-determination. Over against the other, the mighty barrier, self-determination is powerless, abstract and formal. The powerlessness is then also seen both in the thought and in the word, for pure self-thinking is an empty reflecting, and the word a name that has no power over the object. ... But when the talk is not of the psychological but of the ontological, not of the relative but of the absolute subjectivity, nominalism must disappear. It is the absolute ideality of the almighty word on which the emphasis falls.'' (Grundid. L. II. pp. 336--37).

The task is a logical development of form; the logical task is: to develop accurately and strictly the form, the only possible form, in which an absolute subjectivity can be thought. Since subjectivity is to correspond in unconditioned unity to the universal, the particular and the individual alike, its logical essentiality must be posited in a syllogism, a middle with two extremes that are opposed to infinity and yet interpenetrate each other, thus in a threefold form. The task is to carry the development of form to the greatest possible completion, not to keep it in the grayish, but to state the moments so definitely that it can be subjected to a critique.

What will happen, now, when a philosophical thinker, a real
% --- p. 129 ---
\opage{129}\setcounter{footnote}{0}%
logician, one day takes it into his head to revise such an exposition? Here there will certainly be no complaint that it is only forms that logic offers, that it cannot also give us ``life'' and ``fullness'' and ``love'' and ``the three flames''; no, here it will be asked whether the form is correct, whether the development of form is many-sided, consistent, stringent. Should it then be found that the form is deficient, correctives will be applied --- logical, be it noted, not theological correctives.

But if it is now not a logician but, on the contrary, a theologian who, on the basis of what he understands by theism and subjectivity, takes it into his head to appraise the theological yield of the development of form: what then? what verdict would we then get? A severe verdict! The logical, the merely logical, has after all no theological fullness; what, then, is this logical? The logical that does not swell with theological fullness is lean as ``the lean fare of propaedeutics''; it is ``the kenological.'' ``The whole,'' says the occasional piece (pp. 75--76), ``is the emptiest formalism, hollow nuts and shadow-play up and down the walls, in the face of which one comes to think, instead of the antilogical, of the \emph{kenological}, because ethical reality lies outside the system.'' The premises of this verdict are by no means built on any logical assessment of the development of form; for in that case something might perhaps still be achieved here by a scientific discussion. No, the critique we have to do with here springs psychologically from ``the good will,'' ontologically from a speculation that puts images in the place of concepts, and is absolutely heterogeneous with scientific critique.

``To seek the absolute,'' it says (p. 74), ``in a sphere that is higher than the logical-physical lies in any case far outside the intentions that appear in the logic of the `fundamental ideas' (be it noted: of the fundamental ideas) treated here. The universe one looks into here is without ethical inhabitants, like one of those uninhabited globes in
% --- p. 130 ---
\opage{130}\setcounter{footnote}{0}%
cosmic space where only natural existences are found, and a corresponding logical and mathematical system.''

For such a misapprehension of the logical-concrete, such a confusion of images and concepts, logic is, however, so well prepared that the author of the occasional piece, merely by turning a leaf, could have found the answer that his ethically pathetic objections so sorely need. ``When the principle of particularity'' --- the talk is precisely of a glance into the absolute --- (Grundid. Lgk. II p. 345) ``sets the universal trembling, when thought shines in the word, and the word, with the light of thought in the color-gleam of relativity, disperses the darkness and opens an entrance for cognition: what visions then present themselves to intuition, what figures will the glance of thought then meet in the realm of ideas? Surely it would never be the `playing possibilities,' `plastic prototypes,' `magical realities' of speculative dogmatics, a `heavenly host' of the `visions' that cannot be content with `a revelation \textit{ad intra},' but are in importunate tension to get out and `be released to revelation \textit{ad extra}'? By no means!... `Magical visions,' speculative hallucinations, are not suited to a divine subjectivity, and poorly enough to a human one. With a host of such visions the individual may perhaps become a fermenting mystic, but not a poet, still less a thinker. The visions that dawn upon the poet, the prototypes that hover before the thinker, always have, as inflows from the wellspring of the ideas, however strong the current may run, the distinguishing mark that they come with the word, arrange themselves in thoughts and are explicated in judgments; even in the mythical consciousness the visions, despite all their magic, must subordinate themselves to the power of the word. But visions that subordinate themselves to the power of the word do not have the `creaturely' forwardness, the apocryphal importunity, that they should immediately `desire to be released to a revelation \textit{ad extra}'; no, on the contrary! they clarify themselves precisely into revelation
% --- p. 131 ---
\opage{131}\setcounter{footnote}{0}%
\textit{ad intra}, for the word of knowledge in the depth of the ideas is revelation \textit{ad intra}.''

2) \emph{The theism in Grundideernes Logik is only a dualism.}

On this objection I can be brief. The question of by what right Grundideernes Logik has put its unity-of-opposition of knowledge and power in place of the Hegelian identity of thinking and being might well deserve a discussion; and it will be a true honor for me to treat this point if anyone with a scientific eye for logic and logical investigations has something to bring forward. But the Martensenian eye is so unfavorable, not merely to my logic but to all logic, that I now almost venture to say in advance: everything that Bishop Martensen --- from the standpoint he has hitherto maintained --- might in future have to communicate to the public about either my logic or that of others will, as far as I am concerned, be laid \textit{ad acta}.

Only as a sample of the peculiarity of the Bishop-Martensenian critique shall we, for the reader's amusement, cite an example. ``We come,'' says Bishop Martensen (Occasional piece p. 56), ``to a new dualism in the system, and get the opposition between the theoretical and the practical as an opposition within theoretical philosophy itself, two divine principles, the knowing one and the mighty one, while we are nevertheless given the explanation that `it is one and the same power which, in mighty self-opposition, is ideal in thinking and judging, real in working and producing.'\,'' What the episcopal critique here dismisses as an explanatory remark is in the logic a fundamental thought that is developed, worked out and shaped in purely logical determinations of concepts. Should there really be something here to correct, a scientific critique would of course necessarily have to enter into logical investigations; but Bishop Martensen's critique and scientific critique are related quite properly as the two absolutely
% --- p. 132 ---
\opage{132}\setcounter{footnote}{0}%
heterogeneous trees, of one of which the fruits were forbidden to man. Instead of paying attention to determinations of thought, reflections, conceptual turns and such formalities, the Bishop, as usual, goes the broad speculative way of the total view, and so entertains his readers with the following explanation. ``We could perhaps,'' he says, ``make the relation clear to ourselves by imagining the practical God or the principle of nature as the \emph{strong blind man} who carries out all the works of thought, and the theoretical, the knowing God as the \emph{seeing lame man} who comprehends, judges and infers, anticipates and subjectively determines everything the blind man punctually carries out, but can himself do nothing whatever, move neither hand nor foot. And the identity, the unity-of-opposition, or `the absolute truth,' we could perhaps imagine thus: that it is the principle of nature or the strong blind man who carries on his back the seeing lame man, who in his `hovering-over reflection' must doubtless hold to the Aristotelian maxim that pure theory is the most blessed. In this way we do after all get a certain identity. For in the lame man we can recognize the mighty knowledge, insofar as the blind man must carry out everything that this paralytic cripple thinks and indicates; and in the blind man we can recognize the knowing power, insofar as the concept not only `hovers over him' but is also instinctively `indwelling' or immanent in him.''

Now this is an image, and what follows is also an image; by ``love'' and ``the three flames,'' does not the Bishop have his whole logic in images! Bishop Martensen is on this occasion even so condescending that he condescends outright to what Grundideernes Logik II (pp. 194--201) has characterized as polemic by opinion.\footnote{That, however, theological polemic by opinion does not go quite so easily at the present moment as it did seventeen years ago, Bishop Martensen will himself be able to see from the fact that N. Lindberg already, in his review of the ``occasional piece'' (Dansk Kirketidende, May 5 and 12), has pointed out with logically categorical certainty several of the confusions on which the distortion relies.}
% --- p. 133 ---
\opage{133}\setcounter{footnote}{0}%
From this one now sees what the position, not merely of my philosophy but of all philosophy, must become when Bishop Martensen one of these days gets his Homeric-Aristotelian-dogmatic universal monarchy established.

3) \emph{The theism in Grundideernes Logik is a naturalism pure and simple.}

It is an objection that recurs again and again in the occasional piece, that the design of Grundideernes Logik cannot possibly be theistic, since the principle of theism is a supernatural principle, and the logic has once and for all excluded the supernatural from science. To this it must be remarked that it goes with the concept of the supernatural as with so many others: it depends on determining its original character and the region to which it belongs.

That the cognition of a higher nature, a higher order of things than this sensuous present world, is not to be excluded from science, anyone can easily convince himself who will merely cast a glance at Grundideernes Logik I (e.g., pp. 452--456) and II (e.g., pp. 450--465). On the other hand, for an understanding of what the ``pronounced theism'' of the logic essentially means, and what it aims at, it will be necessary to determine the relation between logic and philosophy of religion somewhat more closely.

One of the great fundamental questions of modern philosophy is, in my view, the question of the limit of knowledge. That human knowledge has its limits, no one of course denies: it is only the relation between the absolute and the relative side of the limitation that is at issue here. From
% --- p. 134 ---
\opage{134}\setcounter{footnote}{0}%
the standpoint of the critique of reason, Kant set an absolute limit to knowledge, in that he conceived the a priori forms of cognition so exclusively subjectively that they could not be applied at all to the absolutely objective (the thing in itself). The contradiction in which the Kantian philosophy entangled itself through its dualism is illuminated from the logical side in Grundideernes Logik (I. pp. 76--85). When Schelling afterwards went beyond Fichte by making intuitable the objective validity of the subjective categories, and Hegel, in his ``Logic as Science,'' developed the unity of the subjective and the objective into the immanent concept of the idea, philosophical knowledge won back its absolute. But in speculative joy over the absolute and the absolutely absolute, one forgot to make a distinction between the human and the divine knowledge of the absolute: man's knowledge of God became, in an uncritical-mystical way, God's knowledge of himself; the essential limit was obliterated. In Grundideernes Logik the limit is set anew: not in the Kantian sense, such that the objective validity of the categories is denied, but in the ontological sense, such that the difference between human and divine knowledge can be seen in the ideal of our knowledge.

How decisive the logical question of the limit is for the cognition of God, and thereby for the philosophy of religion, is now easy to see. If the absolute limit is taken away, the mystery disappears, and God's essence becomes transparent in the idea. To take one's standpoint in the absolute idea is then to take one's scientific standpoint in the divine center. From a theocentric standpoint the God of faith then quite rightly becomes a God of knowledge; but faith disappears, religion becomes science. If, on the other hand, the limit is posited as absolute, then an absolute mystery is posited at the same time; God's essence is then no longer an essence transparent to man's gaze: a theocentric standpoint is impossible.

As surely as there is an absolute limit of knowledge, an absolute mystery, so surely must there be a principle of faith
% --- p. 135 ---
\opage{135}\setcounter{footnote}{0}%
absolutely heterogeneous with the principle of knowledge; as surely as knowledge and faith constitute two absolutely heterogeneous regions in our higher conscious life, so surely will the one and same God also appear to us in a different form in the light of knowledge than in the light of faith. It is from this point of view that I have drawn attention to a Pauline distinction between the God of knowledge (Rom. 1:19--20) and the God of faith. That such a distinction is valid can also be seen when we consider the corresponding distinction between the philosophical and the evangelical ideal of wisdom. When Paul (1 Cor. 1:22) says of the Greeks that they seek wisdom, he certainly overlooks the relative difference within Greek wisdom itself, the difference, e.g., between an Aristippean and a Platonic, between an Epicurean and a Stoic wisdom. And when he expressly adds that the wisdom he proclaims, in preaching the Crucified, is foolishness to the Greeks, that is, to the Greek lovers of wisdom (οἱ φιλόσοφοι), he does not engage at all in determining the philosophical difference among these lovers of wisdom. The wisdom Paul proclaims is and remains foolishness, just as much to a Socrates and a Plato as to an Aristippus and an Epicurus; in other words: the evangelical ideal of wisdom is absolutely heterogeneous with the philosophical, the ideal of faith absolutely heterogeneous with the ideal of knowledge.

Now if God as God can be cognized both in knowledge and in faith, but in such a way that the cognition of God in knowledge must rest on rational necessity, whereas the cognition of God in faith rests on unconditioned freedom, then the relation between logic and the philosophy of religion appears to us in a new light. If God's essence reveals itself in two forms, that of necessity and that of freedom, that of knowledge and that of faith: then there must surely be able to be two kinds of theism, a scientific and a religious one, a theism of knowledge and a theism of faith. That the two kinds of theism must
% --- p. 136 ---
\opage{136}\setcounter{footnote}{0}%
be related to each other as the spheres of cognition to which they belong, and that the theism of knowledge must in principle be absolutely heterogeneous with the theism of faith, goes without saying. But the absolute heterogeneity by no means excludes the possibility that theism in the one form may have characteristic marks that correspond exactly to theism in the other. On the contrary! it lies in the nature of the case that the characteristic marks in the theism of knowledge must mirror the characteristic marks in the theism of faith as determinations of necessity determinations of freedom, as concepts of knowledge concepts of faith.

On this presupposition one will understand: how a Christian view of the world will be able, not immediately but mediately, to gain influence on philosophy as an autonomous science; how regard to the theism of faith can be a guide in a strictly scientific development of the theism of knowledge; how the theism ``pronounced'' in Grundideernes Logik can --- without at any point overstepping the limit of the logical-physical --- be determined by, and determinative for, a corresponding philosophy of religion.

But what, then, is philosophy of religion? to what sphere must it be assigned? Insofar as the philosophy of religion is a science, it must naturally belong to the domain of knowledge; but insofar as it has religion for its content, it must orient us in the sphere of faith. By this two-sided task its two-sided character as a \emph{border science} is determined. The philosophizing consciousness is, as goes without saying, the subjective medium in which the heterogeneous concepts meet; from this it follows that the unity of consciousness must presuppose a formal unity of concepts corresponding to the determination of the limit, that, in other words, the absolutely heterogeneous concepts must become formally homogeneous in the limit itself. What such a formal homogeneity means can easily be shown. We speak of absolutely heterogeneous principles, but a principle is
% --- p. 137 ---
\opage{137}\setcounter{footnote}{0}%
after all always a principle \textit{in abstracto}. We speak of heterogeneous cognitions, heterogeneous concepts, etc.; but when we abstract from what is peculiar, cognition = cognition, concept = concept.

As a border science, the philosophy of religion gets its concepts of faith from the life of faith and its concepts of knowledge from science. In the latter respect it appears what a logical theism means; in the former respect it depends on religious psychology. It is of no use for the philosopher of religion to want to help science along by assurances about his own faith, his own individual experiences. The question here is not where the philosopher has his experiences from; the question here is about the psychological value of the experiences in their relation to the principle of faith. The apparent contradiction, that the philosophy of religion on the conditions laid down must become a science that treats both of what falls \emph{outside} and of what falls \emph{inside} science, is easily removed. Should anyone assert that there cannot possibly be a knowledge of a cognition of faith when the cognition of faith is by its nature to fall outside knowledge, let him then begin with the beginning and assert the impossibility that there can be a consciousness of the earth when the real earth is by its nature to fall outside consciousness. An objection such as this, that the concepts of faith, by reflecting themselves in knowledge, must absolutely become concepts of knowledge, means nothing; the task of science is precisely to know the heterogeneous as heterogeneous.

As a border science, the philosophy of religion must show us how one and the same concept can become heterogeneous with itself according as it is developed into a concept of knowledge or into a concept of faith. To the concept of faith creation there corresponds, e.g., a concept of knowledge creation, the latter developed from necessity, the former from freedom; to the concept of faith God-man corresponds the concept of knowledge God-man, etc. How different in kind these
% --- p. 138 ---
\opage{138}\setcounter{footnote}{0}%
concepts are can be seen in Strauß, Feuerbach and Kierkegaard. There must now be shown a many-sidedly articulated reciprocal determination between the heterogeneous concepts; this happens in such a way that the concepts of knowledge serve as an intellectual substratum, while the concepts of faith, on the other hand, step into the foreground and are developed systematically. From the reciprocal relation between the concepts of the two spheres there then shines forth the reciprocal relation between the scientific and the religious theism, the theism whose concept of God can be developed with logical necessity in a Grundideernes Logik, and the theism whose revealed God is an object of faith.

But how is it possible that agreement can be brought about between the absolutely heterogeneous concepts? From whichever side we consider them, they seem after all to be in irreconcilable conflict. If we place ourselves with Strauß and Feuerbach on the standpoint of knowledge, the concepts of faith dissolve into sheer contradictions; if we hold fast with Kierkegaard to the standpoint of faith, while at the same time conceding the contradictions shown by science, then we come to assert faith by defying the contradictions and making do with the paradox. It looks rather dark. What if we tried it with an amicable agreement between the contending parties, a sort of concordat? Bishop Martensen, who has so many concordats in his head that he can scarcely hold them, has after all, out of his ``good will,'' offered to serve us with one for temporary use. Yes, if we wanted to be theologians, we would have to give thanks for the offer; but we do not want to be theologians, no, not for a thousand concordats! Theology and concordat are inseparable: Bishop Martensen, as a theologian, cannot possibly be rid of the concordats; he has them not only outside himself, he has them ``also within himself'': he has them in the dogmatics. If it is objected to us that, by our own admission, there is no concordat in Bishop Martensen's ``Christian Dogmatics,'' and that Bishop Martensen's dogmatics must nonetheless surely be theology, since in it we have meant to attack all theology,
% --- p. 139 ---
\opage{139}\setcounter{footnote}{0}%
then we answer: there is not merely one concordat but many kinds of concordats in ``the Christian Dogmatics''; yet not φανερῶς, but κατὰ κρύψιν, for they are all mediated, mediated into gruel and hidden in the gruel.

The absolute mystery makes every concordat between faith and knowledge not merely superfluous but impossible. For the fact that the concepts of faith, in their absolutely decisive opposition to the concepts of knowledge, are burdened with contradictions acquires, in consequence of the absolute mystery, an entirely different significance in the philosophy of religion than in theology. In my little work ``Om Holbergs Kirkehistorie og Theologie'' I have carefully sought to point out the contradictions in which orthodox theology entangled itself when, by its ``humble inconsistencies,'' it took, as the phrase goes, reason captive under the obedience of faith. The contradiction is that theology regards itself not merely as a science but even, as Bishop Martensen expresses it, as ``the central science''; for that the central science, the science with the central, the theocentric-central reason, falls out with reason and wants to take reason captive is undeniably a sign of scientific despair. Orthodox theology has a sharp eye for the characterization of the concepts of faith; but since it is not attentive to the difference between a scientific and a religious motivation, the motivation constantly comes into conflict with the characterization.

As a border science, the philosophy of religion has an entirely different explanation than theology of the contradictions characteristic of the concepts of faith. That the concepts of faith, when they are developed from their own principle, contain contradictions cannot be explained, as from a scholastic-orthodox standpoint, by the claim that the original and necessary categories of sinful reason are in conflict with revealed truth, nor, as from a critical-scientific standpoint,
% --- p. 140 ---
\opage{140}\setcounter{footnote}{0}%
by the claim that the concepts of faith are themselves untrue; no, the contradictions are explained, in consequence of the principle, expressly by the fact that the missing intermediate determinations lie hidden in the absolute mystery. Since, then, there can here be as little question of pressing the concepts of faith under a scientific standard as the concepts of knowledge under a religious one, it will be seen from this that in the philosophy of religion there cannot even be any thought of a concordat between faith and knowledge. Whoever, notwithstanding this, has such clear eyes that he can really see in any of my writings what a Bishop Martensen can see, either ``a concordat'' or something resembling ``a concordat,'' he can see more than I can see; he can see more than Schelling's new clothes: he can see ``the Emperor's new clothes.''

Here, as has been said, is a decisive dilemma: \emph{either theology must go, or it must take its place as the highest science.} Bishop Martensen holds to the latter, I to the former.

Now if Bishop Martensen is right, then he is surely the man who will know how to make his right prevail, to put his Homeric-Aristotelian: Let one be ruler! into force and establish the theological universal monarchy. If that happens, then my objections will quite rightly, as the occasional piece (p. 142) indeed also prophesies, have to be regarded as ``a literary scandal (\textit{parturiunt montes}).'' But then the festivities in the monarchy will also be splendid: on the palm-strewn way, to the ringing of the bells, ``the dogmatics with the image of Christ'' will move in procession through richly adorned streets, solemnly upward toward the objective universal church with the objective universal tower.

But should it turn out that Bishop Martensen is not right; should it turn out that with his concept of the real sciences he is not able to introduce ``a thorough subordination in the realm of ideas,'' and thus --- despite his ``\emph{good will}'' ---
% --- p. 141 ---
\opage{141}\setcounter{footnote}{0}%
is not able to establish the monarchy either: then it would indeed be possible that the Bishop, in view of the many and great obstacles that tower up against his \emph{Let one be ruler!}, might come to think that the turn was now his to change standpoint. Should this be the outcome of the dispute, and should any weight be attached to my endeavors: then I offer with a well-disposed mind my \textit{bona officia}, and would truly regard it as a great honor if I might once more come to work together with so gifted a man as Bishop Martensen.

\begin{center}
\rule[0.5ex]{3em}{0.4pt}\,$\diamond$\,\rule[0.5ex]{3em}{0.4pt}
\end{center}

\section*{Errata.}
\noindent
Page 17, line 4 from top: legitim; read: legitim'';\\
--- 43, line 17 from bottom: saadan read: saadant\\
--- 46, note, line 7 from bottom: \textit{Sanctus} read: \textit{Sancti}\\
--- 61, line 10 from top: hedder det om read: hedder det, som ovenfor anført, om\\
--- 61, note, line 1 from bottom: S. 3 read: S. 7.

\end{document}
